\documentclass[11pt,a4paper]{article}
\usepackage[a4paper,margin=20mm]{geometry}
\usepackage{fontspec}
\setmainfont{Noto Serif}
\setsansfont{Noto Sans}
\usepackage{amsmath,amssymb}
\usepackage{booktabs,tabularx,array}
\usepackage[table]{xcolor}
\usepackage{hyperref}
\definecolor{primary}{RGB}{27,79,114}
\definecolor{accent}{RGB}{40,116,166}
\definecolor{rowtint}{RGB}{235,245,251}
\hypersetup{colorlinks=true,linkcolor=primary,urlcolor=primary}
\setlength{\parskip}{5pt plus 2pt minus 1pt}
\setlength{\parindent}{0pt}
\begin{document}
\begin{center}
  {\LARGE\bfseries\color{primary} PHIẾU HỌC TẬP TRẮC NGHIỆM TOÁN 6}\\[4pt]
  {\normalsize\itshape\color{accent} Bộ sách: Kết nối tri thức với cuộc sống (Tập 1 \& Tập 2 – 9 Chương)}\\[4pt]
  {\small Cấu trúc mỗi chương: 20 câu Dễ (Câu 1--20) $\bullet$ 20 câu Trung bình (Câu 21--40) $\bullet$ 10 câu Khó (Câu 41--50) -- Tổng 50 câu/chương}
\end{center}
\vspace{6pt}\hrule\vspace{10pt}
\section*{\color{primary} CHƯƠNG I: TẬP HỢP CÁC SỐ TỰ NHIÊN}
\textit{Trọng tâm kiến thức SGK: Tập hợp, cách ghi số tự nhiên, thứ tự trong tập hợp số tự nhiên, các phép tính cộng, trừ, nhân, chia và lũy thừa với số mũ tự nhiên.}\par\vspace{4pt}
\subsection*{\color{accent} Phần I. Mức độ Dễ (Nhận biết -- Câu 1 đến Câu 20)}
\textbf{\color{primary} Câu 1.} Cho tập hợp A = \{x $\in$ $\mathbb{N}$ | x < 5\}. Số phần tử của tập hợp A là bao nhiêu?\par
\hspace*{1.5em}\textbf{A.} 4 phần tử\quad\quad \textbf{B.} 3 phần tử\quad\quad \textbf{C.} 6 phần tử\quad\quad \textbf{D.} 5 phần tử\par\vspace{2pt}
\textbf{\color{primary} Câu 2.} Hãy trả lời câu hỏi cơ bản sau: Số liền sau của số tự nhiên 137 là số nào?\par
\hspace*{1.5em}\textbf{A.} 138\quad\quad \textbf{B.} 139\quad\quad \textbf{C.} 137\quad\quad \textbf{D.} 136\par\vspace{2pt}
\textbf{\color{primary} Câu 3.} Trong bài học thuộc Chương I, xét nội dung: Kết quả của phép nhân hai lũy thừa 4$^{4}$ $\cdot$ 4$^{3}$ viết dưới dạng một lũy thừa là. Đáp án đúng là:\par
\hspace*{1.5em}\textbf{A.} 4$^{6}$\quad\quad \textbf{B.} 8$^{7}$\quad\quad \textbf{C.} 4$^{12}$\quad\quad \textbf{D.} 4$^{7}$\par\vspace{2pt}
\textbf{\color{primary} Câu 4.} Dựa vào kiến thức chuẩn SGK của Chương I (Tập hợp các số tự nhiên), trong số tự nhiên 6782, chữ số 7 có giá trị bằng bao nhiêu:\par
\hspace*{1.5em}\textbf{A.} 7\quad\quad \textbf{B.} 7000\quad\quad \textbf{C.} 70\quad\quad \textbf{D.} 700\par\vspace{2pt}
\textbf{\color{primary} Câu 5.} Chọn kết quả đúng nhất khi xét: Giá trị của lũy thừa 3$^{4}$ bằng:\par
\hspace*{1.5em}\textbf{A.} 12\quad\quad \textbf{B.} 27\quad\quad \textbf{C.} 81\quad\quad \textbf{D.} 83\par\vspace{2pt}
\textbf{\color{primary} Câu 6.} Khi làm bài tập cơ bản về tập hợp các số tự nhiên, một học sinh gặp câu hỏi: Cho tập hợp A = \{x $\in$ $\mathbb{N}$ | x < 10\}. Số phần tử của tập hợp A là bao nhiêu. Kết quả nào dưới đây đúng?\par
\hspace*{1.5em}\textbf{A.} 10 phần tử\quad\quad \textbf{B.} 9 phần tử\quad\quad \textbf{C.} 8 phần tử\quad\quad \textbf{D.} 11 phần tử\par\vspace{2pt}
\textbf{\color{primary} Câu 7.} Xác định giá trị hoặc khẳng định đúng cho bài toán sau: Số liền sau của số tự nhiên 222 là số nào:\par
\hspace*{1.5em}\textbf{A.} 223\quad\quad \textbf{B.} 222\quad\quad \textbf{C.} 224\quad\quad \textbf{D.} 221\par\vspace{2pt}
\textbf{\color{primary} Câu 8.} Áp dụng trực tiếp định nghĩa và quy tắc trong Chương I: Kết quả của phép nhân hai lũy thừa 5$^{9}$ $\cdot$ 5$^{3}$ viết dưới dạng một lũy thừa là:\par
\hspace*{1.5em}\textbf{A.} 5$^{11}$\quad\quad \textbf{B.} 5$^{27}$\quad\quad \textbf{C.} 5$^{12}$\quad\quad \textbf{D.} 10$^{12}$\par\vspace{2pt}
\textbf{\color{primary} Câu 9.} Kiểm tra nhanh kiến thức nền tảng: Trong số tự nhiên 5782, chữ số 7 có giá trị bằng bao nhiêu?\par
\hspace*{1.5em}\textbf{A.} 700\quad\quad \textbf{B.} 7000\quad\quad \textbf{C.} 70\quad\quad \textbf{D.} 7\par\vspace{2pt}
\textbf{\color{primary} Câu 10.} Cho dữ kiện nhận biết sau đây: Giá trị của lũy thừa 2$^{3}$ bằng. Phương án chính xác là:\par
\hspace*{1.5em}\textbf{A.} 10\quad\quad \textbf{B.} 4\quad\quad \textbf{C.} 8\quad\quad \textbf{D.} 6\par\vspace{2pt}
\textbf{\color{primary} Câu 11.} Hãy chỉ ra kết quả đúng trong bốn lựa chọn A, B, C, D đối với yêu cầu: Cho tập hợp A = \{x $\in$ $\mathbb{N}$ | x < 15\}. Số phần tử của tập hợp A là bao nhiêu:\par
\hspace*{1.5em}\textbf{A.} 15 phần tử\quad\quad \textbf{B.} 16 phần tử\quad\quad \textbf{C.} 14 phần tử\quad\quad \textbf{D.} 13 phần tử\par\vspace{2pt}
\textbf{\color{primary} Câu 12.} Câu hỏi ôn tập lý thuyết và kỹ năng cơ bản (Chương I): Số liền sau của số tự nhiên 307 là số nào:\par
\hspace*{1.5em}\textbf{A.} 308\quad\quad \textbf{B.} 309\quad\quad \textbf{C.} 306\quad\quad \textbf{D.} 307\par\vspace{2pt}
\textbf{\color{primary} Câu 13.} Theo nội dung sách giáo khoa Kết nối tri thức với cuộc sống (Chương I), kết quả của phép nhân hai lũy thừa 2$^{14}$ $\cdot$ 2$^{3}$ viết dưới dạng một lũy thừa là:\par
\hspace*{1.5em}\textbf{A.} 2$^{42}$\quad\quad \textbf{B.} 2$^{16}$\quad\quad \textbf{C.} 2$^{17}$\quad\quad \textbf{D.} 4$^{17}$\par\vspace{2pt}
\textbf{\color{primary} Câu 14.} Tìm đáp số đúng của bài toán cơ bản sau: Trong số tự nhiên 4782, chữ số 7 có giá trị bằng bao nhiêu:\par
\hspace*{1.5em}\textbf{A.} 70\quad\quad \textbf{B.} 700\quad\quad \textbf{C.} 7\quad\quad \textbf{D.} 7000\par\vspace{2pt}
\textbf{\color{primary} Câu 15.} Xét đối tượng toán học trong Chương I: Giá trị của lũy thừa 4$^{5}$ bằng. Giá trị hoặc tính chất tương ứng là:\par
\hspace*{1.5em}\textbf{A.} 1024\quad\quad \textbf{B.} 256\quad\quad \textbf{C.} 1026\quad\quad \textbf{D.} 20\par\vspace{2pt}
\textbf{\color{primary} Câu 16.} Hoàn thành phép xác định cơ bản sau đây: Cho tập hợp A = \{x $\in$ $\mathbb{N}$ | x < 20\}. Số phần tử của tập hợp A là bao nhiêu:\par
\hspace*{1.5em}\textbf{A.} 21 phần tử\quad\quad \textbf{B.} 19 phần tử\quad\quad \textbf{C.} 20 phần tử\quad\quad \textbf{D.} 18 phần tử\par\vspace{2pt}
\textbf{\color{primary} Câu 17.} Trong các khẳng định hoặc giá trị dưới đây, đâu là đáp án đúng cho câu hỏi: Số liền sau của số tự nhiên 392 là số nào?\par
\hspace*{1.5em}\textbf{A.} 392\quad\quad \textbf{B.} 394\quad\quad \textbf{C.} 393\quad\quad \textbf{D.} 391\par\vspace{2pt}
\textbf{\color{primary} Câu 18.} Bài tập khởi động Chương I: Kết quả của phép nhân hai lũy thừa 3$^{19}$ $\cdot$ 3$^{3}$ viết dưới dạng một lũy thừa là:\par
\hspace*{1.5em}\textbf{A.} 3$^{21}$\quad\quad \textbf{B.} 6$^{22}$\quad\quad \textbf{C.} 3$^{57}$\quad\quad \textbf{D.} 3$^{22}$\par\vspace{2pt}
\textbf{\color{primary} Câu 19.} Quan sát dữ kiện đề bài: Trong số tự nhiên 3782, chữ số 7 có giá trị bằng bao nhiêu. Hãy chọn phương án đúng:\par
\hspace*{1.5em}\textbf{A.} 7000\quad\quad \textbf{B.} 700\quad\quad \textbf{C.} 70\quad\quad \textbf{D.} 7\par\vspace{2pt}
\textbf{\color{primary} Câu 20.} Sử dụng công thức hoặc tính chất đầu bài của Chương I để giải quyết yêu cầu: Giá trị của lũy thừa 3$^{4}$ bằng:\par
\hspace*{1.5em}\textbf{A.} 27\quad\quad \textbf{B.} 81\quad\quad \textbf{C.} 83\quad\quad \textbf{D.} 12\par\vspace{2pt}
\subsection*{\color{accent} Phần II. Mức độ Trung bình (Thông hiểu -- Câu 21 đến Câu 40)}
\textbf{\color{primary} Câu 21.} Tính nhanh giá trị của biểu thức P = 25 $\cdot$ 64 + 25 $\cdot$ 36.:\par
\hspace*{1.5em}\textbf{A.} 2750\quad\quad \textbf{B.} 2250\quad\quad \textbf{C.} 2500\quad\quad \textbf{D.} 2564\par\vspace{2pt}
\textbf{\color{primary} Câu 22.} Vận dụng quy tắc trong Chương I (Tập hợp các số tự nhiên) để giải bài toán: Tìm số tự nhiên x, biết: 2 $\cdot$ x + 12 = 56.:\par
\hspace*{1.5em}\textbf{A.} x = 19\quad\quad \textbf{B.} x = 27\quad\quad \textbf{C.} x = 44\quad\quad \textbf{D.} x = 22\par\vspace{2pt}
\textbf{\color{primary} Câu 23.} Biến đổi và tính toán theo kiến thức Chương I: Tập hợp B = \{x $\in$ $\mathbb{N}^*$ | x $\le$ 24\} có bao nhiêu phần tử. Kết quả thu được là:\par
\hspace*{1.5em}\textbf{A.} 23 phần tử\quad\quad \textbf{B.} 24 phần tử\quad\quad \textbf{C.} 25 phần tử\quad\quad \textbf{D.} 22 phần tử\par\vspace{2pt}
\textbf{\color{primary} Câu 24.} Xét bài toán mức độ thông hiểu sau: Tính nhanh giá trị của biểu thức P = 28 $\cdot$ 64 + 28 $\cdot$ 36.. Đáp án đúng là:\par
\hspace*{1.5em}\textbf{A.} 2864\quad\quad \textbf{B.} 3080\quad\quad \textbf{C.} 2520\quad\quad \textbf{D.} 2800\par\vspace{2pt}
\textbf{\color{primary} Câu 25.} Để hoàn thành phiếu học tập Chương I, hãy xác định: Tìm số tự nhiên x, biết: 2 $\cdot$ x + 18 = 74.:\par
\hspace*{1.5em}\textbf{A.} x = 25\quad\quad \textbf{B.} x = 33\quad\quad \textbf{C.} x = 28\quad\quad \textbf{D.} x = 56\par\vspace{2pt}
\textbf{\color{primary} Câu 26.} Phân tích dữ kiện bài toán: Tập hợp B = \{x $\in$ $\mathbb{N}^*$ | x $\le$ 33\} có bao nhiêu phần tử. Giá trị cần tìm bằng:\par
\hspace*{1.5em}\textbf{A.} 33 phần tử\quad\quad \textbf{B.} 32 phần tử\quad\quad \textbf{C.} 34 phần tử\quad\quad \textbf{D.} 31 phần tử\par\vspace{2pt}
\textbf{\color{primary} Câu 27.} Áp dụng tính chất đã học ở Chương I vào bài tập sau: Tính nhanh giá trị của biểu thức P = 31 $\cdot$ 64 + 31 $\cdot$ 36.:\par
\hspace*{1.5em}\textbf{A.} 3164\quad\quad \textbf{B.} 3100\quad\quad \textbf{C.} 3410\quad\quad \textbf{D.} 2790\par\vspace{2pt}
\textbf{\color{primary} Câu 28.} Thực hiện các bước suy luận hợp lý để giải quyết yêu cầu: Tìm số tự nhiên x, biết: 2 $\cdot$ x + 24 = 92.:\par
\hspace*{1.5em}\textbf{A.} x = 31\quad\quad \textbf{B.} x = 68\quad\quad \textbf{C.} x = 39\quad\quad \textbf{D.} x = 34\par\vspace{2pt}
\textbf{\color{primary} Câu 29.} Cho bài toán tự luyện (Tập hợp các số tự nhiên): Tập hợp B = \{x $\in$ $\mathbb{N}^*$ | x $\le$ 42\} có bao nhiêu phần tử. Phương án nào sau đây chính xác?\par
\hspace*{1.5em}\textbf{A.} 43 phần tử\quad\quad \textbf{B.} 41 phần tử\quad\quad \textbf{C.} 40 phần tử\quad\quad \textbf{D.} 42 phần tử\par\vspace{2pt}
\textbf{\color{primary} Câu 30.} Kết hợp các phép biến đổi trong Chương I, hãy tìm kết quả của bài toán: Tính nhanh giá trị của biểu thức P = 34 $\cdot$ 64 + 34 $\cdot$ 36.:\par
\hspace*{1.5em}\textbf{A.} 3060\quad\quad \textbf{B.} 3400\quad\quad \textbf{C.} 3740\quad\quad \textbf{D.} 3464\par\vspace{2pt}
\textbf{\color{primary} Câu 31.} Một bài kiểm tra định kỳ nội dung Chương I nêu câu hỏi: Tìm số tự nhiên x, biết: 2 $\cdot$ x + 30 = 110.. Hãy chọn đáp số đúng:\par
\hspace*{1.5em}\textbf{A.} x = 40\quad\quad \textbf{B.} x = 37\quad\quad \textbf{C.} x = 80\quad\quad \textbf{D.} x = 45\par\vspace{2pt}
\textbf{\color{primary} Câu 32.} Tìm lời giải đúng cho bài toán thông hiểu: Tập hợp B = \{x $\in$ $\mathbb{N}^*$ | x $\le$ 51\} có bao nhiêu phần tử:\par
\hspace*{1.5em}\textbf{A.} 51 phần tử\quad\quad \textbf{B.} 49 phần tử\quad\quad \textbf{C.} 52 phần tử\quad\quad \textbf{D.} 50 phần tử\par\vspace{2pt}
\textbf{\color{primary} Câu 33.} Sử dụng mối liên hệ giữa các đại lượng trong Chương I: Tính nhanh giá trị của biểu thức P = 37 $\cdot$ 64 + 37 $\cdot$ 36.:\par
\hspace*{1.5em}\textbf{A.} 3700\quad\quad \textbf{B.} 3330\quad\quad \textbf{C.} 3764\quad\quad \textbf{D.} 4070\par\vspace{2pt}
\textbf{\color{primary} Câu 34.} Hãy tính toán và chọn phương án đúng đối với bài tập: Tìm số tự nhiên x, biết: 2 $\cdot$ x + 36 = 128.:\par
\hspace*{1.5em}\textbf{A.} x = 51\quad\quad \textbf{B.} x = 92\quad\quad \textbf{C.} x = 46\quad\quad \textbf{D.} x = 43\par\vspace{2pt}
\textbf{\color{primary} Câu 35.} Dựa trên các tính chất trọng tâm của Chương I, xét bài toán: Tập hợp B = \{x $\in$ $\mathbb{N}^*$ | x $\le$ 60\} có bao nhiêu phần tử:\par
\hspace*{1.5em}\textbf{A.} 60 phần tử\quad\quad \textbf{B.} 61 phần tử\quad\quad \textbf{C.} 59 phần tử\quad\quad \textbf{D.} 58 phần tử\par\vspace{2pt}
\textbf{\color{primary} Câu 36.} Thực hành giải toán (Tập hợp các số tự nhiên): Tính nhanh giá trị của biểu thức P = 40 $\cdot$ 64 + 40 $\cdot$ 36.. Đáp án nào dưới đây thỏa mãn?\par
\hspace*{1.5em}\textbf{A.} 4064\quad\quad \textbf{B.} 4400\quad\quad \textbf{C.} 3600\quad\quad \textbf{D.} 4000\par\vspace{2pt}
\textbf{\color{primary} Câu 37.} Kiểm tra kỹ năng biến đổi và suy luận qua bài toán: Tìm số tự nhiên x, biết: 2 $\cdot$ x + 42 = 146.:\par
\hspace*{1.5em}\textbf{A.} x = 52\quad\quad \textbf{B.} x = 49\quad\quad \textbf{C.} x = 104\quad\quad \textbf{D.} x = 57\par\vspace{2pt}
\textbf{\color{primary} Câu 38.} Cho giả thiết thuộc phạm vi Chương I: Tập hợp B = \{x $\in$ $\mathbb{N}^*$ | x $\le$ 69\} có bao nhiêu phần tử. Kết luận đúng là:\par
\hspace*{1.5em}\textbf{A.} 68 phần tử\quad\quad \textbf{B.} 67 phần tử\quad\quad \textbf{C.} 69 phần tử\quad\quad \textbf{D.} 70 phần tử\par\vspace{2pt}
\textbf{\color{primary} Câu 39.} Rút gọn hoặc tính giá trị theo yêu cầu sau của Chương I: Tính nhanh giá trị của biểu thức P = 43 $\cdot$ 64 + 43 $\cdot$ 36.:\par
\hspace*{1.5em}\textbf{A.} 4300\quad\quad \textbf{B.} 3870\quad\quad \textbf{C.} 4730\quad\quad \textbf{D.} 4364\par\vspace{2pt}
\textbf{\color{primary} Câu 40.} Xác định phương án đúng sau khi thực hiện phép tính / suy luận: Tìm số tự nhiên x, biết: 2 $\cdot$ x + 48 = 164.:\par
\hspace*{1.5em}\textbf{A.} x = 63\quad\quad \textbf{B.} x = 58\quad\quad \textbf{C.} x = 116\quad\quad \textbf{D.} x = 55\par\vspace{2pt}
\subsection*{\color{accent} Phần III. Mức độ Khó (Vận dụng \& Vận dụng cao -- Câu 41 đến Câu 50)}
\textbf{\color{primary} Câu 41.} Tìm số tự nhiên n thỏa mãn: 2$^{n + 1}$ + 2$^{n}$ = 12.:\par
\hspace*{1.5em}\textbf{A.} n = 2\quad\quad \textbf{B.} n = 4\quad\quad \textbf{C.} n = 3\quad\quad \textbf{D.} n = 1\par\vspace{2pt}
\textbf{\color{primary} Câu 42.} Bài toán vận dụng (Chương I – Tập hợp các số tự nhiên): Tính tổng dãy số cách đều S = 2 + 4 + 6 + ... + 44.. Kết quả đúng là:\par
\hspace*{1.5em}\textbf{A.} S = 529\quad\quad \textbf{B.} S = 506\quad\quad \textbf{C.} S = 528\quad\quad \textbf{D.} S = 484\par\vspace{2pt}
\textbf{\color{primary} Câu 43.} Vận dụng kiến thức nâng cao trong phạm vi Chương I để giải bài toán: Tìm số tự nhiên n thỏa mãn: 2$^{n + 1}$ + 2$^{n}$ = 48.:\par
\hspace*{1.5em}\textbf{A.} n = 3\quad\quad \textbf{B.} n = 6\quad\quad \textbf{C.} n = 4\quad\quad \textbf{D.} n = 5\par\vspace{2pt}
\textbf{\color{primary} Câu 44.} Phân tích và giải bài toán tổng hợp sau của Chương I: Tính tổng dãy số cách đều S = 2 + 4 + 6 + ... + 52.:\par
\hspace*{1.5em}\textbf{A.} S = 676\quad\quad \textbf{B.} S = 728\quad\quad \textbf{C.} S = 702\quad\quad \textbf{D.} S = 729\par\vspace{2pt}
\textbf{\color{primary} Câu 45.} Thử thách tư duy trong phạm vi kiến thức Chương I: Tìm số tự nhiên n thỏa mãn: 2$^{n + 1}$ + 2$^{n}$ = 192.. Đáp số chính xác là:\par
\hspace*{1.5em}\textbf{A.} n = 5\quad\quad \textbf{B.} n = 6\quad\quad \textbf{C.} n = 7\quad\quad \textbf{D.} n = 8\par\vspace{2pt}
\textbf{\color{primary} Câu 46.} Sử dụng phương pháp giải toán của Chương I để tìm đáp án cho bài toán: Tính tổng dãy số cách đều S = 2 + 4 + 6 + ... + 60.:\par
\hspace*{1.5em}\textbf{A.} S = 960\quad\quad \textbf{B.} S = 930\quad\quad \textbf{C.} S = 961\quad\quad \textbf{D.} S = 900\par\vspace{2pt}
\textbf{\color{primary} Câu 47.} Xét bài toán vận dụng thực tiễn / nâng cao thuộc Chương I: Tìm số tự nhiên n thỏa mãn: 2$^{n + 1}$ + 2$^{n}$ = 768.:\par
\hspace*{1.5em}\textbf{A.} n = 9\quad\quad \textbf{B.} n = 7\quad\quad \textbf{C.} n = 8\quad\quad \textbf{D.} n = 10\par\vspace{2pt}
\textbf{\color{primary} Câu 48.} Tìm kết quả cuối cùng của bài toán phân loại học sinh khá giỏi (Chương I): Tính tổng dãy số cách đều S = 2 + 4 + 6 + ... + 68.:\par
\hspace*{1.5em}\textbf{A.} S = 1190\quad\quad \textbf{B.} S = 1156\quad\quad \textbf{C.} S = 1225\quad\quad \textbf{D.} S = 1224\par\vspace{2pt}
\textbf{\color{primary} Câu 49.} Kết hợp các tính chất đã học trong Chương I, hãy giải quyết bài toán: Tìm số tự nhiên n thỏa mãn: 2$^{n + 1}$ + 2$^{n}$ = 3072.:\par
\hspace*{1.5em}\textbf{A.} n = 10\quad\quad \textbf{B.} n = 9\quad\quad \textbf{C.} n = 12\quad\quad \textbf{D.} n = 11\par\vspace{2pt}
\textbf{\color{primary} Câu 50.} Câu hỏi vận dụng cao cuối Chương I: Tính tổng dãy số cách đều S = 2 + 4 + 6 + ... + 76.. Phương án đúng là:\par
\hspace*{1.5em}\textbf{A.} S = 1520\quad\quad \textbf{B.} S = 1521\quad\quad \textbf{C.} S = 1444\quad\quad \textbf{D.} S = 1482\par\vspace{2pt}
\subsection*{\color{primary} Bảng đáp án Chương I: Tập hợp các số tự nhiên}
\begin{center}
\renewcommand{\arraystretch}{1.25}
\begin{tabularx}{\textwidth}{*{10}{>{\centering\arraybackslash}X}}
\toprule
\textbf{\color{primary} Cột 1} & \textbf{\color{primary} Cột 2} & \textbf{\color{primary} Cột 3} & \textbf{\color{primary} Cột 4} & \textbf{\color{primary} Cột 5} & \textbf{\color{primary} Cột 6} & \textbf{\color{primary} Cột 7} & \textbf{\color{primary} Cột 8} & \textbf{\color{primary} Cột 9} & \textbf{\color{primary} Cột 10} \\
\midrule
\textbf{1.D} & \textbf{2.A} & \textbf{3.D} & \textbf{4.D} & \textbf{5.C} & \textbf{6.A} & \textbf{7.A} & \textbf{8.C} & \textbf{9.A} & \textbf{10.C} \\
\rowcolor{rowtint} \textbf{11.A} & \textbf{12.A} & \textbf{13.C} & \textbf{14.B} & \textbf{15.A} & \textbf{16.C} & \textbf{17.C} & \textbf{18.D} & \textbf{19.B} & \textbf{20.B} \\
\textbf{21.C} & \textbf{22.D} & \textbf{23.B} & \textbf{24.D} & \textbf{25.C} & \textbf{26.A} & \textbf{27.B} & \textbf{28.D} & \textbf{29.D} & \textbf{30.B} \\
\rowcolor{rowtint} \textbf{31.A} & \textbf{32.A} & \textbf{33.A} & \textbf{34.C} & \textbf{35.A} & \textbf{36.D} & \textbf{37.A} & \textbf{38.C} & \textbf{39.A} & \textbf{40.B} \\
\textbf{41.A} & \textbf{42.B} & \textbf{43.C} & \textbf{44.C} & \textbf{45.B} & \textbf{46.B} & \textbf{47.C} & \textbf{48.A} & \textbf{49.A} & \textbf{50.D} \\
\bottomrule
\end{tabularx}
\end{center}\vspace{10pt}
\section*{\color{primary} CHƯƠNG II: TÍNH CHIA HẾT TRONG TẬP HỢP CÁC SỐ TỰ NHIÊN}
\textit{Trọng tâm kiến thức SGK: Quan hệ chia hết và tính chất, dấu hiệu chia hết cho 2, 3, 5, 9, số nguyên tố, phân tích một số ra thừa số nguyên tố, ƯCLN và BCNN.}\par\vspace{4pt}
\subsection*{\color{accent} Phần I. Mức độ Dễ (Nhận biết -- Câu 1 đến Câu 20)}
\textbf{\color{primary} Câu 1.} Trong các số sau, số nào chia hết cho cả 2 và 5?\par
\hspace*{1.5em}\textbf{A.} 120\quad\quad \textbf{B.} 122\quad\quad \textbf{C.} 129\quad\quad \textbf{D.} 125\par\vspace{2pt}
\textbf{\color{primary} Câu 2.} Hãy trả lời câu hỏi cơ bản sau: Số nào sau đây là ước của 30?\par
\hspace*{1.5em}\textbf{A.} 29\quad\quad \textbf{B.} 60\quad\quad \textbf{C.} 5\quad\quad \textbf{D.} 31\par\vspace{2pt}
\textbf{\color{primary} Câu 3.} Trong bài học thuộc Chương II, xét nội dung: Trong các số tự nhiên từ 12 đến 32, số nào sau đây là số nguyên tố. Đáp án đúng là:\par
\hspace*{1.5em}\textbf{A.} 17\quad\quad \textbf{B.} 15\quad\quad \textbf{C.} 27\quad\quad \textbf{D.} 21\par\vspace{2pt}
\textbf{\color{primary} Câu 4.} Dựa vào kiến thức chuẩn SGK của Chương II (Tính chia hết trong tập hợp các số tự nhiên), trong các số sau, số nào chia hết cho cả 2 và 5:\par
\hspace*{1.5em}\textbf{A.} 159\quad\quad \textbf{B.} 152\quad\quad \textbf{C.} 150\quad\quad \textbf{D.} 155\par\vspace{2pt}
\textbf{\color{primary} Câu 5.} Chọn phương án trả lời đúng nhất: Số nào sau đây là ước của 48?\par
\hspace*{1.5em}\textbf{A.} 8\quad\quad \textbf{B.} 96\quad\quad \textbf{C.} 47\quad\quad \textbf{D.} 49\par\vspace{2pt}
\textbf{\color{primary} Câu 6.} Khi làm bài tập cơ bản về tính chia hết trong tập hợp các số tự nhiên, một học sinh gặp câu hỏi: Trong các số tự nhiên từ 15 đến 35, số nào sau đây là số nguyên tố. Kết quả nào dưới đây đúng?\par
\hspace*{1.5em}\textbf{A.} 15\quad\quad \textbf{B.} 27\quad\quad \textbf{C.} 21\quad\quad \textbf{D.} 29\par\vspace{2pt}
\textbf{\color{primary} Câu 7.} Xác định giá trị hoặc khẳng định đúng cho bài toán sau: Trong các số sau, số nào chia hết cho cả 2 và 5:\par
\hspace*{1.5em}\textbf{A.} 189\quad\quad \textbf{B.} 185\quad\quad \textbf{C.} 182\quad\quad \textbf{D.} 180\par\vspace{2pt}
\textbf{\color{primary} Câu 8.} Áp dụng trực tiếp định nghĩa và quy tắc trong Chương II: Số nào sau đây là ước của 66:\par
\hspace*{1.5em}\textbf{A.} 65\quad\quad \textbf{B.} 11\quad\quad \textbf{C.} 132\quad\quad \textbf{D.} 67\par\vspace{2pt}
\textbf{\color{primary} Câu 9.} Kiểm tra nhanh kiến thức nền tảng: Trong các số tự nhiên từ 18 đến 38, số nào sau đây là số nguyên tố?\par
\hspace*{1.5em}\textbf{A.} 21\quad\quad \textbf{B.} 15\quad\quad \textbf{C.} 27\quad\quad \textbf{D.} 11\par\vspace{2pt}
\textbf{\color{primary} Câu 10.} Cho dữ kiện nhận biết sau đây: Trong các số sau, số nào chia hết cho cả 2 và 5. Phương án chính xác là:\par
\hspace*{1.5em}\textbf{A.} 215\quad\quad \textbf{B.} 210\quad\quad \textbf{C.} 219\quad\quad \textbf{D.} 212\par\vspace{2pt}
\textbf{\color{primary} Câu 11.} Hãy chỉ ra kết quả đúng trong bốn lựa chọn A, B, C, D đối với yêu cầu: Số nào sau đây là ước của 84:\par
\hspace*{1.5em}\textbf{A.} 14\quad\quad \textbf{B.} 168\quad\quad \textbf{C.} 83\quad\quad \textbf{D.} 85\par\vspace{2pt}
\textbf{\color{primary} Câu 12.} Câu hỏi ôn tập lý thuyết và kỹ năng cơ bản (Chương II): Trong các số tự nhiên từ 21 đến 41, số nào sau đây là số nguyên tố:\par
\hspace*{1.5em}\textbf{A.} 19\quad\quad \textbf{B.} 27\quad\quad \textbf{C.} 15\quad\quad \textbf{D.} 21\par\vspace{2pt}
\textbf{\color{primary} Câu 13.} Theo nội dung sách giáo khoa Kết nối tri thức với cuộc sống (Chương II), trong các số sau, số nào chia hết cho cả 2 và 5:\par
\hspace*{1.5em}\textbf{A.} 245\quad\quad \textbf{B.} 240\quad\quad \textbf{C.} 249\quad\quad \textbf{D.} 242\par\vspace{2pt}
\textbf{\color{primary} Câu 14.} Tìm đáp số đúng của bài toán cơ bản sau: Số nào sau đây là ước của 102:\par
\hspace*{1.5em}\textbf{A.} 101\quad\quad \textbf{B.} 204\quad\quad \textbf{C.} 17\quad\quad \textbf{D.} 103\par\vspace{2pt}
\textbf{\color{primary} Câu 15.} Xét đối tượng toán học trong Chương II: Trong các số tự nhiên từ 24 đến 44, số nào sau đây là số nguyên tố. Giá trị hoặc tính chất tương ứng là:\par
\hspace*{1.5em}\textbf{A.} 31\quad\quad \textbf{B.} 21\quad\quad \textbf{C.} 15\quad\quad \textbf{D.} 27\par\vspace{2pt}
\textbf{\color{primary} Câu 16.} Hoàn thành phép xác định cơ bản sau đây: Trong các số sau, số nào chia hết cho cả 2 và 5:\par
\hspace*{1.5em}\textbf{A.} 272\quad\quad \textbf{B.} 275\quad\quad \textbf{C.} 270\quad\quad \textbf{D.} 279\par\vspace{2pt}
\textbf{\color{primary} Câu 17.} Trong các khẳng định hoặc giá trị dưới đây, đâu là đáp án đúng cho câu hỏi: Số nào sau đây là ước của 120?\par
\hspace*{1.5em}\textbf{A.} 119\quad\quad \textbf{B.} 240\quad\quad \textbf{C.} 121\quad\quad \textbf{D.} 20\par\vspace{2pt}
\textbf{\color{primary} Câu 18.} Bài tập khởi động Chương II: Trong các số tự nhiên từ 27 đến 47, số nào sau đây là số nguyên tố:\par
\hspace*{1.5em}\textbf{A.} 21\quad\quad \textbf{B.} 15\quad\quad \textbf{C.} 27\quad\quad \textbf{D.} 13\par\vspace{2pt}
\textbf{\color{primary} Câu 19.} Quan sát dữ kiện đề bài: Trong các số sau, số nào chia hết cho cả 2 và 5. Hãy chọn phương án đúng:\par
\hspace*{1.5em}\textbf{A.} 302\quad\quad \textbf{B.} 309\quad\quad \textbf{C.} 305\quad\quad \textbf{D.} 300\par\vspace{2pt}
\textbf{\color{primary} Câu 20.} Sử dụng công thức hoặc tính chất đầu bài của Chương II để giải quyết yêu cầu: Số nào sau đây là ước của 138:\par
\hspace*{1.5em}\textbf{A.} 276\quad\quad \textbf{B.} 139\quad\quad \textbf{C.} 137\quad\quad \textbf{D.} 23\par\vspace{2pt}
\subsection*{\color{accent} Phần II. Mức độ Trung bình (Thông hiểu -- Câu 21 đến Câu 40)}
\textbf{\color{primary} Câu 21.} Tìm ƯCLN(12, 18).:\par
\hspace*{1.5em}\textbf{A.} 6\quad\quad \textbf{B.} 36\quad\quad \textbf{C.} 12\quad\quad \textbf{D.} 3\par\vspace{2pt}
\textbf{\color{primary} Câu 22.} Vận dụng quy tắc trong Chương II (Tính chia hết trong tập hợp các số tự nhiên) để giải bài toán: Tìm BCNN(8, 12).:\par
\hspace*{1.5em}\textbf{A.} 24\quad\quad \textbf{B.} 12\quad\quad \textbf{C.} 4\quad\quad \textbf{D.} 48\par\vspace{2pt}
\textbf{\color{primary} Câu 23.} Biến đổi và tính toán theo kiến thức Chương II: Phân tích số 72 ra thừa số nguyên tố, khẳng định nào đúng khi xét ước nguyên tố nhỏ nhất của nó. Kết quả thu được là:\par
\begin{itemize}\setlength{\itemsep}{1pt}\setlength{\parskip}{0pt}
  \item[] \textbf{A.} Ước nguyên tố nhỏ nhất là 1
  \item[] \textbf{B.} Ước nguyên tố nhỏ nhất là 2
  \item[] \textbf{C.} Ước nguyên tố nhỏ nhất là 4
  \item[] \textbf{D.} Ước nguyên tố nhỏ nhất là 3
\end{itemize}\vspace{2pt}
\textbf{\color{primary} Câu 24.} Xét bài toán mức độ thông hiểu sau: Tìm ƯCLN(48, 72).. Đáp án đúng là:\par
\hspace*{1.5em}\textbf{A.} 48\quad\quad \textbf{B.} 24\quad\quad \textbf{C.} 144\quad\quad \textbf{D.} 12\par\vspace{2pt}
\textbf{\color{primary} Câu 25.} Để hoàn thành phiếu học tập Chương II, hãy xác định: Tìm BCNN(20, 30).:\par
\hspace*{1.5em}\textbf{A.} 30\quad\quad \textbf{B.} 120\quad\quad \textbf{C.} 60\quad\quad \textbf{D.} 10\par\vspace{2pt}
\textbf{\color{primary} Câu 26.} Phân tích dữ kiện bài toán: Phân tích số 72 ra thừa số nguyên tố, khẳng định nào đúng khi xét ước nguyên tố nhỏ nhất của nó. Giá trị cần tìm bằng:\par
\begin{itemize}\setlength{\itemsep}{1pt}\setlength{\parskip}{0pt}
  \item[] \textbf{A.} Ước nguyên tố nhỏ nhất là 3
  \item[] \textbf{B.} Ước nguyên tố nhỏ nhất là 4
  \item[] \textbf{C.} Ước nguyên tố nhỏ nhất là 2
  \item[] \textbf{D.} Ước nguyên tố nhỏ nhất là 1
\end{itemize}\vspace{2pt}
\textbf{\color{primary} Câu 27.} Áp dụng tính chất đã học ở Chương II vào bài tập sau: Tìm ƯCLN(84, 126).:\par
\hspace*{1.5em}\textbf{A.} 84\quad\quad \textbf{B.} 252\quad\quad \textbf{C.} 42\quad\quad \textbf{D.} 21\par\vspace{2pt}
\textbf{\color{primary} Câu 28.} Thực hiện các bước suy luận hợp lý để giải quyết yêu cầu: Tìm BCNN(32, 48).:\par
\hspace*{1.5em}\textbf{A.} 16\quad\quad \textbf{B.} 192\quad\quad \textbf{C.} 48\quad\quad \textbf{D.} 96\par\vspace{2pt}
\textbf{\color{primary} Câu 29.} Cho bài toán tự luyện (Tính chia hết trong tập hợp các số tự nhiên): Phân tích số 180 ra thừa số nguyên tố, khẳng định nào đúng khi xét ước nguyên tố nhỏ nhất của nó. Phương án nào sau đây chính xác?\par
\begin{itemize}\setlength{\itemsep}{1pt}\setlength{\parskip}{0pt}
  \item[] \textbf{A.} Ước nguyên tố nhỏ nhất là 3
  \item[] \textbf{B.} Ước nguyên tố nhỏ nhất là 2
  \item[] \textbf{C.} Ước nguyên tố nhỏ nhất là 4
  \item[] \textbf{D.} Ước nguyên tố nhỏ nhất là 1
\end{itemize}\vspace{2pt}
\textbf{\color{primary} Câu 30.} Kết hợp các phép biến đổi trong Chương II, hãy tìm kết quả của bài toán: Tìm ƯCLN(120, 180).:\par
\hspace*{1.5em}\textbf{A.} 30\quad\quad \textbf{B.} 360\quad\quad \textbf{C.} 60\quad\quad \textbf{D.} 120\par\vspace{2pt}
\textbf{\color{primary} Câu 31.} Một bài kiểm tra định kỳ nội dung Chương II nêu câu hỏi: Tìm BCNN(44, 66).. Hãy chọn đáp số đúng:\par
\hspace*{1.5em}\textbf{A.} 66\quad\quad \textbf{B.} 264\quad\quad \textbf{C.} 132\quad\quad \textbf{D.} 22\par\vspace{2pt}
\textbf{\color{primary} Câu 32.} Tìm lời giải đúng cho bài toán thông hiểu: Phân tích số 288 ra thừa số nguyên tố, khẳng định nào đúng khi xét ước nguyên tố nhỏ nhất của nó:\par
\begin{itemize}\setlength{\itemsep}{1pt}\setlength{\parskip}{0pt}
  \item[] \textbf{A.} Ước nguyên tố nhỏ nhất là 3
  \item[] \textbf{B.} Ước nguyên tố nhỏ nhất là 4
  \item[] \textbf{C.} Ước nguyên tố nhỏ nhất là 1
  \item[] \textbf{D.} Ước nguyên tố nhỏ nhất là 2
\end{itemize}\vspace{2pt}
\textbf{\color{primary} Câu 33.} Sử dụng mối liên hệ giữa các đại lượng trong Chương II: Tìm ƯCLN(156, 234).:\par
\hspace*{1.5em}\textbf{A.} 156\quad\quad \textbf{B.} 468\quad\quad \textbf{C.} 78\quad\quad \textbf{D.} 39\par\vspace{2pt}
\textbf{\color{primary} Câu 34.} Hãy tính toán và chọn phương án đúng đối với bài tập: Tìm BCNN(56, 84).:\par
\hspace*{1.5em}\textbf{A.} 28\quad\quad \textbf{B.} 336\quad\quad \textbf{C.} 84\quad\quad \textbf{D.} 168\par\vspace{2pt}
\textbf{\color{primary} Câu 35.} Dựa trên các tính chất trọng tâm của Chương II, xét bài toán: Phân tích số 396 ra thừa số nguyên tố, khẳng định nào đúng khi xét ước nguyên tố nhỏ nhất của nó:\par
\begin{itemize}\setlength{\itemsep}{1pt}\setlength{\parskip}{0pt}
  \item[] \textbf{A.} Ước nguyên tố nhỏ nhất là 1
  \item[] \textbf{B.} Ước nguyên tố nhỏ nhất là 2
  \item[] \textbf{C.} Ước nguyên tố nhỏ nhất là 3
  \item[] \textbf{D.} Ước nguyên tố nhỏ nhất là 4
\end{itemize}\vspace{2pt}
\textbf{\color{primary} Câu 36.} Thực hành giải toán (Tính chia hết trong tập hợp các số tự nhiên): Tìm ƯCLN(192, 288).. Đáp án nào dưới đây thỏa mãn?\par
\hspace*{1.5em}\textbf{A.} 48\quad\quad \textbf{B.} 192\quad\quad \textbf{C.} 96\quad\quad \textbf{D.} 576\par\vspace{2pt}
\textbf{\color{primary} Câu 37.} Kiểm tra kỹ năng biến đổi và suy luận qua bài toán: Tìm BCNN(68, 102).:\par
\hspace*{1.5em}\textbf{A.} 408\quad\quad \textbf{B.} 34\quad\quad \textbf{C.} 204\quad\quad \textbf{D.} 102\par\vspace{2pt}
\textbf{\color{primary} Câu 38.} Cho giả thiết thuộc phạm vi Chương II: Phân tích số 504 ra thừa số nguyên tố, khẳng định nào đúng khi xét ước nguyên tố nhỏ nhất của nó. Kết luận đúng là:\par
\begin{itemize}\setlength{\itemsep}{1pt}\setlength{\parskip}{0pt}
  \item[] \textbf{A.} Ước nguyên tố nhỏ nhất là 2
  \item[] \textbf{B.} Ước nguyên tố nhỏ nhất là 1
  \item[] \textbf{C.} Ước nguyên tố nhỏ nhất là 3
  \item[] \textbf{D.} Ước nguyên tố nhỏ nhất là 4
\end{itemize}\vspace{2pt}
\textbf{\color{primary} Câu 39.} Rút gọn hoặc tính giá trị theo yêu cầu sau của Chương II: Tìm ƯCLN(228, 342).:\par
\hspace*{1.5em}\textbf{A.} 684\quad\quad \textbf{B.} 114\quad\quad \textbf{C.} 228\quad\quad \textbf{D.} 57\par\vspace{2pt}
\textbf{\color{primary} Câu 40.} Xác định phương án đúng sau khi thực hiện phép tính / suy luận: Tìm BCNN(80, 120).:\par
\hspace*{1.5em}\textbf{A.} 40\quad\quad \textbf{B.} 120\quad\quad \textbf{C.} 240\quad\quad \textbf{D.} 480\par\vspace{2pt}
\subsection*{\color{accent} Phần III. Mức độ Khó (Vận dụng \& Vận dụng cao -- Câu 41 đến Câu 50)}
\textbf{\color{primary} Câu 41.} Một đội y tế có 24 bác sĩ và 36 điều dưỡng. Có thể chia đội y tế đó nhiều nhất thành mấy tổ sao cho số bác sĩ và điều dưỡng được chia đều vào các tổ?\par
\hspace*{1.5em}\textbf{A.} 6 tổ\quad\quad \textbf{B.} 12 tổ\quad\quad \textbf{C.} 24 tổ\quad\quad \textbf{D.} 4 tổ\par\vspace{2pt}
\textbf{\color{primary} Câu 42.} Bài toán vận dụng (Chương II – Tính chia hết trong tập hợp các số tự nhiên): Tìm số tự nhiên nhỏ nhất khác 0 biết rằng số đó chia hết cho cả 8, 12 và 16.. Kết quả đúng là:\par
\hspace*{1.5em}\textbf{A.} 32\quad\quad \textbf{B.} 48\quad\quad \textbf{C.} 24\quad\quad \textbf{D.} 96\par\vspace{2pt}
\textbf{\color{primary} Câu 43.} Vận dụng kiến thức nâng cao trong phạm vi Chương II để giải bài toán: Một đội y tế có 48 bác sĩ và 72 điều dưỡng. Có thể chia đội y tế đó nhiều nhất thành mấy tổ sao cho số bác sĩ và điều dưỡng được chia đều vào các tổ:\par
\hspace*{1.5em}\textbf{A.} 48 tổ\quad\quad \textbf{B.} 24 tổ\quad\quad \textbf{C.} 8 tổ\quad\quad \textbf{D.} 12 tổ\par\vspace{2pt}
\textbf{\color{primary} Câu 44.} Phân tích và giải bài toán tổng hợp sau của Chương II: Tìm số tự nhiên nhỏ nhất khác 0 biết rằng số đó chia hết cho cả 6, 9 và 12.:\par
\hspace*{1.5em}\textbf{A.} 72\quad\quad \textbf{B.} 36\quad\quad \textbf{C.} 24\quad\quad \textbf{D.} 18\par\vspace{2pt}
\textbf{\color{primary} Câu 45.} Thử thách tư duy trong phạm vi kiến thức Chương II: Một đội y tế có 24 bác sĩ và 36 điều dưỡng. Có thể chia đội y tế đó nhiều nhất thành mấy tổ sao cho số bác sĩ và điều dưỡng được chia đều vào các tổ. Đáp số chính xác là:\par
\hspace*{1.5em}\textbf{A.} 12 tổ\quad\quad \textbf{B.} 6 tổ\quad\quad \textbf{C.} 24 tổ\quad\quad \textbf{D.} 4 tổ\par\vspace{2pt}
\textbf{\color{primary} Câu 46.} Sử dụng phương pháp giải toán của Chương II để tìm đáp án cho bài toán: Tìm số tự nhiên nhỏ nhất khác 0 biết rằng số đó chia hết cho cả 10, 15 và 20.:\par
\hspace*{1.5em}\textbf{A.} 40\quad\quad \textbf{B.} 120\quad\quad \textbf{C.} 60\quad\quad \textbf{D.} 30\par\vspace{2pt}
\textbf{\color{primary} Câu 47.} Xét bài toán vận dụng thực tiễn / nâng cao thuộc Chương II: Một đội y tế có 48 bác sĩ và 72 điều dưỡng. Có thể chia đội y tế đó nhiều nhất thành mấy tổ sao cho số bác sĩ và điều dưỡng được chia đều vào các tổ:\par
\hspace*{1.5em}\textbf{A.} 12 tổ\quad\quad \textbf{B.} 24 tổ\quad\quad \textbf{C.} 8 tổ\quad\quad \textbf{D.} 48 tổ\par\vspace{2pt}
\textbf{\color{primary} Câu 48.} Tìm kết quả cuối cùng của bài toán phân loại học sinh khá giỏi (Chương II): Tìm số tự nhiên nhỏ nhất khác 0 biết rằng số đó chia hết cho cả 8, 12 và 16.:\par
\hspace*{1.5em}\textbf{A.} 48\quad\quad \textbf{B.} 32\quad\quad \textbf{C.} 24\quad\quad \textbf{D.} 96\par\vspace{2pt}
\textbf{\color{primary} Câu 49.} Kết hợp các tính chất đã học trong Chương II, hãy giải quyết bài toán: Một đội y tế có 24 bác sĩ và 36 điều dưỡng. Có thể chia đội y tế đó nhiều nhất thành mấy tổ sao cho số bác sĩ và điều dưỡng được chia đều vào các tổ:\par
\hspace*{1.5em}\textbf{A.} 12 tổ\quad\quad \textbf{B.} 6 tổ\quad\quad \textbf{C.} 24 tổ\quad\quad \textbf{D.} 4 tổ\par\vspace{2pt}
\textbf{\color{primary} Câu 50.} Câu hỏi vận dụng cao cuối Chương II: Tìm số tự nhiên nhỏ nhất khác 0 biết rằng số đó chia hết cho cả 6, 9 và 12.. Phương án đúng là:\par
\hspace*{1.5em}\textbf{A.} 18\quad\quad \textbf{B.} 24\quad\quad \textbf{C.} 36\quad\quad \textbf{D.} 72\par\vspace{2pt}
\subsection*{\color{primary} Bảng đáp án Chương II: Tính chia hết trong tập hợp các số tự nhiên}
\begin{center}
\renewcommand{\arraystretch}{1.25}
\begin{tabularx}{\textwidth}{*{10}{>{\centering\arraybackslash}X}}
\toprule
\textbf{\color{primary} Cột 1} & \textbf{\color{primary} Cột 2} & \textbf{\color{primary} Cột 3} & \textbf{\color{primary} Cột 4} & \textbf{\color{primary} Cột 5} & \textbf{\color{primary} Cột 6} & \textbf{\color{primary} Cột 7} & \textbf{\color{primary} Cột 8} & \textbf{\color{primary} Cột 9} & \textbf{\color{primary} Cột 10} \\
\midrule
\textbf{1.A} & \textbf{2.C} & \textbf{3.A} & \textbf{4.C} & \textbf{5.A} & \textbf{6.D} & \textbf{7.D} & \textbf{8.B} & \textbf{9.D} & \textbf{10.B} \\
\rowcolor{rowtint} \textbf{11.A} & \textbf{12.A} & \textbf{13.B} & \textbf{14.C} & \textbf{15.A} & \textbf{16.C} & \textbf{17.D} & \textbf{18.D} & \textbf{19.D} & \textbf{20.D} \\
\textbf{21.A} & \textbf{22.A} & \textbf{23.B} & \textbf{24.B} & \textbf{25.C} & \textbf{26.C} & \textbf{27.C} & \textbf{28.D} & \textbf{29.B} & \textbf{30.C} \\
\rowcolor{rowtint} \textbf{31.C} & \textbf{32.D} & \textbf{33.C} & \textbf{34.D} & \textbf{35.B} & \textbf{36.C} & \textbf{37.C} & \textbf{38.A} & \textbf{39.B} & \textbf{40.C} \\
\textbf{41.B} & \textbf{42.B} & \textbf{43.B} & \textbf{44.B} & \textbf{45.A} & \textbf{46.C} & \textbf{47.B} & \textbf{48.A} & \textbf{49.A} & \textbf{50.C} \\
\bottomrule
\end{tabularx}
\end{center}\vspace{10pt}
\section*{\color{primary} CHƯƠNG III: SỐ NGUYÊN}
\textit{Trọng tâm kiến thức SGK: Tập hợp các số nguyên, thứ tự trong tập hợp số nguyên, phép cộng, trừ, nhân, chia hết, ước và bội của một số nguyên, quy tắc dấu ngoặc.}\par\vspace{4pt}
\subsection*{\color{accent} Phần I. Mức độ Dễ (Nhận biết -- Câu 1 đến Câu 20)}
\textbf{\color{primary} Câu 1.} Số đối của số nguyên -15 là:\par
\hspace*{1.5em}\textbf{A.} 14\quad\quad \textbf{B.} 0\quad\quad \textbf{C.} 15\quad\quad \textbf{D.} -15\par\vspace{2pt}
\textbf{\color{primary} Câu 2.} Thực hiện yêu cầu nhận biết sau: Kết quả của phép tính (-12) + (-6) là.\par
\hspace*{1.5em}\textbf{A.} 6\quad\quad \textbf{B.} -6\quad\quad \textbf{C.} 18\quad\quad \textbf{D.} -18\par\vspace{2pt}
\textbf{\color{primary} Câu 3.} Trong bài học thuộc Chương III, xét nội dung: Tích của hai số nguyên âm (-5) $\cdot$ (-6) bằng. Đáp án đúng là:\par
\hspace*{1.5em}\textbf{A.} 30\quad\quad \textbf{B.} -11\quad\quad \textbf{C.} 11\quad\quad \textbf{D.} -30\par\vspace{2pt}
\textbf{\color{primary} Câu 4.} Dựa vào kiến thức chuẩn SGK của Chương III (Số nguyên), số đối của số nguyên -24 là:\par
\hspace*{1.5em}\textbf{A.} 0\quad\quad \textbf{B.} 23\quad\quad \textbf{C.} -24\quad\quad \textbf{D.} 24\par\vspace{2pt}
\textbf{\color{primary} Câu 5.} Chọn kết quả đúng nhất khi xét: Kết quả của phép tính (-18) + (-9) là:\par
\hspace*{1.5em}\textbf{A.} 9\quad\quad \textbf{B.} -9\quad\quad \textbf{C.} -27\quad\quad \textbf{D.} 27\par\vspace{2pt}
\textbf{\color{primary} Câu 6.} Khi làm bài tập cơ bản về số nguyên, một học sinh gặp câu hỏi: Tích của hai số nguyên âm (-8) $\cdot$ (-9) bằng. Kết quả nào dưới đây đúng?\par
\hspace*{1.5em}\textbf{A.} -17\quad\quad \textbf{B.} 72\quad\quad \textbf{C.} -72\quad\quad \textbf{D.} 17\par\vspace{2pt}
\textbf{\color{primary} Câu 7.} Xác định giá trị hoặc khẳng định đúng cho bài toán sau: Số đối của số nguyên -33 là:\par
\hspace*{1.5em}\textbf{A.} -33\quad\quad \textbf{B.} 0\quad\quad \textbf{C.} 33\quad\quad \textbf{D.} 32\par\vspace{2pt}
\textbf{\color{primary} Câu 8.} Áp dụng trực tiếp định nghĩa và quy tắc trong Chương III: Kết quả của phép tính (-24) + (-12) là:\par
\hspace*{1.5em}\textbf{A.} -12\quad\quad \textbf{B.} -36\quad\quad \textbf{C.} 36\quad\quad \textbf{D.} 12\par\vspace{2pt}
\textbf{\color{primary} Câu 9.} Kiểm tra nhanh kiến thức nền tảng: Tích của hai số nguyên âm (-11) $\cdot$ (-12) bằng:\par
\hspace*{1.5em}\textbf{A.} -132\quad\quad \textbf{B.} 23\quad\quad \textbf{C.} 132\quad\quad \textbf{D.} -23\par\vspace{2pt}
\textbf{\color{primary} Câu 10.} Cho dữ kiện nhận biết sau đây: Số đối của số nguyên -42 là. Phương án chính xác là:\par
\hspace*{1.5em}\textbf{A.} 42\quad\quad \textbf{B.} 0\quad\quad \textbf{C.} -42\quad\quad \textbf{D.} 41\par\vspace{2pt}
\textbf{\color{primary} Câu 11.} Hãy chỉ ra kết quả đúng trong bốn lựa chọn A, B, C, D đối với yêu cầu: Kết quả của phép tính (-30) + (-15) là:\par
\hspace*{1.5em}\textbf{A.} -15\quad\quad \textbf{B.} -45\quad\quad \textbf{C.} 45\quad\quad \textbf{D.} 15\par\vspace{2pt}
\textbf{\color{primary} Câu 12.} Câu hỏi ôn tập lý thuyết và kỹ năng cơ bản (Chương III): Tích của hai số nguyên âm (-14) $\cdot$ (-15) bằng:\par
\hspace*{1.5em}\textbf{A.} 210\quad\quad \textbf{B.} -210\quad\quad \textbf{C.} -29\quad\quad \textbf{D.} 29\par\vspace{2pt}
\textbf{\color{primary} Câu 13.} Theo nội dung sách giáo khoa Kết nối tri thức với cuộc sống (Chương III), số đối của số nguyên -51 là:\par
\hspace*{1.5em}\textbf{A.} 0\quad\quad \textbf{B.} -51\quad\quad \textbf{C.} 50\quad\quad \textbf{D.} 51\par\vspace{2pt}
\textbf{\color{primary} Câu 14.} Tìm đáp số đúng của bài toán cơ bản sau: Kết quả của phép tính (-36) + (-18) là:\par
\hspace*{1.5em}\textbf{A.} 18\quad\quad \textbf{B.} -18\quad\quad \textbf{C.} -54\quad\quad \textbf{D.} 54\par\vspace{2pt}
\textbf{\color{primary} Câu 15.} Xét đối tượng toán học trong Chương III: Tích của hai số nguyên âm (-17) $\cdot$ (-18) bằng. Giá trị hoặc tính chất tương ứng là:\par
\hspace*{1.5em}\textbf{A.} 306\quad\quad \textbf{B.} -306\quad\quad \textbf{C.} -35\quad\quad \textbf{D.} 35\par\vspace{2pt}
\textbf{\color{primary} Câu 16.} Hoàn thành phép xác định cơ bản sau đây: Số đối của số nguyên -60 là:\par
\hspace*{1.5em}\textbf{A.} 0\quad\quad \textbf{B.} 60\quad\quad \textbf{C.} -60\quad\quad \textbf{D.} 59\par\vspace{2pt}
\textbf{\color{primary} Câu 17.} Trong các khẳng định hoặc giá trị dưới đây, đâu là đáp án đúng cho câu hỏi: Kết quả của phép tính (-42) + (-21) là?\par
\hspace*{1.5em}\textbf{A.} 63\quad\quad \textbf{B.} 21\quad\quad \textbf{C.} -63\quad\quad \textbf{D.} -21\par\vspace{2pt}
\textbf{\color{primary} Câu 18.} Bài tập khởi động Chương III: Tích của hai số nguyên âm (-20) $\cdot$ (-21) bằng:\par
\hspace*{1.5em}\textbf{A.} 41\quad\quad \textbf{B.} -420\quad\quad \textbf{C.} 420\quad\quad \textbf{D.} -41\par\vspace{2pt}
\textbf{\color{primary} Câu 19.} Quan sát dữ kiện đề bài: Số đối của số nguyên -69 là. Hãy chọn phương án đúng:\par
\hspace*{1.5em}\textbf{A.} 0\quad\quad \textbf{B.} 69\quad\quad \textbf{C.} -69\quad\quad \textbf{D.} 68\par\vspace{2pt}
\textbf{\color{primary} Câu 20.} Sử dụng công thức hoặc tính chất đầu bài của Chương III để giải quyết yêu cầu: Kết quả của phép tính (-48) + (-24) là:\par
\hspace*{1.5em}\textbf{A.} 24\quad\quad \textbf{B.} -24\quad\quad \textbf{C.} -72\quad\quad \textbf{D.} 72\par\vspace{2pt}
\subsection*{\color{accent} Phần II. Mức độ Trung bình (Thông hiểu -- Câu 21 đến Câu 40)}
\textbf{\color{primary} Câu 21.} Rút gọn và tính giá trị biểu thức A = (100 - 47) - (100 + 53).:\par
\hspace*{1.5em}\textbf{A.} 200\quad\quad \textbf{B.} 100\quad\quad \textbf{C.} -6\quad\quad \textbf{D.} -100\par\vspace{2pt}
\textbf{\color{primary} Câu 22.} Vận dụng quy tắc trong Chương III (Số nguyên) để giải bài toán: Tìm số nguyên x, biết: 2x + 16 = -8.:\par
\hspace*{1.5em}\textbf{A.} x = -4\quad\quad \textbf{B.} x = -12\quad\quad \textbf{C.} x = 4\quad\quad \textbf{D.} x = 12\par\vspace{2pt}
\textbf{\color{primary} Câu 23.} Biến đổi và tính toán theo kiến thức Chương III: Tính giá trị của biểu thức M = (-4) $\cdot$ 7 + (-4) $\cdot$ (93).. Kết quả thu được là:\par
\hspace*{1.5em}\textbf{A.} -400\quad\quad \textbf{B.} -440\quad\quad \textbf{C.} 400\quad\quad \textbf{D.} -360\par\vspace{2pt}
\textbf{\color{primary} Câu 24.} Xét bài toán mức độ thông hiểu sau: Rút gọn và tính giá trị biểu thức A = (115 - 47) - (115 + 53).. Đáp án đúng là:\par
\hspace*{1.5em}\textbf{A.} -6\quad\quad \textbf{B.} 230\quad\quad \textbf{C.} -100\quad\quad \textbf{D.} 100\par\vspace{2pt}
\textbf{\color{primary} Câu 25.} Để hoàn thành phiếu học tập Chương III, hãy xác định: Tìm số nguyên x, biết: 2x + 22 = -14.:\par
\hspace*{1.5em}\textbf{A.} x = -4\quad\quad \textbf{B.} x = 4\quad\quad \textbf{C.} x = -18\quad\quad \textbf{D.} x = 18\par\vspace{2pt}
\textbf{\color{primary} Câu 26.} Phân tích dữ kiện bài toán: Tính giá trị của biểu thức M = (-7) $\cdot$ 10 + (-7) $\cdot$ (90).. Giá trị cần tìm bằng:\par
\hspace*{1.5em}\textbf{A.} -630\quad\quad \textbf{B.} -770\quad\quad \textbf{C.} 700\quad\quad \textbf{D.} -700\par\vspace{2pt}
\textbf{\color{primary} Câu 27.} Áp dụng tính chất đã học ở Chương III vào bài tập sau: Rút gọn và tính giá trị biểu thức A = (130 - 47) - (130 + 53).:\par
\hspace*{1.5em}\textbf{A.} -6\quad\quad \textbf{B.} 260\quad\quad \textbf{C.} 100\quad\quad \textbf{D.} -100\par\vspace{2pt}
\textbf{\color{primary} Câu 28.} Thực hiện các bước suy luận hợp lý để giải quyết yêu cầu: Tìm số nguyên x, biết: 2x + 28 = -20.:\par
\hspace*{1.5em}\textbf{A.} x = 4\quad\quad \textbf{B.} x = 24\quad\quad \textbf{C.} x = -24\quad\quad \textbf{D.} x = -4\par\vspace{2pt}
\textbf{\color{primary} Câu 29.} Cho bài toán tự luyện (Số nguyên): Tính giá trị của biểu thức M = (-10) $\cdot$ 13 + (-10) $\cdot$ (87).. Phương án nào sau đây chính xác?\par
\hspace*{1.5em}\textbf{A.} -900\quad\quad \textbf{B.} -1000\quad\quad \textbf{C.} 1000\quad\quad \textbf{D.} -1100\par\vspace{2pt}
\textbf{\color{primary} Câu 30.} Kết hợp các phép biến đổi trong Chương III, hãy tìm kết quả của bài toán: Rút gọn và tính giá trị biểu thức A = (145 - 47) - (145 + 53).:\par
\hspace*{1.5em}\textbf{A.} -6\quad\quad \textbf{B.} 290\quad\quad \textbf{C.} -100\quad\quad \textbf{D.} 100\par\vspace{2pt}
\textbf{\color{primary} Câu 31.} Một bài kiểm tra định kỳ nội dung Chương III nêu câu hỏi: Tìm số nguyên x, biết: 2x + 34 = -26.. Hãy chọn đáp số đúng:\par
\hspace*{1.5em}\textbf{A.} x = 4\quad\quad \textbf{B.} x = -30\quad\quad \textbf{C.} x = -4\quad\quad \textbf{D.} x = 30\par\vspace{2pt}
\textbf{\color{primary} Câu 32.} Tìm lời giải đúng cho bài toán thông hiểu: Tính giá trị của biểu thức M = (-13) $\cdot$ 16 + (-13) $\cdot$ (84).:\par
\hspace*{1.5em}\textbf{A.} -1430\quad\quad \textbf{B.} -1170\quad\quad \textbf{C.} -1300\quad\quad \textbf{D.} 1300\par\vspace{2pt}
\textbf{\color{primary} Câu 33.} Sử dụng mối liên hệ giữa các đại lượng trong Chương III: Rút gọn và tính giá trị biểu thức A = (160 - 47) - (160 + 53).:\par
\hspace*{1.5em}\textbf{A.} -6\quad\quad \textbf{B.} 320\quad\quad \textbf{C.} 100\quad\quad \textbf{D.} -100\par\vspace{2pt}
\textbf{\color{primary} Câu 34.} Hãy tính toán và chọn phương án đúng đối với bài tập: Tìm số nguyên x, biết: 2x + 40 = -32.:\par
\hspace*{1.5em}\textbf{A.} x = 4\quad\quad \textbf{B.} x = -4\quad\quad \textbf{C.} x = -36\quad\quad \textbf{D.} x = 36\par\vspace{2pt}
\textbf{\color{primary} Câu 35.} Dựa trên các tính chất trọng tâm của Chương III, xét bài toán: Tính giá trị của biểu thức M = (-16) $\cdot$ 19 + (-16) $\cdot$ (81).:\par
\hspace*{1.5em}\textbf{A.} -1440\quad\quad \textbf{B.} -1600\quad\quad \textbf{C.} -1760\quad\quad \textbf{D.} 1600\par\vspace{2pt}
\textbf{\color{primary} Câu 36.} Thực hành giải toán (Số nguyên): Rút gọn và tính giá trị biểu thức A = (175 - 47) - (175 + 53).. Đáp án nào dưới đây thỏa mãn?\par
\hspace*{1.5em}\textbf{A.} -100\quad\quad \textbf{B.} -6\quad\quad \textbf{C.} 350\quad\quad \textbf{D.} 100\par\vspace{2pt}
\textbf{\color{primary} Câu 37.} Kiểm tra kỹ năng biến đổi và suy luận qua bài toán: Tìm số nguyên x, biết: 2x + 46 = -38.:\par
\hspace*{1.5em}\textbf{A.} x = -42\quad\quad \textbf{B.} x = 4\quad\quad \textbf{C.} x = 42\quad\quad \textbf{D.} x = -4\par\vspace{2pt}
\textbf{\color{primary} Câu 38.} Cho giả thiết thuộc phạm vi Chương III: Tính giá trị của biểu thức M = (-19) $\cdot$ 22 + (-19) $\cdot$ (78).. Kết luận đúng là:\par
\hspace*{1.5em}\textbf{A.} -1710\quad\quad \textbf{B.} -2090\quad\quad \textbf{C.} 1900\quad\quad \textbf{D.} -1900\par\vspace{2pt}
\textbf{\color{primary} Câu 39.} Rút gọn hoặc tính giá trị theo yêu cầu sau của Chương III: Rút gọn và tính giá trị biểu thức A = (190 - 47) - (190 + 53).:\par
\hspace*{1.5em}\textbf{A.} 100\quad\quad \textbf{B.} 380\quad\quad \textbf{C.} -100\quad\quad \textbf{D.} -6\par\vspace{2pt}
\textbf{\color{primary} Câu 40.} Xác định phương án đúng sau khi thực hiện phép tính / suy luận: Tìm số nguyên x, biết: 2x + 52 = -44.:\par
\hspace*{1.5em}\textbf{A.} x = -48\quad\quad \textbf{B.} x = -4\quad\quad \textbf{C.} x = 4\quad\quad \textbf{D.} x = 48\par\vspace{2pt}
\subsection*{\color{accent} Phần III. Mức độ Khó (Vận dụng \& Vận dụng cao -- Câu 41 đến Câu 50)}
\textbf{\color{primary} Câu 41.} Có bao nhiêu số nguyên x thỏa mãn (x - 2) là ước nguyên của 6?\par
\hspace*{1.5em}\textbf{A.} 4 số nguyên\quad\quad \textbf{B.} 2 số nguyên\quad\quad \textbf{C.} 6 số nguyên\quad\quad \textbf{D.} 8 số nguyên\par\vspace{2pt}
\textbf{\color{primary} Câu 42.} Bài toán vận dụng (Chương III – Số nguyên): Tính tổng các số nguyên x thỏa mãn -13 $\le$ x < 13.. Kết quả đúng là:\par
\hspace*{1.5em}\textbf{A.} 13\quad\quad \textbf{B.} -13\quad\quad \textbf{C.} 0\quad\quad \textbf{D.} -26\par\vspace{2pt}
\textbf{\color{primary} Câu 43.} Vận dụng kiến thức nâng cao trong phạm vi Chương III để giải bài toán: Có bao nhiêu số nguyên x thỏa mãn (x - 4) là ước nguyên của 6:\par
\hspace*{1.5em}\textbf{A.} 2 số nguyên\quad\quad \textbf{B.} 6 số nguyên\quad\quad \textbf{C.} 4 số nguyên\quad\quad \textbf{D.} 8 số nguyên\par\vspace{2pt}
\textbf{\color{primary} Câu 44.} Phân tích và giải bài toán tổng hợp sau của Chương III: Tính tổng các số nguyên x thỏa mãn -15 $\le$ x < 15.:\par
\hspace*{1.5em}\textbf{A.} 15\quad\quad \textbf{B.} -30\quad\quad \textbf{C.} -15\quad\quad \textbf{D.} 0\par\vspace{2pt}
\textbf{\color{primary} Câu 45.} Thử thách tư duy trong phạm vi kiến thức Chương III: Có bao nhiêu số nguyên x thỏa mãn (x - 6) là ước nguyên của 6. Đáp số chính xác là:\par
\hspace*{1.5em}\textbf{A.} 8 số nguyên\quad\quad \textbf{B.} 2 số nguyên\quad\quad \textbf{C.} 6 số nguyên\quad\quad \textbf{D.} 4 số nguyên\par\vspace{2pt}
\textbf{\color{primary} Câu 46.} Sử dụng phương pháp giải toán của Chương III để tìm đáp án cho bài toán: Tính tổng các số nguyên x thỏa mãn -17 $\le$ x < 17.:\par
\hspace*{1.5em}\textbf{A.} -17\quad\quad \textbf{B.} 17\quad\quad \textbf{C.} -34\quad\quad \textbf{D.} 0\par\vspace{2pt}
\textbf{\color{primary} Câu 47.} Xét bài toán vận dụng thực tiễn / nâng cao thuộc Chương III: Có bao nhiêu số nguyên x thỏa mãn (x - 8) là ước nguyên của 6:\par
\hspace*{1.5em}\textbf{A.} 2 số nguyên\quad\quad \textbf{B.} 6 số nguyên\quad\quad \textbf{C.} 8 số nguyên\quad\quad \textbf{D.} 4 số nguyên\par\vspace{2pt}
\textbf{\color{primary} Câu 48.} Tìm kết quả cuối cùng của bài toán phân loại học sinh khá giỏi (Chương III): Tính tổng các số nguyên x thỏa mãn -19 $\le$ x < 19.:\par
\hspace*{1.5em}\textbf{A.} -38\quad\quad \textbf{B.} -19\quad\quad \textbf{C.} 19\quad\quad \textbf{D.} 0\par\vspace{2pt}
\textbf{\color{primary} Câu 49.} Kết hợp các tính chất đã học trong Chương III, hãy giải quyết bài toán: Có bao nhiêu số nguyên x thỏa mãn (x - 10) là ước nguyên của 6:\par
\hspace*{1.5em}\textbf{A.} 4 số nguyên\quad\quad \textbf{B.} 2 số nguyên\quad\quad \textbf{C.} 8 số nguyên\quad\quad \textbf{D.} 6 số nguyên\par\vspace{2pt}
\textbf{\color{primary} Câu 50.} Câu hỏi vận dụng cao cuối Chương III: Tính tổng các số nguyên x thỏa mãn -21 $\le$ x < 21.. Phương án đúng là:\par
\hspace*{1.5em}\textbf{A.} 21\quad\quad \textbf{B.} -21\quad\quad \textbf{C.} -42\quad\quad \textbf{D.} 0\par\vspace{2pt}
\subsection*{\color{primary} Bảng đáp án Chương III: Số nguyên}
\begin{center}
\renewcommand{\arraystretch}{1.25}
\begin{tabularx}{\textwidth}{*{10}{>{\centering\arraybackslash}X}}
\toprule
\textbf{\color{primary} Cột 1} & \textbf{\color{primary} Cột 2} & \textbf{\color{primary} Cột 3} & \textbf{\color{primary} Cột 4} & \textbf{\color{primary} Cột 5} & \textbf{\color{primary} Cột 6} & \textbf{\color{primary} Cột 7} & \textbf{\color{primary} Cột 8} & \textbf{\color{primary} Cột 9} & \textbf{\color{primary} Cột 10} \\
\midrule
\textbf{1.C} & \textbf{2.D} & \textbf{3.A} & \textbf{4.D} & \textbf{5.C} & \textbf{6.B} & \textbf{7.C} & \textbf{8.B} & \textbf{9.C} & \textbf{10.A} \\
\rowcolor{rowtint} \textbf{11.B} & \textbf{12.A} & \textbf{13.D} & \textbf{14.C} & \textbf{15.A} & \textbf{16.B} & \textbf{17.C} & \textbf{18.C} & \textbf{19.B} & \textbf{20.C} \\
\textbf{21.D} & \textbf{22.B} & \textbf{23.A} & \textbf{24.C} & \textbf{25.C} & \textbf{26.D} & \textbf{27.D} & \textbf{28.C} & \textbf{29.B} & \textbf{30.C} \\
\rowcolor{rowtint} \textbf{31.B} & \textbf{32.C} & \textbf{33.D} & \textbf{34.C} & \textbf{35.B} & \textbf{36.A} & \textbf{37.A} & \textbf{38.D} & \textbf{39.C} & \textbf{40.A} \\
\textbf{41.D} & \textbf{42.B} & \textbf{43.D} & \textbf{44.C} & \textbf{45.A} & \textbf{46.A} & \textbf{47.C} & \textbf{48.B} & \textbf{49.C} & \textbf{50.B} \\
\bottomrule
\end{tabularx}
\end{center}\vspace{10pt}
\section*{\color{primary} CHƯƠNG IV: MỘT SỐ HÌNH PHẲNG TRONG THỰC TIỄN}
\textit{Trọng tâm kiến thức SGK: Nhận biết tính chất và tính chu vi, diện tích của hình tam giác đều, hình vuông, hình lục giác đều, hình chữ nhật, hình thoi, hình bình hành, hình thang cân (không dùng hình vẽ).}\par\vspace{4pt}
\subsection*{\color{accent} Phần I. Mức độ Dễ (Nhận biết -- Câu 1 đến Câu 20)}
\textbf{\color{primary} Câu 1.} Một hình vuông có độ dài cạnh bằng 4 cm. Chu vi của hình vuông đó bằng:\par
\hspace*{1.5em}\textbf{A.} 8 cm\quad\quad \textbf{B.} 16 cm\quad\quad \textbf{C.} 18 cm\quad\quad \textbf{D.} 16 cm (1)\par\vspace{2pt}
\textbf{\color{primary} Câu 2.} Thực hiện yêu cầu nhận biết sau: Một hình tam giác đều có cạnh bằng 6 cm. Chu vi của tam giác đều đó là.\par
\hspace*{1.5em}\textbf{A.} 36 cm\quad\quad \textbf{B.} 18 cm (1)\quad\quad \textbf{C.} 18 cm\quad\quad \textbf{D.} 24 cm\par\vspace{2pt}
\textbf{\color{primary} Câu 3.} Trong bài học thuộc Chương IV, xét nội dung: Một hình lục giác đều có độ dài một cạnh là 5 dm. Chu vi của hình lục giác đều đó bằng. Đáp án đúng là:\par
\hspace*{1.5em}\textbf{A.} 25 dm\quad\quad \textbf{B.} 40 dm\quad\quad \textbf{C.} 30 dm\quad\quad \textbf{D.} 20 dm\par\vspace{2pt}
\textbf{\color{primary} Câu 4.} Dựa vào kiến thức chuẩn SGK của Chương IV (Một số hình phẳng trong thực tiễn), một hình vuông có độ dài cạnh bằng 7 cm. Chu vi của hình vuông đó bằng:\par
\hspace*{1.5em}\textbf{A.} 49 cm\quad\quad \textbf{B.} 30 cm\quad\quad \textbf{C.} 28 cm\quad\quad \textbf{D.} 14 cm\par\vspace{2pt}
\textbf{\color{primary} Câu 5.} Chọn kết quả đúng nhất khi xét: Một hình tam giác đều có cạnh bằng 9 cm. Chu vi của tam giác đều đó là:\par
\hspace*{1.5em}\textbf{A.} 81 cm\quad\quad \textbf{B.} 54 cm\quad\quad \textbf{C.} 36 cm\quad\quad \textbf{D.} 27 cm\par\vspace{2pt}
\textbf{\color{primary} Câu 6.} Khi làm bài tập cơ bản về một số hình phẳng trong thực tiễn, một học sinh gặp câu hỏi: Một hình lục giác đều có độ dài một cạnh là 8 dm. Chu vi của hình lục giác đều đó bằng. Kết quả nào dưới đây đúng?\par
\hspace*{1.5em}\textbf{A.} 64 dm\quad\quad \textbf{B.} 40 dm\quad\quad \textbf{C.} 48 dm\quad\quad \textbf{D.} 32 dm\par\vspace{2pt}
\textbf{\color{primary} Câu 7.} Xác định giá trị hoặc khẳng định đúng cho bài toán sau: Một hình vuông có độ dài cạnh bằng 10 cm. Chu vi của hình vuông đó bằng:\par
\hspace*{1.5em}\textbf{A.} 40 cm\quad\quad \textbf{B.} 20 cm\quad\quad \textbf{C.} 42 cm\quad\quad \textbf{D.} 100 cm\par\vspace{2pt}
\textbf{\color{primary} Câu 8.} Áp dụng trực tiếp định nghĩa và quy tắc trong Chương IV: Một hình tam giác đều có cạnh bằng 12 cm. Chu vi của tam giác đều đó là:\par
\hspace*{1.5em}\textbf{A.} 36 cm\quad\quad \textbf{B.} 48 cm\quad\quad \textbf{C.} 72 cm\quad\quad \textbf{D.} 144 cm\par\vspace{2pt}
\textbf{\color{primary} Câu 9.} Kiểm tra nhanh kiến thức nền tảng: Một hình lục giác đều có độ dài một cạnh là 11 dm. Chu vi của hình lục giác đều đó bằng:\par
\hspace*{1.5em}\textbf{A.} 55 dm\quad\quad \textbf{B.} 66 dm\quad\quad \textbf{C.} 88 dm\quad\quad \textbf{D.} 44 dm\par\vspace{2pt}
\textbf{\color{primary} Câu 10.} Cho dữ kiện nhận biết sau đây: Một hình vuông có độ dài cạnh bằng 13 cm. Chu vi của hình vuông đó bằng. Phương án chính xác là:\par
\hspace*{1.5em}\textbf{A.} 169 cm\quad\quad \textbf{B.} 26 cm\quad\quad \textbf{C.} 52 cm\quad\quad \textbf{D.} 54 cm\par\vspace{2pt}
\textbf{\color{primary} Câu 11.} Hãy chỉ ra kết quả đúng trong bốn lựa chọn A, B, C, D đối với yêu cầu: Một hình tam giác đều có cạnh bằng 15 cm. Chu vi của tam giác đều đó là:\par
\hspace*{1.5em}\textbf{A.} 45 cm\quad\quad \textbf{B.} 60 cm\quad\quad \textbf{C.} 225 cm\quad\quad \textbf{D.} 90 cm\par\vspace{2pt}
\textbf{\color{primary} Câu 12.} Câu hỏi ôn tập lý thuyết và kỹ năng cơ bản (Chương IV): Một hình lục giác đều có độ dài một cạnh là 14 dm. Chu vi của hình lục giác đều đó bằng:\par
\hspace*{1.5em}\textbf{A.} 70 dm\quad\quad \textbf{B.} 56 dm\quad\quad \textbf{C.} 112 dm\quad\quad \textbf{D.} 84 dm\par\vspace{2pt}
\textbf{\color{primary} Câu 13.} Theo nội dung sách giáo khoa Kết nối tri thức với cuộc sống (Chương IV), một hình vuông có độ dài cạnh bằng 16 cm. Chu vi của hình vuông đó bằng:\par
\hspace*{1.5em}\textbf{A.} 66 cm\quad\quad \textbf{B.} 256 cm\quad\quad \textbf{C.} 64 cm\quad\quad \textbf{D.} 32 cm\par\vspace{2pt}
\textbf{\color{primary} Câu 14.} Tìm đáp số đúng của bài toán cơ bản sau: Một hình tam giác đều có cạnh bằng 18 cm. Chu vi của tam giác đều đó là:\par
\hspace*{1.5em}\textbf{A.} 54 cm\quad\quad \textbf{B.} 72 cm\quad\quad \textbf{C.} 108 cm\quad\quad \textbf{D.} 324 cm\par\vspace{2pt}
\textbf{\color{primary} Câu 15.} Xét đối tượng toán học trong Chương IV: Một hình lục giác đều có độ dài một cạnh là 17 dm. Chu vi của hình lục giác đều đó bằng. Giá trị hoặc tính chất tương ứng là:\par
\hspace*{1.5em}\textbf{A.} 136 dm\quad\quad \textbf{B.} 85 dm\quad\quad \textbf{C.} 102 dm\quad\quad \textbf{D.} 68 dm\par\vspace{2pt}
\textbf{\color{primary} Câu 16.} Hoàn thành phép xác định cơ bản sau đây: Một hình vuông có độ dài cạnh bằng 19 cm. Chu vi của hình vuông đó bằng:\par
\hspace*{1.5em}\textbf{A.} 38 cm\quad\quad \textbf{B.} 78 cm\quad\quad \textbf{C.} 76 cm\quad\quad \textbf{D.} 361 cm\par\vspace{2pt}
\textbf{\color{primary} Câu 17.} Trong các khẳng định hoặc giá trị dưới đây, đâu là đáp án đúng cho câu hỏi: Một hình tam giác đều có cạnh bằng 21 cm. Chu vi của tam giác đều đó là?\par
\hspace*{1.5em}\textbf{A.} 126 cm\quad\quad \textbf{B.} 63 cm\quad\quad \textbf{C.} 84 cm\quad\quad \textbf{D.} 441 cm\par\vspace{2pt}
\textbf{\color{primary} Câu 18.} Bài tập khởi động Chương IV: Một hình lục giác đều có độ dài một cạnh là 20 dm. Chu vi của hình lục giác đều đó bằng:\par
\hspace*{1.5em}\textbf{A.} 80 dm\quad\quad \textbf{B.} 160 dm\quad\quad \textbf{C.} 120 dm\quad\quad \textbf{D.} 100 dm\par\vspace{2pt}
\textbf{\color{primary} Câu 19.} Quan sát dữ kiện đề bài: Một hình vuông có độ dài cạnh bằng 22 cm. Chu vi của hình vuông đó bằng. Hãy chọn phương án đúng:\par
\hspace*{1.5em}\textbf{A.} 88 cm\quad\quad \textbf{B.} 90 cm\quad\quad \textbf{C.} 484 cm\quad\quad \textbf{D.} 44 cm\par\vspace{2pt}
\textbf{\color{primary} Câu 20.} Sử dụng công thức hoặc tính chất đầu bài của Chương IV để giải quyết yêu cầu: Một hình tam giác đều có cạnh bằng 24 cm. Chu vi của tam giác đều đó là:\par
\hspace*{1.5em}\textbf{A.} 144 cm\quad\quad \textbf{B.} 96 cm\quad\quad \textbf{C.} 576 cm\quad\quad \textbf{D.} 72 cm\par\vspace{2pt}
\subsection*{\color{accent} Phần II. Mức độ Trung bình (Thông hiểu -- Câu 21 đến Câu 40)}
\textbf{\color{primary} Câu 21.} Một hình thoi có độ dài hai đường chéo lần lượt là 6 cm và 5 cm. Diện tích của hình thoi đó bằng:\par
\hspace*{1.5em}\textbf{A.} 11 cm$^2$\quad\quad \textbf{B.} 15 cm$^2$\quad\quad \textbf{C.} 30 cm$^2$\quad\quad \textbf{D.} 19 cm$^2$\par\vspace{2pt}
\textbf{\color{primary} Câu 22.} Vận dụng quy tắc trong Chương IV (Một số hình phẳng trong thực tiễn) để giải bài toán: Một hình bình hành có độ dài đáy là 9 cm và chiều cao tương ứng là 5 cm. Diện tích hình bình hành đó là:\par
\hspace*{1.5em}\textbf{A.} 14 cm$^2$\quad\quad \textbf{B.} 28 cm$^2$\quad\quad \textbf{C.} 45 cm$^2$\quad\quad \textbf{D.} 22 cm$^2$\par\vspace{2pt}
\textbf{\color{primary} Câu 23.} Biến đổi và tính toán theo kiến thức Chương IV: Một hình thang cân có hai đáy dài 8 m và 14 m, chiều cao bằng 6 m. Diện tích hình thang cân đó là. Kết quả thu được là:\par
\hspace*{1.5em}\textbf{A.} 66 m$^2$\quad\quad \textbf{B.} 28 m$^2$\quad\quad \textbf{C.} 72 m$^2$\quad\quad \textbf{D.} 132 m$^2$\par\vspace{2pt}
\textbf{\color{primary} Câu 24.} Xét bài toán mức độ thông hiểu sau: Một hình thoi có độ dài hai đường chéo lần lượt là 12 cm và 8 cm. Diện tích của hình thoi đó bằng. Đáp án đúng là:\par
\hspace*{1.5em}\textbf{A.} 20 cm$^2$\quad\quad \textbf{B.} 52 cm$^2$\quad\quad \textbf{C.} 96 cm$^2$\quad\quad \textbf{D.} 48 cm$^2$\par\vspace{2pt}
\textbf{\color{primary} Câu 25.} Để hoàn thành phiếu học tập Chương IV, hãy xác định: Một hình bình hành có độ dài đáy là 12 cm và chiều cao tương ứng là 8 cm. Diện tích hình bình hành đó là:\par
\hspace*{1.5em}\textbf{A.} 20 cm$^2$\quad\quad \textbf{B.} 96 cm$^2$\quad\quad \textbf{C.} 48 cm$^2$\quad\quad \textbf{D.} 40 cm$^2$\par\vspace{2pt}
\textbf{\color{primary} Câu 26.} Phân tích dữ kiện bài toán: Một hình thang cân có hai đáy dài 11 m và 17 m, chiều cao bằng 9 m. Diện tích hình thang cân đó là. Giá trị cần tìm bằng:\par
\hspace*{1.5em}\textbf{A.} 132 m$^2$\quad\quad \textbf{B.} 252 m$^2$\quad\quad \textbf{C.} 37 m$^2$\quad\quad \textbf{D.} 126 m$^2$\par\vspace{2pt}
\textbf{\color{primary} Câu 27.} Áp dụng tính chất đã học ở Chương IV vào bài tập sau: Một hình thoi có độ dài hai đường chéo lần lượt là 18 cm và 11 cm. Diện tích của hình thoi đó bằng:\par
\hspace*{1.5em}\textbf{A.} 198 cm$^2$\quad\quad \textbf{B.} 103 cm$^2$\quad\quad \textbf{C.} 29 cm$^2$\quad\quad \textbf{D.} 99 cm$^2$\par\vspace{2pt}
\textbf{\color{primary} Câu 28.} Thực hiện các bước suy luận hợp lý để giải quyết yêu cầu: Một hình bình hành có độ dài đáy là 15 cm và chiều cao tương ứng là 11 cm. Diện tích hình bình hành đó là:\par
\hspace*{1.5em}\textbf{A.} 165 cm$^2$\quad\quad \textbf{B.} 52 cm$^2$\quad\quad \textbf{C.} 82 cm$^2$\quad\quad \textbf{D.} 26 cm$^2$\par\vspace{2pt}
\textbf{\color{primary} Câu 29.} Cho bài toán tự luyện (Một số hình phẳng trong thực tiễn): Một hình thang cân có hai đáy dài 14 m và 20 m, chiều cao bằng 12 m. Diện tích hình thang cân đó là. Phương án nào sau đây chính xác?\par
\hspace*{1.5em}\textbf{A.} 46 m$^2$\quad\quad \textbf{B.} 204 m$^2$\quad\quad \textbf{C.} 408 m$^2$\quad\quad \textbf{D.} 210 m$^2$\par\vspace{2pt}
\textbf{\color{primary} Câu 30.} Kết hợp các phép biến đổi trong Chương IV, hãy tìm kết quả của bài toán: Một hình thoi có độ dài hai đường chéo lần lượt là 24 cm và 14 cm. Diện tích của hình thoi đó bằng:\par
\hspace*{1.5em}\textbf{A.} 168 cm$^2$\quad\quad \textbf{B.} 336 cm$^2$\quad\quad \textbf{C.} 172 cm$^2$\quad\quad \textbf{D.} 38 cm$^2$\par\vspace{2pt}
\textbf{\color{primary} Câu 31.} Một bài kiểm tra định kỳ nội dung Chương IV nêu câu hỏi: Một hình bình hành có độ dài đáy là 18 cm và chiều cao tương ứng là 14 cm. Diện tích hình bình hành đó là. Hãy chọn đáp số đúng:\par
\hspace*{1.5em}\textbf{A.} 32 cm$^2$\quad\quad \textbf{B.} 252 cm$^2$\quad\quad \textbf{C.} 126 cm$^2$\quad\quad \textbf{D.} 64 cm$^2$\par\vspace{2pt}
\textbf{\color{primary} Câu 32.} Tìm lời giải đúng cho bài toán thông hiểu: Một hình thang cân có hai đáy dài 17 m và 23 m, chiều cao bằng 15 m. Diện tích hình thang cân đó là:\par
\hspace*{1.5em}\textbf{A.} 600 m$^2$\quad\quad \textbf{B.} 306 m$^2$\quad\quad \textbf{C.} 300 m$^2$\quad\quad \textbf{D.} 55 m$^2$\par\vspace{2pt}
\textbf{\color{primary} Câu 33.} Sử dụng mối liên hệ giữa các đại lượng trong Chương IV: Một hình thoi có độ dài hai đường chéo lần lượt là 30 cm và 17 cm. Diện tích của hình thoi đó bằng:\par
\hspace*{1.5em}\textbf{A.} 47 cm$^2$\quad\quad \textbf{B.} 510 cm$^2$\quad\quad \textbf{C.} 259 cm$^2$\quad\quad \textbf{D.} 255 cm$^2$\par\vspace{2pt}
\textbf{\color{primary} Câu 34.} Hãy tính toán và chọn phương án đúng đối với bài tập: Một hình bình hành có độ dài đáy là 21 cm và chiều cao tương ứng là 17 cm. Diện tích hình bình hành đó là:\par
\hspace*{1.5em}\textbf{A.} 178 cm$^2$\quad\quad \textbf{B.} 38 cm$^2$\quad\quad \textbf{C.} 76 cm$^2$\quad\quad \textbf{D.} 357 cm$^2$\par\vspace{2pt}
\textbf{\color{primary} Câu 35.} Dựa trên các tính chất trọng tâm của Chương IV, xét bài toán: Một hình thang cân có hai đáy dài 20 m và 26 m, chiều cao bằng 18 m. Diện tích hình thang cân đó là:\par
\hspace*{1.5em}\textbf{A.} 420 m$^2$\quad\quad \textbf{B.} 414 m$^2$\quad\quad \textbf{C.} 828 m$^2$\quad\quad \textbf{D.} 64 m$^2$\par\vspace{2pt}
\textbf{\color{primary} Câu 36.} Thực hành giải toán (Một số hình phẳng trong thực tiễn): Một hình thoi có độ dài hai đường chéo lần lượt là 36 cm và 20 cm. Diện tích của hình thoi đó bằng. Đáp án nào dưới đây thỏa mãn?\par
\hspace*{1.5em}\textbf{A.} 364 cm$^2$\quad\quad \textbf{B.} 360 cm$^2$\quad\quad \textbf{C.} 720 cm$^2$\quad\quad \textbf{D.} 56 cm$^2$\par\vspace{2pt}
\textbf{\color{primary} Câu 37.} Kiểm tra kỹ năng biến đổi và suy luận qua bài toán: Một hình bình hành có độ dài đáy là 24 cm và chiều cao tương ứng là 20 cm. Diện tích hình bình hành đó là:\par
\hspace*{1.5em}\textbf{A.} 240 cm$^2$\quad\quad \textbf{B.} 480 cm$^2$\quad\quad \textbf{C.} 44 cm$^2$\quad\quad \textbf{D.} 88 cm$^2$\par\vspace{2pt}
\textbf{\color{primary} Câu 38.} Cho giả thiết thuộc phạm vi Chương IV: Một hình thang cân có hai đáy dài 23 m và 29 m, chiều cao bằng 21 m. Diện tích hình thang cân đó là. Kết luận đúng là:\par
\hspace*{1.5em}\textbf{A.} 73 m$^2$\quad\quad \textbf{B.} 546 m$^2$\quad\quad \textbf{C.} 552 m$^2$\quad\quad \textbf{D.} 1092 m$^2$\par\vspace{2pt}
\textbf{\color{primary} Câu 39.} Rút gọn hoặc tính giá trị theo yêu cầu sau của Chương IV: Một hình thoi có độ dài hai đường chéo lần lượt là 42 cm và 23 cm. Diện tích của hình thoi đó bằng:\par
\hspace*{1.5em}\textbf{A.} 483 cm$^2$\quad\quad \textbf{B.} 966 cm$^2$\quad\quad \textbf{C.} 487 cm$^2$\quad\quad \textbf{D.} 65 cm$^2$\par\vspace{2pt}
\textbf{\color{primary} Câu 40.} Xác định phương án đúng sau khi thực hiện phép tính / suy luận: Một hình bình hành có độ dài đáy là 27 cm và chiều cao tương ứng là 23 cm. Diện tích hình bình hành đó là:\par
\hspace*{1.5em}\textbf{A.} 100 cm$^2$\quad\quad \textbf{B.} 50 cm$^2$\quad\quad \textbf{C.} 621 cm$^2$\quad\quad \textbf{D.} 310 cm$^2$\par\vspace{2pt}
\subsection*{\color{accent} Phần III. Mức độ Khó (Vận dụng \& Vận dụng cao -- Câu 41 đến Câu 50)}
\textbf{\color{primary} Câu 41.} Một mảnh sân hình chữ nhật có chiều dài 12 m và chiều rộng 8 m. Người ta lát sân bằng các viên gạch hình vuông cạnh 50 cm (không cắt gạch). Số viên gạch cần dùng là:\par
\hspace*{1.5em}\textbf{A.} 192 viên\quad\quad \textbf{B.} 768 viên\quad\quad \textbf{C.} 384 viên\quad\quad \textbf{D.} 96 viên\par\vspace{2pt}
\textbf{\color{primary} Câu 42.} Bài toán vận dụng (Chương IV – Một số hình phẳng trong thực tiễn): Một khu vườn hình chữ nhật có chu vi bằng 44 m, chiều dài hơn chiều rộng 6 m. Diện tích khu vườn đó bằng. Kết quả đúng là:\par
\hspace*{1.5em}\textbf{A.} 112 m$^2$\quad\quad \textbf{B.} 88 m$^2$\quad\quad \textbf{C.} 104 m$^2$\quad\quad \textbf{D.} 124 m$^2$\par\vspace{2pt}
\textbf{\color{primary} Câu 43.} Vận dụng kiến thức nâng cao trong phạm vi Chương IV để giải bài toán: Một mảnh sân hình chữ nhật có chiều dài 16 m và chiều rộng 12 m. Người ta lát sân bằng các viên gạch hình vuông cạnh 50 cm (không cắt gạch). Số viên gạch cần dùng là:\par
\hspace*{1.5em}\textbf{A.} 384 viên\quad\quad \textbf{B.} 192 viên\quad\quad \textbf{C.} 768 viên\quad\quad \textbf{D.} 1536 viên\par\vspace{2pt}
\textbf{\color{primary} Câu 44.} Phân tích và giải bài toán tổng hợp sau của Chương IV: Một khu vườn hình chữ nhật có chu vi bằng 52 m, chiều dài hơn chiều rộng 4 m. Diện tích khu vườn đó bằng:\par
\hspace*{1.5em}\textbf{A.} 165 m$^2$\quad\quad \textbf{B.} 157 m$^2$\quad\quad \textbf{C.} 177 m$^2$\quad\quad \textbf{D.} 104 m$^2$\par\vspace{2pt}
\textbf{\color{primary} Câu 45.} Thử thách tư duy trong phạm vi kiến thức Chương IV: Một mảnh sân hình chữ nhật có chiều dài 20 m và chiều rộng 16 m. Người ta lát sân bằng các viên gạch hình vuông cạnh 50 cm (không cắt gạch). Số viên gạch cần dùng là. Đáp số chính xác là:\par
\hspace*{1.5em}\textbf{A.} 640 viên\quad\quad \textbf{B.} 2560 viên\quad\quad \textbf{C.} 320 viên\quad\quad \textbf{D.} 1280 viên\par\vspace{2pt}
\textbf{\color{primary} Câu 46.} Sử dụng phương pháp giải toán của Chương IV để tìm đáp án cho bài toán: Một khu vườn hình chữ nhật có chu vi bằng 60 m, chiều dài hơn chiều rộng 8 m. Diện tích khu vườn đó bằng:\par
\hspace*{1.5em}\textbf{A.} 221 m$^2$\quad\quad \textbf{B.} 201 m$^2$\quad\quad \textbf{C.} 120 m$^2$\quad\quad \textbf{D.} 209 m$^2$\par\vspace{2pt}
\textbf{\color{primary} Câu 47.} Xét bài toán vận dụng thực tiễn / nâng cao thuộc Chương IV: Một mảnh sân hình chữ nhật có chiều dài 24 m và chiều rộng 20 m. Người ta lát sân bằng các viên gạch hình vuông cạnh 50 cm (không cắt gạch). Số viên gạch cần dùng là:\par
\hspace*{1.5em}\textbf{A.} 480 viên\quad\quad \textbf{B.} 1920 viên\quad\quad \textbf{C.} 960 viên\quad\quad \textbf{D.} 3840 viên\par\vspace{2pt}
\textbf{\color{primary} Câu 48.} Tìm kết quả cuối cùng của bài toán phân loại học sinh khá giỏi (Chương IV): Một khu vườn hình chữ nhật có chu vi bằng 68 m, chiều dài hơn chiều rộng 6 m. Diện tích khu vườn đó bằng:\par
\hspace*{1.5em}\textbf{A.} 272 m$^2$\quad\quad \textbf{B.} 280 m$^2$\quad\quad \textbf{C.} 292 m$^2$\quad\quad \textbf{D.} 136 m$^2$\par\vspace{2pt}
\textbf{\color{primary} Câu 49.} Kết hợp các tính chất đã học trong Chương IV, hãy giải quyết bài toán: Một mảnh sân hình chữ nhật có chiều dài 28 m và chiều rộng 24 m. Người ta lát sân bằng các viên gạch hình vuông cạnh 50 cm (không cắt gạch). Số viên gạch cần dùng là:\par
\hspace*{1.5em}\textbf{A.} 2688 viên\quad\quad \textbf{B.} 1344 viên\quad\quad \textbf{C.} 672 viên\quad\quad \textbf{D.} 5376 viên\par\vspace{2pt}
\textbf{\color{primary} Câu 50.} Câu hỏi vận dụng cao cuối Chương IV: Một khu vườn hình chữ nhật có chu vi bằng 76 m, chiều dài hơn chiều rộng 4 m. Diện tích khu vườn đó bằng. Phương án đúng là:\par
\hspace*{1.5em}\textbf{A.} 349 m$^2$\quad\quad \textbf{B.} 152 m$^2$\quad\quad \textbf{C.} 357 m$^2$\quad\quad \textbf{D.} 369 m$^2$\par\vspace{2pt}
\subsection*{\color{primary} Bảng đáp án Chương IV: Một số hình phẳng trong thực tiễn}
\begin{center}
\renewcommand{\arraystretch}{1.25}
\begin{tabularx}{\textwidth}{*{10}{>{\centering\arraybackslash}X}}
\toprule
\textbf{\color{primary} Cột 1} & \textbf{\color{primary} Cột 2} & \textbf{\color{primary} Cột 3} & \textbf{\color{primary} Cột 4} & \textbf{\color{primary} Cột 5} & \textbf{\color{primary} Cột 6} & \textbf{\color{primary} Cột 7} & \textbf{\color{primary} Cột 8} & \textbf{\color{primary} Cột 9} & \textbf{\color{primary} Cột 10} \\
\midrule
\textbf{1.B} & \textbf{2.C} & \textbf{3.C} & \textbf{4.C} & \textbf{5.D} & \textbf{6.C} & \textbf{7.A} & \textbf{8.A} & \textbf{9.B} & \textbf{10.C} \\
\rowcolor{rowtint} \textbf{11.A} & \textbf{12.D} & \textbf{13.C} & \textbf{14.A} & \textbf{15.C} & \textbf{16.C} & \textbf{17.B} & \textbf{18.C} & \textbf{19.A} & \textbf{20.D} \\
\textbf{21.B} & \textbf{22.C} & \textbf{23.A} & \textbf{24.D} & \textbf{25.B} & \textbf{26.D} & \textbf{27.D} & \textbf{28.A} & \textbf{29.B} & \textbf{30.A} \\
\rowcolor{rowtint} \textbf{31.B} & \textbf{32.C} & \textbf{33.D} & \textbf{34.D} & \textbf{35.B} & \textbf{36.B} & \textbf{37.B} & \textbf{38.B} & \textbf{39.A} & \textbf{40.C} \\
\textbf{41.C} & \textbf{42.A} & \textbf{43.C} & \textbf{44.A} & \textbf{45.D} & \textbf{46.D} & \textbf{47.B} & \textbf{48.B} & \textbf{49.A} & \textbf{50.C} \\
\bottomrule
\end{tabularx}
\end{center}\vspace{10pt}
\section*{\color{primary} CHƯƠNG V: TÍNH ĐỐI XỨNG CỦA HÌNH PHẲNG TRONG TỰ NHIÊN}
\textit{Trọng tâm kiến thức SGK: Hình có trục đối xứng, hình có tâm đối xứng và vai trò của tính đối xứng trong thế giới tự nhiên (nhận diện qua tính chất hình học và chữ cái in hoa, không dùng hình vẽ).}\par\vspace{4pt}
\subsection*{\color{accent} Phần I. Mức độ Dễ (Nhận biết -- Câu 1 đến Câu 20)}
\textbf{\color{primary} Câu 1.} Hình đa giác đều có 3 cạnh bằng nhau thì có bao nhiêu trục đối xứng?\par
\hspace*{1.5em}\textbf{A.} 6 trục đối xứng\quad\quad \textbf{B.} 1 trục đối xứng\quad\quad \textbf{C.} 3 trục đối xứng\quad\quad \textbf{D.} 2 trục đối xứng\par\vspace{2pt}
\textbf{\color{primary} Câu 2.} Hãy trả lời câu hỏi cơ bản sau: Trong các hình phẳng sau đây, hình nào vừa có trục đối xứng vừa có tâm đối xứng?\par
\begin{itemize}\setlength{\itemsep}{1pt}\setlength{\parskip}{0pt}
  \item[] \textbf{A.} Hình tam giác thường
  \item[] \textbf{B.} Hình tam giác đều
  \item[] \textbf{C.} Hình chữ nhật
  \item[] \textbf{D.} Hình thang cân (hai đáy khác nhau)
\end{itemize}\vspace{2pt}
\textbf{\color{primary} Câu 3.} Trong bài học thuộc Chương V, xét nội dung: Trong các chữ cái in hoa (kiểu chữ không chân chuẩn): A, H, N, P, chữ cái nào có cả trục đối xứng và tâm đối xứng. Đáp án đúng là:\par
\hspace*{1.5em}\textbf{A.} Chữ O\quad\quad \textbf{B.} Chữ A\quad\quad \textbf{C.} Chữ F\quad\quad \textbf{D.} Chữ P\par\vspace{2pt}
\textbf{\color{primary} Câu 4.} Dựa vào kiến thức chuẩn SGK của Chương V (Tính đối xứng của hình phẳng trong tự nhiên), hình đa giác đều có 6 cạnh bằng nhau thì có bao nhiêu trục đối xứng:\par
\hspace*{1.5em}\textbf{A.} 5 trục đối xứng\quad\quad \textbf{B.} 6 trục đối xứng\quad\quad \textbf{C.} 12 trục đối xứng\quad\quad \textbf{D.} 1 trục đối xứng\par\vspace{2pt}
\textbf{\color{primary} Câu 5.} Chọn phương án trả lời đúng nhất: Trong các hình phẳng sau đây, hình nào vừa có trục đối xứng vừa có tâm đối xứng?\par
\begin{itemize}\setlength{\itemsep}{1pt}\setlength{\parskip}{0pt}
  \item[] \textbf{A.} Hình tròn
  \item[] \textbf{B.} Hình tam giác thường
  \item[] \textbf{C.} Hình thang cân (hai đáy khác nhau)
  \item[] \textbf{D.} Hình tam giác đều
\end{itemize}\vspace{2pt}
\textbf{\color{primary} Câu 6.} Khi làm bài tập cơ bản về tính đối xứng của hình phẳng trong tự nhiên, một học sinh gặp câu hỏi: Trong các chữ cái in hoa (kiểu chữ không chân chuẩn): A, H, N, P, chữ cái nào có cả trục đối xứng và tâm đối xứng. Kết quả nào dưới đây đúng?\par
\hspace*{1.5em}\textbf{A.} Chữ P\quad\quad \textbf{B.} Chữ F\quad\quad \textbf{C.} Chữ I\quad\quad \textbf{D.} Chữ A\par\vspace{2pt}
\textbf{\color{primary} Câu 7.} Xác định giá trị hoặc khẳng định đúng cho bài toán sau: Hình đa giác đều có 9 cạnh bằng nhau thì có bao nhiêu trục đối xứng:\par
\hspace*{1.5em}\textbf{A.} 1 trục đối xứng\quad\quad \textbf{B.} 18 trục đối xứng\quad\quad \textbf{C.} 9 trục đối xứng\quad\quad \textbf{D.} 8 trục đối xứng\par\vspace{2pt}
\textbf{\color{primary} Câu 8.} Áp dụng trực tiếp định nghĩa và quy tắc trong Chương V: Trong các hình phẳng sau đây, hình nào vừa có trục đối xứng vừa có tâm đối xứng:\par
\begin{itemize}\setlength{\itemsep}{1pt}\setlength{\parskip}{0pt}
  \item[] \textbf{A.} Hình thoi
  \item[] \textbf{B.} Hình thang cân (hai đáy khác nhau)
  \item[] \textbf{C.} Hình tam giác đều
  \item[] \textbf{D.} Hình tam giác thường
\end{itemize}\vspace{2pt}
\textbf{\color{primary} Câu 9.} Kiểm tra nhanh kiến thức nền tảng: Trong các chữ cái in hoa (kiểu chữ không chân chuẩn): A, H, N, P, chữ cái nào có cả trục đối xứng và tâm đối xứng?\par
\hspace*{1.5em}\textbf{A.} Chữ F\quad\quad \textbf{B.} Chữ A\quad\quad \textbf{C.} Chữ P\quad\quad \textbf{D.} Chữ H\par\vspace{2pt}
\textbf{\color{primary} Câu 10.} Cho dữ kiện nhận biết sau đây: Hình đa giác đều có 12 cạnh bằng nhau thì có bao nhiêu trục đối xứng. Phương án chính xác là:\par
\hspace*{1.5em}\textbf{A.} 11 trục đối xứng\quad\quad \textbf{B.} 24 trục đối xứng\quad\quad \textbf{C.} 12 trục đối xứng\quad\quad \textbf{D.} 1 trục đối xứng\par\vspace{2pt}
\textbf{\color{primary} Câu 11.} Hãy chỉ ra kết quả đúng trong bốn lựa chọn A, B, C, D đối với yêu cầu: Trong các hình phẳng sau đây, hình nào vừa có trục đối xứng vừa có tâm đối xứng:\par
\begin{itemize}\setlength{\itemsep}{1pt}\setlength{\parskip}{0pt}
  \item[] \textbf{A.} Hình tam giác đều
  \item[] \textbf{B.} Hình thang cân (hai đáy khác nhau)
  \item[] \textbf{C.} Hình vuông
  \item[] \textbf{D.} Hình tam giác thường
\end{itemize}\vspace{2pt}
\textbf{\color{primary} Câu 12.} Câu hỏi ôn tập lý thuyết và kỹ năng cơ bản (Chương V): Trong các chữ cái in hoa (kiểu chữ không chân chuẩn): A, H, N, P, chữ cái nào có cả trục đối xứng và tâm đối xứng:\par
\hspace*{1.5em}\textbf{A.} Chữ A\quad\quad \textbf{B.} Chữ F\quad\quad \textbf{C.} Chữ X\quad\quad \textbf{D.} Chữ P\par\vspace{2pt}
\textbf{\color{primary} Câu 13.} Theo nội dung sách giáo khoa Kết nối tri thức với cuộc sống (Chương V), hình đa giác đều có 15 cạnh bằng nhau thì có bao nhiêu trục đối xứng:\par
\hspace*{1.5em}\textbf{A.} 15 trục đối xứng\quad\quad \textbf{B.} 30 trục đối xứng\quad\quad \textbf{C.} 14 trục đối xứng\quad\quad \textbf{D.} 1 trục đối xứng\par\vspace{2pt}
\textbf{\color{primary} Câu 14.} Tìm đáp số đúng của bài toán cơ bản sau: Trong các hình phẳng sau đây, hình nào vừa có trục đối xứng vừa có tâm đối xứng:\par
\begin{itemize}\setlength{\itemsep}{1pt}\setlength{\parskip}{0pt}
  \item[] \textbf{A.} Hình tam giác đều
  \item[] \textbf{B.} Hình tam giác thường
  \item[] \textbf{C.} Hình thang cân (hai đáy khác nhau)
  \item[] \textbf{D.} Hình lục giác đều
\end{itemize}\vspace{2pt}
\textbf{\color{primary} Câu 15.} Xét đối tượng toán học trong Chương V: Trong các chữ cái in hoa (kiểu chữ không chân chuẩn): A, H, N, P, chữ cái nào có cả trục đối xứng và tâm đối xứng. Giá trị hoặc tính chất tương ứng là:\par
\hspace*{1.5em}\textbf{A.} Chữ F\quad\quad \textbf{B.} Chữ P\quad\quad \textbf{C.} Chữ O\quad\quad \textbf{D.} Chữ A\par\vspace{2pt}
\textbf{\color{primary} Câu 16.} Hoàn thành phép xác định cơ bản sau đây: Hình đa giác đều có 18 cạnh bằng nhau thì có bao nhiêu trục đối xứng:\par
\hspace*{1.5em}\textbf{A.} 1 trục đối xứng\quad\quad \textbf{B.} 18 trục đối xứng\quad\quad \textbf{C.} 36 trục đối xứng\quad\quad \textbf{D.} 17 trục đối xứng\par\vspace{2pt}
\textbf{\color{primary} Câu 17.} Trong các khẳng định hoặc giá trị dưới đây, đâu là đáp án đúng cho câu hỏi: Trong các hình phẳng sau đây, hình nào vừa có trục đối xứng vừa có tâm đối xứng?\par
\begin{itemize}\setlength{\itemsep}{1pt}\setlength{\parskip}{0pt}
  \item[] \textbf{A.} Hình tam giác đều
  \item[] \textbf{B.} Hình chữ nhật
  \item[] \textbf{C.} Hình tam giác thường
  \item[] \textbf{D.} Hình thang cân (hai đáy khác nhau)
\end{itemize}\vspace{2pt}
\textbf{\color{primary} Câu 18.} Bài tập khởi động Chương V: Trong các chữ cái in hoa (kiểu chữ không chân chuẩn): A, H, N, P, chữ cái nào có cả trục đối xứng và tâm đối xứng:\par
\hspace*{1.5em}\textbf{A.} Chữ F\quad\quad \textbf{B.} Chữ P\quad\quad \textbf{C.} Chữ I\quad\quad \textbf{D.} Chữ A\par\vspace{2pt}
\textbf{\color{primary} Câu 19.} Quan sát dữ kiện đề bài: Hình đa giác đều có 21 cạnh bằng nhau thì có bao nhiêu trục đối xứng. Hãy chọn phương án đúng:\par
\hspace*{1.5em}\textbf{A.} 42 trục đối xứng\quad\quad \textbf{B.} 1 trục đối xứng\quad\quad \textbf{C.} 21 trục đối xứng\quad\quad \textbf{D.} 20 trục đối xứng\par\vspace{2pt}
\textbf{\color{primary} Câu 20.} Sử dụng công thức hoặc tính chất đầu bài của Chương V để giải quyết yêu cầu: Trong các hình phẳng sau đây, hình nào vừa có trục đối xứng vừa có tâm đối xứng:\par
\begin{itemize}\setlength{\itemsep}{1pt}\setlength{\parskip}{0pt}
  \item[] \textbf{A.} Hình tam giác thường
  \item[] \textbf{B.} Hình tròn
  \item[] \textbf{C.} Hình tam giác đều
  \item[] \textbf{D.} Hình thang cân (hai đáy khác nhau)
\end{itemize}\vspace{2pt}
\subsection*{\color{accent} Phần II. Mức độ Trung bình (Thông hiểu -- Câu 21 đến Câu 40)}
\textbf{\color{primary} Câu 21.} Tổng số trục đối xứng của 2 hình vuông rời nhau và 1 hình chữ nhật (không phải hình vuông) rời nhau là bao nhiêu?\par
\hspace*{1.5em}\textbf{A.} 6 trục\quad\quad \textbf{B.} 12 trục\quad\quad \textbf{C.} 10 trục\quad\quad \textbf{D.} 9 trục\par\vspace{2pt}
\textbf{\color{primary} Câu 22.} Vận dụng quy tắc trong Chương V (Tính đối xứng của hình phẳng trong tự nhiên) để giải bài toán: Khẳng định nào dưới đây là ĐÚNG về tâm đối xứng của một đoạn thẳng AB có độ dài bằng 10 cm:\par
\begin{itemize}\setlength{\itemsep}{1pt}\setlength{\parskip}{0pt}
  \item[] \textbf{A.} Đoạn thẳng AB có vô số tâm đối xứng
  \item[] \textbf{B.} Tâm đối xứng là điểm A cách B một khoảng 10 cm
  \item[] \textbf{C.} Tâm đối xứng là trung điểm I của đoạn thẳng AB, cách A một khoảng 5 cm
  \item[] \textbf{D.} Đoạn thẳng AB không có tâm đối xứng
\end{itemize}\vspace{2pt}
\textbf{\color{primary} Câu 23.} Biến đổi và tính toán theo kiến thức Chương V: Tổng số trục đối xứng của 1 hình lục giác đều và 3 hình tam giác đều rời nhau là. Kết quả thu được là:\par
\hspace*{1.5em}\textbf{A.} 15 trục\quad\quad \textbf{B.} 18 trục\quad\quad \textbf{C.} 15 trục (1)\quad\quad \textbf{D.} 12 trục\par\vspace{2pt}
\textbf{\color{primary} Câu 24.} Xét bài toán mức độ thông hiểu sau: Tổng số trục đối xứng của 5 hình vuông rời nhau và 1 hình chữ nhật (không phải hình vuông) rời nhau là bao nhiêu. Đáp án đúng là:\par
\hspace*{1.5em}\textbf{A.} 22 trục\quad\quad \textbf{B.} 12 trục\quad\quad \textbf{C.} 24 trục\quad\quad \textbf{D.} 21 trục\par\vspace{2pt}
\textbf{\color{primary} Câu 25.} Để hoàn thành phiếu học tập Chương V, hãy xác định: Khẳng định nào dưới đây là ĐÚNG về tâm đối xứng của một đoạn thẳng AB có độ dài bằng 16 cm:\par
\begin{itemize}\setlength{\itemsep}{1pt}\setlength{\parskip}{0pt}
  \item[] \textbf{A.} Tâm đối xứng là trung điểm I của đoạn thẳng AB, cách A một khoảng 8 cm
  \item[] \textbf{B.} Đoạn thẳng AB có vô số tâm đối xứng
  \item[] \textbf{C.} Tâm đối xứng là điểm A cách B một khoảng 16 cm
  \item[] \textbf{D.} Đoạn thẳng AB không có tâm đối xứng
\end{itemize}\vspace{2pt}
\textbf{\color{primary} Câu 26.} Phân tích dữ kiện bài toán: Tổng số trục đối xứng của 1 hình lục giác đều và 6 hình tam giác đều rời nhau là. Giá trị cần tìm bằng:\par
\hspace*{1.5em}\textbf{A.} 24 trục\quad\quad \textbf{B.} 21 trục\quad\quad \textbf{C.} 30 trục\quad\quad \textbf{D.} 18 trục\par\vspace{2pt}
\textbf{\color{primary} Câu 27.} Áp dụng tính chất đã học ở Chương V vào bài tập sau: Tổng số trục đối xứng của 8 hình vuông rời nhau và 1 hình chữ nhật (không phải hình vuông) rời nhau là bao nhiêu:\par
\hspace*{1.5em}\textbf{A.} 36 trục\quad\quad \textbf{B.} 18 trục\quad\quad \textbf{C.} 34 trục\quad\quad \textbf{D.} 33 trục\par\vspace{2pt}
\textbf{\color{primary} Câu 28.} Thực hiện các bước suy luận hợp lý để giải quyết yêu cầu: Khẳng định nào dưới đây là ĐÚNG về tâm đối xứng của một đoạn thẳng AB có độ dài bằng 22 cm:\par
\begin{itemize}\setlength{\itemsep}{1pt}\setlength{\parskip}{0pt}
  \item[] \textbf{A.} Đoạn thẳng AB có vô số tâm đối xứng
  \item[] \textbf{B.} Tâm đối xứng là trung điểm I của đoạn thẳng AB, cách A một khoảng 11 cm
  \item[] \textbf{C.} Tâm đối xứng là điểm A cách B một khoảng 22 cm
  \item[] \textbf{D.} Đoạn thẳng AB không có tâm đối xứng
\end{itemize}\vspace{2pt}
\textbf{\color{primary} Câu 29.} Cho bài toán tự luyện (Tính đối xứng của hình phẳng trong tự nhiên): Tổng số trục đối xứng của 1 hình lục giác đều và 9 hình tam giác đều rời nhau là. Phương án nào sau đây chính xác?\par
\hspace*{1.5em}\textbf{A.} 42 trục\quad\quad \textbf{B.} 33 trục\quad\quad \textbf{C.} 30 trục\quad\quad \textbf{D.} 24 trục\par\vspace{2pt}
\textbf{\color{primary} Câu 30.} Kết hợp các phép biến đổi trong Chương V, hãy tìm kết quả của bài toán: Tổng số trục đối xứng của 11 hình vuông rời nhau và 1 hình chữ nhật (không phải hình vuông) rời nhau là bao nhiêu:\par
\hspace*{1.5em}\textbf{A.} 48 trục\quad\quad \textbf{B.} 24 trục\quad\quad \textbf{C.} 45 trục\quad\quad \textbf{D.} 46 trục\par\vspace{2pt}
\textbf{\color{primary} Câu 31.} Một bài kiểm tra định kỳ nội dung Chương V nêu câu hỏi: Khẳng định nào dưới đây là ĐÚNG về tâm đối xứng của một đoạn thẳng AB có độ dài bằng 28 cm. Hãy chọn đáp số đúng:\par
\begin{itemize}\setlength{\itemsep}{1pt}\setlength{\parskip}{0pt}
  \item[] \textbf{A.} Tâm đối xứng là trung điểm I của đoạn thẳng AB, cách A một khoảng 14 cm
  \item[] \textbf{B.} Đoạn thẳng AB không có tâm đối xứng
  \item[] \textbf{C.} Đoạn thẳng AB có vô số tâm đối xứng
  \item[] \textbf{D.} Tâm đối xứng là điểm A cách B một khoảng 28 cm
\end{itemize}\vspace{2pt}
\textbf{\color{primary} Câu 32.} Tìm lời giải đúng cho bài toán thông hiểu: Tổng số trục đối xứng của 1 hình lục giác đều và 12 hình tam giác đều rời nhau là:\par
\hspace*{1.5em}\textbf{A.} 42 trục\quad\quad \textbf{B.} 30 trục\quad\quad \textbf{C.} 54 trục\quad\quad \textbf{D.} 39 trục\par\vspace{2pt}
\textbf{\color{primary} Câu 33.} Sử dụng mối liên hệ giữa các đại lượng trong Chương V: Tổng số trục đối xứng của 14 hình vuông rời nhau và 1 hình chữ nhật (không phải hình vuông) rời nhau là bao nhiêu:\par
\hspace*{1.5em}\textbf{A.} 60 trục\quad\quad \textbf{B.} 57 trục\quad\quad \textbf{C.} 58 trục\quad\quad \textbf{D.} 30 trục\par\vspace{2pt}
\textbf{\color{primary} Câu 34.} Hãy tính toán và chọn phương án đúng đối với bài tập: Khẳng định nào dưới đây là ĐÚNG về tâm đối xứng của một đoạn thẳng AB có độ dài bằng 34 cm:\par
\begin{itemize}\setlength{\itemsep}{1pt}\setlength{\parskip}{0pt}
  \item[] \textbf{A.} Đoạn thẳng AB không có tâm đối xứng
  \item[] \textbf{B.} Tâm đối xứng là điểm A cách B một khoảng 34 cm
  \item[] \textbf{C.} Đoạn thẳng AB có vô số tâm đối xứng
  \item[] \textbf{D.} Tâm đối xứng là trung điểm I của đoạn thẳng AB, cách A một khoảng 17 cm
\end{itemize}\vspace{2pt}
\textbf{\color{primary} Câu 35.} Dựa trên các tính chất trọng tâm của Chương V, xét bài toán: Tổng số trục đối xứng của 1 hình lục giác đều và 15 hình tam giác đều rời nhau là:\par
\hspace*{1.5em}\textbf{A.} 36 trục\quad\quad \textbf{B.} 48 trục\quad\quad \textbf{C.} 66 trục\quad\quad \textbf{D.} 51 trục\par\vspace{2pt}
\textbf{\color{primary} Câu 36.} Thực hành giải toán (Tính đối xứng của hình phẳng trong tự nhiên): Tổng số trục đối xứng của 17 hình vuông rời nhau và 1 hình chữ nhật (không phải hình vuông) rời nhau là bao nhiêu. Đáp án nào dưới đây thỏa mãn?\par
\hspace*{1.5em}\textbf{A.} 69 trục\quad\quad \textbf{B.} 72 trục\quad\quad \textbf{C.} 36 trục\quad\quad \textbf{D.} 70 trục\par\vspace{2pt}
\textbf{\color{primary} Câu 37.} Kiểm tra kỹ năng biến đổi và suy luận qua bài toán: Khẳng định nào dưới đây là ĐÚNG về tâm đối xứng của một đoạn thẳng AB có độ dài bằng 40 cm:\par
\begin{itemize}\setlength{\itemsep}{1pt}\setlength{\parskip}{0pt}
  \item[] \textbf{A.} Đoạn thẳng AB không có tâm đối xứng
  \item[] \textbf{B.} Tâm đối xứng là trung điểm I của đoạn thẳng AB, cách A một khoảng 20 cm
  \item[] \textbf{C.} Đoạn thẳng AB có vô số tâm đối xứng
  \item[] \textbf{D.} Tâm đối xứng là điểm A cách B một khoảng 40 cm
\end{itemize}\vspace{2pt}
\textbf{\color{primary} Câu 38.} Cho giả thiết thuộc phạm vi Chương V: Tổng số trục đối xứng của 1 hình lục giác đều và 18 hình tam giác đều rời nhau là. Kết luận đúng là:\par
\hspace*{1.5em}\textbf{A.} 78 trục\quad\quad \textbf{B.} 60 trục\quad\quad \textbf{C.} 42 trục\quad\quad \textbf{D.} 57 trục\par\vspace{2pt}
\textbf{\color{primary} Câu 39.} Rút gọn hoặc tính giá trị theo yêu cầu sau của Chương V: Tổng số trục đối xứng của 20 hình vuông rời nhau và 1 hình chữ nhật (không phải hình vuông) rời nhau là bao nhiêu:\par
\hspace*{1.5em}\textbf{A.} 84 trục\quad\quad \textbf{B.} 42 trục\quad\quad \textbf{C.} 81 trục\quad\quad \textbf{D.} 82 trục\par\vspace{2pt}
\textbf{\color{primary} Câu 40.} Xác định phương án đúng sau khi thực hiện phép tính / suy luận: Khẳng định nào dưới đây là ĐÚNG về tâm đối xứng của một đoạn thẳng AB có độ dài bằng 46 cm:\par
\begin{itemize}\setlength{\itemsep}{1pt}\setlength{\parskip}{0pt}
  \item[] \textbf{A.} Đoạn thẳng AB có vô số tâm đối xứng
  \item[] \textbf{B.} Tâm đối xứng là trung điểm I của đoạn thẳng AB, cách A một khoảng 23 cm
  \item[] \textbf{C.} Đoạn thẳng AB không có tâm đối xứng
  \item[] \textbf{D.} Tâm đối xứng là điểm A cách B một khoảng 46 cm
\end{itemize}\vspace{2pt}
\subsection*{\color{accent} Phần III. Mức độ Khó (Vận dụng \& Vận dụng cao -- Câu 41 đến Câu 50)}
\textbf{\color{primary} Câu 41.} Trên trục số nguyên, điểm A biểu diễn số nguyên -10 và điểm B biểu diễn số nguyên 6. Điểm M là tâm đối xứng của đoạn thẳng AB thì M biểu diễn số nguyên nào?\par
\hspace*{1.5em}\textbf{A.} -2\quad\quad \textbf{B.} 2\quad\quad \textbf{C.} -2 (1)\quad\quad \textbf{D.} 0\par\vspace{2pt}
\textbf{\color{primary} Câu 42.} Bài toán vận dụng (Chương V – Tính đối xứng của hình phẳng trong tự nhiên): Trên trục số nguyên, điểm P biểu diễn số 4. Điểm Q đối xứng với điểm P qua điểm I biểu diễn số -2 sẽ biểu diễn số nguyên nào. Kết quả đúng là:\par
\hspace*{1.5em}\textbf{A.} -6\quad\quad \textbf{B.} -8\quad\quad \textbf{C.} 8\quad\quad \textbf{D.} -4\par\vspace{2pt}
\textbf{\color{primary} Câu 43.} Vận dụng kiến thức nâng cao trong phạm vi Chương V để giải bài toán: Trên trục số nguyên, điểm A biểu diễn số nguyên -14 và điểm B biểu diễn số nguyên 10. Điểm M là tâm đối xứng của đoạn thẳng AB thì M biểu diễn số nguyên nào:\par
\hspace*{1.5em}\textbf{A.} 0\quad\quad \textbf{B.} -4\quad\quad \textbf{C.} 2\quad\quad \textbf{D.} -2\par\vspace{2pt}
\textbf{\color{primary} Câu 44.} Phân tích và giải bài toán tổng hợp sau của Chương V: Trên trục số nguyên, điểm P biểu diễn số 6. Điểm Q đối xứng với điểm P qua điểm I biểu diễn số -4 sẽ biểu diễn số nguyên nào:\par
\hspace*{1.5em}\textbf{A.} 14\quad\quad \textbf{B.} -10\quad\quad \textbf{C.} -6\quad\quad \textbf{D.} -14\par\vspace{2pt}
\textbf{\color{primary} Câu 45.} Thử thách tư duy trong phạm vi kiến thức Chương V: Trên trục số nguyên, điểm A biểu diễn số nguyên -18 và điểm B biểu diễn số nguyên 14. Điểm M là tâm đối xứng của đoạn thẳng AB thì M biểu diễn số nguyên nào. Đáp số chính xác là:\par
\hspace*{1.5em}\textbf{A.} 0\quad\quad \textbf{B.} -2\quad\quad \textbf{C.} -6\quad\quad \textbf{D.} 2\par\vspace{2pt}
\textbf{\color{primary} Câu 46.} Sử dụng phương pháp giải toán của Chương V để tìm đáp án cho bài toán: Trên trục số nguyên, điểm P biểu diễn số 8. Điểm Q đối xứng với điểm P qua điểm I biểu diễn số -6 sẽ biểu diễn số nguyên nào:\par
\hspace*{1.5em}\textbf{A.} -14\quad\quad \textbf{B.} 20\quad\quad \textbf{C.} -20\quad\quad \textbf{D.} -8\par\vspace{2pt}
\textbf{\color{primary} Câu 47.} Xét bài toán vận dụng thực tiễn / nâng cao thuộc Chương V: Trên trục số nguyên, điểm A biểu diễn số nguyên -22 và điểm B biểu diễn số nguyên 18. Điểm M là tâm đối xứng của đoạn thẳng AB thì M biểu diễn số nguyên nào:\par
\hspace*{1.5em}\textbf{A.} -2\quad\quad \textbf{B.} 0\quad\quad \textbf{C.} 2\quad\quad \textbf{D.} -8\par\vspace{2pt}
\textbf{\color{primary} Câu 48.} Tìm kết quả cuối cùng của bài toán phân loại học sinh khá giỏi (Chương V): Trên trục số nguyên, điểm P biểu diễn số 10. Điểm Q đối xứng với điểm P qua điểm I biểu diễn số -8 sẽ biểu diễn số nguyên nào:\par
\hspace*{1.5em}\textbf{A.} -10\quad\quad \textbf{B.} -26\quad\quad \textbf{C.} -18\quad\quad \textbf{D.} 26\par\vspace{2pt}
\textbf{\color{primary} Câu 49.} Kết hợp các tính chất đã học trong Chương V, hãy giải quyết bài toán: Trên trục số nguyên, điểm A biểu diễn số nguyên -26 và điểm B biểu diễn số nguyên 22. Điểm M là tâm đối xứng của đoạn thẳng AB thì M biểu diễn số nguyên nào:\par
\hspace*{1.5em}\textbf{A.} -10\quad\quad \textbf{B.} 2\quad\quad \textbf{C.} 0\quad\quad \textbf{D.} -2\par\vspace{2pt}
\textbf{\color{primary} Câu 50.} Câu hỏi vận dụng cao cuối Chương V: Trên trục số nguyên, điểm P biểu diễn số 12. Điểm Q đối xứng với điểm P qua điểm I biểu diễn số -10 sẽ biểu diễn số nguyên nào. Phương án đúng là:\par
\hspace*{1.5em}\textbf{A.} 32\quad\quad \textbf{B.} -22\quad\quad \textbf{C.} -32\quad\quad \textbf{D.} -12\par\vspace{2pt}
\subsection*{\color{primary} Bảng đáp án Chương V: Tính đối xứng của hình phẳng trong tự nhiên}
\begin{center}
\renewcommand{\arraystretch}{1.25}
\begin{tabularx}{\textwidth}{*{10}{>{\centering\arraybackslash}X}}
\toprule
\textbf{\color{primary} Cột 1} & \textbf{\color{primary} Cột 2} & \textbf{\color{primary} Cột 3} & \textbf{\color{primary} Cột 4} & \textbf{\color{primary} Cột 5} & \textbf{\color{primary} Cột 6} & \textbf{\color{primary} Cột 7} & \textbf{\color{primary} Cột 8} & \textbf{\color{primary} Cột 9} & \textbf{\color{primary} Cột 10} \\
\midrule
\textbf{1.C} & \textbf{2.C} & \textbf{3.A} & \textbf{4.B} & \textbf{5.A} & \textbf{6.C} & \textbf{7.C} & \textbf{8.A} & \textbf{9.D} & \textbf{10.C} \\
\rowcolor{rowtint} \textbf{11.C} & \textbf{12.C} & \textbf{13.A} & \textbf{14.D} & \textbf{15.C} & \textbf{16.B} & \textbf{17.B} & \textbf{18.C} & \textbf{19.C} & \textbf{20.B} \\
\textbf{21.C} & \textbf{22.C} & \textbf{23.A} & \textbf{24.A} & \textbf{25.A} & \textbf{26.A} & \textbf{27.C} & \textbf{28.B} & \textbf{29.B} & \textbf{30.D} \\
\rowcolor{rowtint} \textbf{31.A} & \textbf{32.A} & \textbf{33.C} & \textbf{34.D} & \textbf{35.D} & \textbf{36.D} & \textbf{37.B} & \textbf{38.B} & \textbf{39.D} & \textbf{40.B} \\
\textbf{41.A} & \textbf{42.B} & \textbf{43.D} & \textbf{44.D} & \textbf{45.B} & \textbf{46.C} & \textbf{47.A} & \textbf{48.B} & \textbf{49.D} & \textbf{50.C} \\
\bottomrule
\end{tabularx}
\end{center}\vspace{10pt}
\section*{\color{primary} CHƯƠNG VI: PHÂN SỐ}
\textit{Trọng tâm kiến thức SGK: Mở rộng phân số, phân số bằng nhau, so sánh phân số, hỗn số dương, các phép tính cộng, trừ, nhân, chia phân số và hai bài toán về phân số.}\par\vspace{4pt}
\subsection*{\color{accent} Phần I. Mức độ Dễ (Nhận biết -- Câu 1 đến Câu 20)}
\textbf{\color{primary} Câu 1.} Rút gọn phân số 6/10 về phân số tối giản ta được:\par
\hspace*{1.5em}\textbf{A.} 2/5\quad\quad \textbf{B.} 3/4\quad\quad \textbf{C.} 5/3\quad\quad \textbf{D.} 3/5\par\vspace{2pt}
\textbf{\color{primary} Câu 2.} Thực hiện yêu cầu nhận biết sau: Số đối của phân số -3/7 là.\par
\hspace*{1.5em}\textbf{A.} -7/3\quad\quad \textbf{B.} -3/7\quad\quad \textbf{C.} 3/7\quad\quad \textbf{D.} 7/3\par\vspace{2pt}
\textbf{\color{primary} Câu 3.} Trong bài học thuộc Chương VI, xét nội dung: Viết hỗn số 4 và 3/7 dưới dạng phân số ta được. Đáp án đúng là:\par
\hspace*{1.5em}\textbf{A.} 25/7\quad\quad \textbf{B.} 28/7\quad\quad \textbf{C.} 19/7\quad\quad \textbf{D.} 31/7\par\vspace{2pt}
\textbf{\color{primary} Câu 4.} Dựa vào kiến thức chuẩn SGK của Chương VI (Phân số), rút gọn phân số 15/25 về phân số tối giản ta được:\par
\hspace*{1.5em}\textbf{A.} 2/5\quad\quad \textbf{B.} 3/5\quad\quad \textbf{C.} 5/3\quad\quad \textbf{D.} 3/4\par\vspace{2pt}
\textbf{\color{primary} Câu 5.} Chọn kết quả đúng nhất khi xét: Số đối của phân số -6/13 là:\par
\hspace*{1.5em}\textbf{A.} -13/6\quad\quad \textbf{B.} 13/6\quad\quad \textbf{C.} -6/13\quad\quad \textbf{D.} 6/13\par\vspace{2pt}
\textbf{\color{primary} Câu 6.} Khi làm bài tập cơ bản về phân số, một học sinh gặp câu hỏi: Viết hỗn số 7 và 3/7 dưới dạng phân số ta được. Kết quả nào dưới đây đúng?\par
\hspace*{1.5em}\textbf{A.} 46/7\quad\quad \textbf{B.} 49/7\quad\quad \textbf{C.} 52/7\quad\quad \textbf{D.} 28/7\par\vspace{2pt}
\textbf{\color{primary} Câu 7.} Xác định giá trị hoặc khẳng định đúng cho bài toán sau: Rút gọn phân số 24/40 về phân số tối giản ta được:\par
\hspace*{1.5em}\textbf{A.} 3/5\quad\quad \textbf{B.} 5/3\quad\quad \textbf{C.} 3/4\quad\quad \textbf{D.} 2/5\par\vspace{2pt}
\textbf{\color{primary} Câu 8.} Áp dụng trực tiếp định nghĩa và quy tắc trong Chương VI: Số đối của phân số -9/19 là:\par
\hspace*{1.5em}\textbf{A.} 9/19\quad\quad \textbf{B.} -9/19\quad\quad \textbf{C.} 19/9\quad\quad \textbf{D.} -19/9\par\vspace{2pt}
\textbf{\color{primary} Câu 9.} Kiểm tra nhanh kiến thức nền tảng: Viết hỗn số 10 và 3/7 dưới dạng phân số ta được:\par
\hspace*{1.5em}\textbf{A.} 37/7\quad\quad \textbf{B.} 73/7\quad\quad \textbf{C.} 67/7\quad\quad \textbf{D.} 70/7\par\vspace{2pt}
\textbf{\color{primary} Câu 10.} Cho dữ kiện nhận biết sau đây: Rút gọn phân số 33/55 về phân số tối giản ta được. Phương án chính xác là:\par
\hspace*{1.5em}\textbf{A.} 3/4\quad\quad \textbf{B.} 2/5\quad\quad \textbf{C.} 3/5\quad\quad \textbf{D.} 5/3\par\vspace{2pt}
\textbf{\color{primary} Câu 11.} Hãy chỉ ra kết quả đúng trong bốn lựa chọn A, B, C, D đối với yêu cầu: Số đối của phân số -12/25 là:\par
\hspace*{1.5em}\textbf{A.} -25/12\quad\quad \textbf{B.} 12/25\quad\quad \textbf{C.} 25/12\quad\quad \textbf{D.} -12/25\par\vspace{2pt}
\textbf{\color{primary} Câu 12.} Câu hỏi ôn tập lý thuyết và kỹ năng cơ bản (Chương VI): Viết hỗn số 13 và 3/7 dưới dạng phân số ta được:\par
\hspace*{1.5em}\textbf{A.} 94/7\quad\quad \textbf{B.} 46/7\quad\quad \textbf{C.} 88/7\quad\quad \textbf{D.} 91/7\par\vspace{2pt}
\textbf{\color{primary} Câu 13.} Theo nội dung sách giáo khoa Kết nối tri thức với cuộc sống (Chương VI), rút gọn phân số 42/70 về phân số tối giản ta được:\par
\hspace*{1.5em}\textbf{A.} 3/5\quad\quad \textbf{B.} 2/5\quad\quad \textbf{C.} 5/3\quad\quad \textbf{D.} 3/4\par\vspace{2pt}
\textbf{\color{primary} Câu 14.} Tìm đáp số đúng của bài toán cơ bản sau: Số đối của phân số -15/31 là:\par
\hspace*{1.5em}\textbf{A.} 15/31\quad\quad \textbf{B.} -31/15\quad\quad \textbf{C.} -15/31\quad\quad \textbf{D.} 31/15\par\vspace{2pt}
\textbf{\color{primary} Câu 15.} Xét đối tượng toán học trong Chương VI: Viết hỗn số 16 và 3/7 dưới dạng phân số ta được. Giá trị hoặc tính chất tương ứng là:\par
\hspace*{1.5em}\textbf{A.} 55/7\quad\quad \textbf{B.} 115/7\quad\quad \textbf{C.} 112/7\quad\quad \textbf{D.} 109/7\par\vspace{2pt}
\textbf{\color{primary} Câu 16.} Hoàn thành phép xác định cơ bản sau đây: Rút gọn phân số 51/85 về phân số tối giản ta được:\par
\hspace*{1.5em}\textbf{A.} 3/5\quad\quad \textbf{B.} 2/5\quad\quad \textbf{C.} 3/4\quad\quad \textbf{D.} 5/3\par\vspace{2pt}
\textbf{\color{primary} Câu 17.} Trong các khẳng định hoặc giá trị dưới đây, đâu là đáp án đúng cho câu hỏi: Số đối của phân số -18/37 là?\par
\hspace*{1.5em}\textbf{A.} -37/18\quad\quad \textbf{B.} 37/18\quad\quad \textbf{C.} 18/37\quad\quad \textbf{D.} -18/37\par\vspace{2pt}
\textbf{\color{primary} Câu 18.} Bài tập khởi động Chương VI: Viết hỗn số 19 và 3/7 dưới dạng phân số ta được:\par
\hspace*{1.5em}\textbf{A.} 64/7\quad\quad \textbf{B.} 130/7\quad\quad \textbf{C.} 133/7\quad\quad \textbf{D.} 136/7\par\vspace{2pt}
\textbf{\color{primary} Câu 19.} Quan sát dữ kiện đề bài: Rút gọn phân số 60/100 về phân số tối giản ta được. Hãy chọn phương án đúng:\par
\hspace*{1.5em}\textbf{A.} 5/3\quad\quad \textbf{B.} 3/4\quad\quad \textbf{C.} 2/5\quad\quad \textbf{D.} 3/5\par\vspace{2pt}
\textbf{\color{primary} Câu 20.} Sử dụng công thức hoặc tính chất đầu bài của Chương VI để giải quyết yêu cầu: Số đối của phân số -21/43 là:\par
\hspace*{1.5em}\textbf{A.} 21/43\quad\quad \textbf{B.} 43/21\quad\quad \textbf{C.} -21/43\quad\quad \textbf{D.} -43/21\par\vspace{2pt}
\subsection*{\color{accent} Phần II. Mức độ Trung bình (Thông hiểu -- Câu 21 đến Câu 40)}
\textbf{\color{primary} Câu 21.} Tìm x, biết: x - 2/3 = 5/3.:\par
\hspace*{1.5em}\textbf{A.} x = 3/3\quad\quad \textbf{B.} x = 7/3\quad\quad \textbf{C.} x = 10/3\quad\quad \textbf{D.} x = -3/3\par\vspace{2pt}
\textbf{\color{primary} Câu 22.} Vận dụng quy tắc trong Chương VI (Phân số) để giải bài toán: Tính giá trị của 3/5 của 45:\par
\hspace*{1.5em}\textbf{A.} 36\quad\quad \textbf{B.} 27\quad\quad \textbf{C.} 75\quad\quad \textbf{D.} 9\par\vspace{2pt}
\textbf{\color{primary} Câu 23.} Biến đổi và tính toán theo kiến thức Chương VI: Biết 2/3 của một số bằng 30. Số đó là. Kết quả thu được là:\par
\hspace*{1.5em}\textbf{A.} 90\quad\quad \textbf{B.} 20\quad\quad \textbf{C.} 36\quad\quad \textbf{D.} 45\par\vspace{2pt}
\textbf{\color{primary} Câu 24.} Xét bài toán mức độ thông hiểu sau: Tìm x, biết: x - 2/9 = 5/9.. Đáp án đúng là:\par
\hspace*{1.5em}\textbf{A.} x = 7/9\quad\quad \textbf{B.} x = 3/9\quad\quad \textbf{C.} x = -3/9\quad\quad \textbf{D.} x = 10/9\par\vspace{2pt}
\textbf{\color{primary} Câu 25.} Để hoàn thành phiếu học tập Chương VI, hãy xác định: Tính giá trị của 3/5 của 90:\par
\hspace*{1.5em}\textbf{A.} 54\quad\quad \textbf{B.} 150\quad\quad \textbf{C.} 72\quad\quad \textbf{D.} 18\par\vspace{2pt}
\textbf{\color{primary} Câu 26.} Phân tích dữ kiện bài toán: Biết 2/3 của một số bằng 60. Số đó là. Giá trị cần tìm bằng:\par
\hspace*{1.5em}\textbf{A.} 72\quad\quad \textbf{B.} 40\quad\quad \textbf{C.} 90\quad\quad \textbf{D.} 180\par\vspace{2pt}
\textbf{\color{primary} Câu 27.} Áp dụng tính chất đã học ở Chương VI vào bài tập sau: Tìm x, biết: x - 2/15 = 5/15.:\par
\hspace*{1.5em}\textbf{A.} x = -3/15\quad\quad \textbf{B.} x = 10/15\quad\quad \textbf{C.} x = 7/15\quad\quad \textbf{D.} x = 3/15\par\vspace{2pt}
\textbf{\color{primary} Câu 28.} Thực hiện các bước suy luận hợp lý để giải quyết yêu cầu: Tính giá trị của 3/5 của 135:\par
\hspace*{1.5em}\textbf{A.} 225\quad\quad \textbf{B.} 108\quad\quad \textbf{C.} 81\quad\quad \textbf{D.} 27\par\vspace{2pt}
\textbf{\color{primary} Câu 29.} Cho bài toán tự luyện (Phân số): Biết 2/3 của một số bằng 90. Số đó là. Phương án nào sau đây chính xác?\par
\hspace*{1.5em}\textbf{A.} 135\quad\quad \textbf{B.} 60\quad\quad \textbf{C.} 108\quad\quad \textbf{D.} 270\par\vspace{2pt}
\textbf{\color{primary} Câu 30.} Kết hợp các phép biến đổi trong Chương VI, hãy tìm kết quả của bài toán: Tìm x, biết: x - 2/21 = 5/21.:\par
\hspace*{1.5em}\textbf{A.} x = -3/21\quad\quad \textbf{B.} x = 10/21\quad\quad \textbf{C.} x = 7/21\quad\quad \textbf{D.} x = 3/21\par\vspace{2pt}
\textbf{\color{primary} Câu 31.} Một bài kiểm tra định kỳ nội dung Chương VI nêu câu hỏi: Tính giá trị của 3/5 của 180. Hãy chọn đáp số đúng:\par
\hspace*{1.5em}\textbf{A.} 36\quad\quad \textbf{B.} 144\quad\quad \textbf{C.} 108\quad\quad \textbf{D.} 300\par\vspace{2pt}
\textbf{\color{primary} Câu 32.} Tìm lời giải đúng cho bài toán thông hiểu: Biết 2/3 của một số bằng 120. Số đó là:\par
\hspace*{1.5em}\textbf{A.} 80\quad\quad \textbf{B.} 360\quad\quad \textbf{C.} 180\quad\quad \textbf{D.} 144\par\vspace{2pt}
\textbf{\color{primary} Câu 33.} Sử dụng mối liên hệ giữa các đại lượng trong Chương VI: Tìm x, biết: x - 2/27 = 5/27.:\par
\hspace*{1.5em}\textbf{A.} x = 10/27\quad\quad \textbf{B.} x = 7/27\quad\quad \textbf{C.} x = -3/27\quad\quad \textbf{D.} x = 3/27\par\vspace{2pt}
\textbf{\color{primary} Câu 34.} Hãy tính toán và chọn phương án đúng đối với bài tập: Tính giá trị của 3/5 của 225:\par
\hspace*{1.5em}\textbf{A.} 135\quad\quad \textbf{B.} 180\quad\quad \textbf{C.} 375\quad\quad \textbf{D.} 45\par\vspace{2pt}
\textbf{\color{primary} Câu 35.} Dựa trên các tính chất trọng tâm của Chương VI, xét bài toán: Biết 2/3 của một số bằng 150. Số đó là:\par
\hspace*{1.5em}\textbf{A.} 225\quad\quad \textbf{B.} 100\quad\quad \textbf{C.} 450\quad\quad \textbf{D.} 180\par\vspace{2pt}
\textbf{\color{primary} Câu 36.} Thực hành giải toán (Phân số): Tìm x, biết: x - 2/33 = 5/33.. Đáp án nào dưới đây thỏa mãn?\par
\hspace*{1.5em}\textbf{A.} x = 3/33\quad\quad \textbf{B.} x = 10/33\quad\quad \textbf{C.} x = 7/33\quad\quad \textbf{D.} x = -3/33\par\vspace{2pt}
\textbf{\color{primary} Câu 37.} Kiểm tra kỹ năng biến đổi và suy luận qua bài toán: Tính giá trị của 3/5 của 270:\par
\hspace*{1.5em}\textbf{A.} 54\quad\quad \textbf{B.} 162\quad\quad \textbf{C.} 450\quad\quad \textbf{D.} 216\par\vspace{2pt}
\textbf{\color{primary} Câu 38.} Cho giả thiết thuộc phạm vi Chương VI: Biết 2/3 của một số bằng 180. Số đó là. Kết luận đúng là:\par
\hspace*{1.5em}\textbf{A.} 120\quad\quad \textbf{B.} 216\quad\quad \textbf{C.} 540\quad\quad \textbf{D.} 270\par\vspace{2pt}
\textbf{\color{primary} Câu 39.} Rút gọn hoặc tính giá trị theo yêu cầu sau của Chương VI: Tìm x, biết: x - 2/39 = 5/39.:\par
\hspace*{1.5em}\textbf{A.} x = 3/39\quad\quad \textbf{B.} x = 7/39\quad\quad \textbf{C.} x = 10/39\quad\quad \textbf{D.} x = -3/39\par\vspace{2pt}
\textbf{\color{primary} Câu 40.} Xác định phương án đúng sau khi thực hiện phép tính / suy luận: Tính giá trị của 3/5 của 315:\par
\hspace*{1.5em}\textbf{A.} 525\quad\quad \textbf{B.} 63\quad\quad \textbf{C.} 252\quad\quad \textbf{D.} 189\par\vspace{2pt}
\subsection*{\color{accent} Phần III. Mức độ Khó (Vận dụng \& Vận dụng cao -- Câu 41 đến Câu 50)}
\textbf{\color{primary} Câu 41.} Tính tổng dãy phân số S = 1/(1$\cdot$2) + 1/(2$\cdot$3) + 1/(3$\cdot$4) + ... + 1/(10$\cdot$11).:\par
\hspace*{1.5em}\textbf{A.} 11/12\quad\quad \textbf{B.} 10/11\quad\quad \textbf{C.} 9/11\quad\quad \textbf{D.} 1/11\par\vspace{2pt}
\textbf{\color{primary} Câu 42.} Bài toán vận dụng (Chương VI – Phân số): Một lớp học có 45 học sinh, trong đó số học sinh giỏi chiếm 1/5 cả lớp, số học sinh khá chiếm 3/4 số học sinh còn lại. Số học sinh khá của lớp đó là. Kết quả đúng là:\par
\hspace*{1.5em}\textbf{A.} 9 học sinh\quad\quad \textbf{B.} 27 học sinh (1)\quad\quad \textbf{C.} 27 học sinh\quad\quad \textbf{D.} 22 học sinh\par\vspace{2pt}
\textbf{\color{primary} Câu 43.} Vận dụng kiến thức nâng cao trong phạm vi Chương VI để giải bài toán: Tính tổng dãy phân số S = 1/(1$\cdot$2) + 1/(2$\cdot$3) + 1/(3$\cdot$4) + ... + 1/(12$\cdot$13).:\par
\hspace*{1.5em}\textbf{A.} 1/13\quad\quad \textbf{B.} 12/13\quad\quad \textbf{C.} 13/14\quad\quad \textbf{D.} 11/13\par\vspace{2pt}
\textbf{\color{primary} Câu 44.} Phân tích và giải bài toán tổng hợp sau của Chương VI: Một lớp học có 40 học sinh, trong đó số học sinh giỏi chiếm 1/5 cả lớp, số học sinh khá chiếm 3/4 số học sinh còn lại. Số học sinh khá của lớp đó là:\par
\hspace*{1.5em}\textbf{A.} 20 học sinh\quad\quad \textbf{B.} 24 học sinh (1)\quad\quad \textbf{C.} 8 học sinh\quad\quad \textbf{D.} 24 học sinh\par\vspace{2pt}
\textbf{\color{primary} Câu 45.} Thử thách tư duy trong phạm vi kiến thức Chương VI: Tính tổng dãy phân số S = 1/(1$\cdot$2) + 1/(2$\cdot$3) + 1/(3$\cdot$4) + ... + 1/(14$\cdot$15).. Đáp số chính xác là:\par
\hspace*{1.5em}\textbf{A.} 15/16\quad\quad \textbf{B.} 1/15\quad\quad \textbf{C.} 13/15\quad\quad \textbf{D.} 14/15\par\vspace{2pt}
\textbf{\color{primary} Câu 46.} Sử dụng phương pháp giải toán của Chương VI để tìm đáp án cho bài toán: Một lớp học có 50 học sinh, trong đó số học sinh giỏi chiếm 1/5 cả lớp, số học sinh khá chiếm 3/4 số học sinh còn lại. Số học sinh khá của lớp đó là:\par
\hspace*{1.5em}\textbf{A.} 25 học sinh\quad\quad \textbf{B.} 30 học sinh\quad\quad \textbf{C.} 10 học sinh\quad\quad \textbf{D.} 30 học sinh (1)\par\vspace{2pt}
\textbf{\color{primary} Câu 47.} Xét bài toán vận dụng thực tiễn / nâng cao thuộc Chương VI: Tính tổng dãy phân số S = 1/(1$\cdot$2) + 1/(2$\cdot$3) + 1/(3$\cdot$4) + ... + 1/(16$\cdot$17).:\par
\hspace*{1.5em}\textbf{A.} 15/17\quad\quad \textbf{B.} 1/17\quad\quad \textbf{C.} 16/17\quad\quad \textbf{D.} 17/18\par\vspace{2pt}
\textbf{\color{primary} Câu 48.} Tìm kết quả cuối cùng của bài toán phân loại học sinh khá giỏi (Chương VI): Một lớp học có 45 học sinh, trong đó số học sinh giỏi chiếm 1/5 cả lớp, số học sinh khá chiếm 3/4 số học sinh còn lại. Số học sinh khá của lớp đó là:\par
\hspace*{1.5em}\textbf{A.} 9 học sinh\quad\quad \textbf{B.} 27 học sinh\quad\quad \textbf{C.} 27 học sinh (1)\quad\quad \textbf{D.} 22 học sinh\par\vspace{2pt}
\textbf{\color{primary} Câu 49.} Kết hợp các tính chất đã học trong Chương VI, hãy giải quyết bài toán: Tính tổng dãy phân số S = 1/(1$\cdot$2) + 1/(2$\cdot$3) + 1/(3$\cdot$4) + ... + 1/(18$\cdot$19).:\par
\hspace*{1.5em}\textbf{A.} 18/19\quad\quad \textbf{B.} 17/19\quad\quad \textbf{C.} 1/19\quad\quad \textbf{D.} 19/20\par\vspace{2pt}
\textbf{\color{primary} Câu 50.} Câu hỏi vận dụng cao cuối Chương VI: Một lớp học có 40 học sinh, trong đó số học sinh giỏi chiếm 1/5 cả lớp, số học sinh khá chiếm 3/4 số học sinh còn lại. Số học sinh khá của lớp đó là. Phương án đúng là:\par
\hspace*{1.5em}\textbf{A.} 24 học sinh\quad\quad \textbf{B.} 8 học sinh\quad\quad \textbf{C.} 24 học sinh (1)\quad\quad \textbf{D.} 20 học sinh\par\vspace{2pt}
\subsection*{\color{primary} Bảng đáp án Chương VI: Phân số}
\begin{center}
\renewcommand{\arraystretch}{1.25}
\begin{tabularx}{\textwidth}{*{10}{>{\centering\arraybackslash}X}}
\toprule
\textbf{\color{primary} Cột 1} & \textbf{\color{primary} Cột 2} & \textbf{\color{primary} Cột 3} & \textbf{\color{primary} Cột 4} & \textbf{\color{primary} Cột 5} & \textbf{\color{primary} Cột 6} & \textbf{\color{primary} Cột 7} & \textbf{\color{primary} Cột 8} & \textbf{\color{primary} Cột 9} & \textbf{\color{primary} Cột 10} \\
\midrule
\textbf{1.D} & \textbf{2.C} & \textbf{3.D} & \textbf{4.B} & \textbf{5.D} & \textbf{6.C} & \textbf{7.A} & \textbf{8.A} & \textbf{9.B} & \textbf{10.C} \\
\rowcolor{rowtint} \textbf{11.B} & \textbf{12.A} & \textbf{13.A} & \textbf{14.A} & \textbf{15.B} & \textbf{16.A} & \textbf{17.C} & \textbf{18.D} & \textbf{19.D} & \textbf{20.A} \\
\textbf{21.B} & \textbf{22.B} & \textbf{23.D} & \textbf{24.A} & \textbf{25.A} & \textbf{26.C} & \textbf{27.C} & \textbf{28.C} & \textbf{29.A} & \textbf{30.C} \\
\rowcolor{rowtint} \textbf{31.C} & \textbf{32.C} & \textbf{33.B} & \textbf{34.A} & \textbf{35.A} & \textbf{36.C} & \textbf{37.B} & \textbf{38.D} & \textbf{39.B} & \textbf{40.D} \\
\textbf{41.B} & \textbf{42.C} & \textbf{43.B} & \textbf{44.D} & \textbf{45.D} & \textbf{46.B} & \textbf{47.C} & \textbf{48.B} & \textbf{49.A} & \textbf{50.A} \\
\bottomrule
\end{tabularx}
\end{center}\vspace{10pt}
\section*{\color{primary} CHƯƠNG VII: SỐ THẬP PHÂN}
\textit{Trọng tâm kiến thức SGK: Số thập phân âm, số đối, so sánh số thập phân, các phép tính với số thập phân, làm tròn và ước lượng, tỉ số và tỉ số phần trăm.}\par\vspace{4pt}
\subsection*{\color{accent} Phần I. Mức độ Dễ (Nhận biết -- Câu 1 đến Câu 20)}
\textbf{\color{primary} Câu 1.} Phân số thập phân -15/10 được viết dưới dạng số thập phân là:\par
\hspace*{1.5em}\textbf{A.} -1.5\quad\quad \textbf{B.} -15.0\quad\quad \textbf{C.} 1.5\quad\quad \textbf{D.} -0.15\par\vspace{2pt}
\textbf{\color{primary} Câu 2.} Thực hiện yêu cầu nhận biết sau: Số đối của số thập phân 13,35 là.\par
\hspace*{1.5em}\textbf{A.} 0,1335\quad\quad \textbf{B.} -13,35\quad\quad \textbf{C.} 13,53\quad\quad \textbf{D.} -13,53\par\vspace{2pt}
\textbf{\color{primary} Câu 3.} Trong bài học thuộc Chương VII, xét nội dung: Làm tròn số thập phân 22,682 đến hàng phần mười ta được. Đáp án đúng là:\par
\hspace*{1.5em}\textbf{A.} 22,7\quad\quad \textbf{B.} 22,6\quad\quad \textbf{C.} 22,68\quad\quad \textbf{D.} 23\par\vspace{2pt}
\textbf{\color{primary} Câu 4.} Dựa vào kiến thức chuẩn SGK của Chương VII (Số thập phân), phân số thập phân -27/10 được viết dưới dạng số thập phân là:\par
\hspace*{1.5em}\textbf{A.} -0.27\quad\quad \textbf{B.} -2.7\quad\quad \textbf{C.} -27.0\quad\quad \textbf{D.} 2.7\par\vspace{2pt}
\textbf{\color{primary} Câu 5.} Chọn kết quả đúng nhất khi xét: Số đối của số thập phân 16,35 là:\par
\hspace*{1.5em}\textbf{A.} -16,35\quad\quad \textbf{B.} 16,53\quad\quad \textbf{C.} -16,53\quad\quad \textbf{D.} 0,1635\par\vspace{2pt}
\textbf{\color{primary} Câu 6.} Khi làm bài tập cơ bản về số thập phân, một học sinh gặp câu hỏi: Làm tròn số thập phân 25,682 đến hàng phần mười ta được. Kết quả nào dưới đây đúng?\par
\hspace*{1.5em}\textbf{A.} 26\quad\quad \textbf{B.} 25,68\quad\quad \textbf{C.} 25,6\quad\quad \textbf{D.} 25,7\par\vspace{2pt}
\textbf{\color{primary} Câu 7.} Xác định giá trị hoặc khẳng định đúng cho bài toán sau: Phân số thập phân -39/10 được viết dưới dạng số thập phân là:\par
\hspace*{1.5em}\textbf{A.} -3.9\quad\quad \textbf{B.} -39.0\quad\quad \textbf{C.} -0.39\quad\quad \textbf{D.} 3.9\par\vspace{2pt}
\textbf{\color{primary} Câu 8.} Áp dụng trực tiếp định nghĩa và quy tắc trong Chương VII: Số đối của số thập phân 19,35 là:\par
\hspace*{1.5em}\textbf{A.} 0,1935\quad\quad \textbf{B.} 19,53\quad\quad \textbf{C.} -19,53\quad\quad \textbf{D.} -19,35\par\vspace{2pt}
\textbf{\color{primary} Câu 9.} Kiểm tra nhanh kiến thức nền tảng: Làm tròn số thập phân 28,682 đến hàng phần mười ta được:\par
\hspace*{1.5em}\textbf{A.} 29\quad\quad \textbf{B.} 28,68\quad\quad \textbf{C.} 28,6\quad\quad \textbf{D.} 28,7\par\vspace{2pt}
\textbf{\color{primary} Câu 10.} Cho dữ kiện nhận biết sau đây: Phân số thập phân -51/10 được viết dưới dạng số thập phân là. Phương án chính xác là:\par
\hspace*{1.5em}\textbf{A.} -0.51\quad\quad \textbf{B.} -51.0\quad\quad \textbf{C.} 5.1\quad\quad \textbf{D.} -5.1\par\vspace{2pt}
\textbf{\color{primary} Câu 11.} Hãy chỉ ra kết quả đúng trong bốn lựa chọn A, B, C, D đối với yêu cầu: Số đối của số thập phân 22,35 là:\par
\hspace*{1.5em}\textbf{A.} -22,35\quad\quad \textbf{B.} -22,53\quad\quad \textbf{C.} 22,53\quad\quad \textbf{D.} 0,2235\par\vspace{2pt}
\textbf{\color{primary} Câu 12.} Câu hỏi ôn tập lý thuyết và kỹ năng cơ bản (Chương VII): Làm tròn số thập phân 31,682 đến hàng phần mười ta được:\par
\hspace*{1.5em}\textbf{A.} 32\quad\quad \textbf{B.} 31,7\quad\quad \textbf{C.} 31,68\quad\quad \textbf{D.} 31,6\par\vspace{2pt}
\textbf{\color{primary} Câu 13.} Theo nội dung sách giáo khoa Kết nối tri thức với cuộc sống (Chương VII), phân số thập phân -63/10 được viết dưới dạng số thập phân là:\par
\hspace*{1.5em}\textbf{A.} -0.63\quad\quad \textbf{B.} -6.3\quad\quad \textbf{C.} -63.0\quad\quad \textbf{D.} 6.3\par\vspace{2pt}
\textbf{\color{primary} Câu 14.} Tìm đáp số đúng của bài toán cơ bản sau: Số đối của số thập phân 25,35 là:\par
\hspace*{1.5em}\textbf{A.} -25,53\quad\quad \textbf{B.} 25,53\quad\quad \textbf{C.} 0,2535\quad\quad \textbf{D.} -25,35\par\vspace{2pt}
\textbf{\color{primary} Câu 15.} Xét đối tượng toán học trong Chương VII: Làm tròn số thập phân 34,682 đến hàng phần mười ta được. Giá trị hoặc tính chất tương ứng là:\par
\hspace*{1.5em}\textbf{A.} 34,68\quad\quad \textbf{B.} 34,6\quad\quad \textbf{C.} 35\quad\quad \textbf{D.} 34,7\par\vspace{2pt}
\textbf{\color{primary} Câu 16.} Hoàn thành phép xác định cơ bản sau đây: Phân số thập phân -75/10 được viết dưới dạng số thập phân là:\par
\hspace*{1.5em}\textbf{A.} 7.5\quad\quad \textbf{B.} -75.0\quad\quad \textbf{C.} -7.5\quad\quad \textbf{D.} -0.75\par\vspace{2pt}
\textbf{\color{primary} Câu 17.} Trong các khẳng định hoặc giá trị dưới đây, đâu là đáp án đúng cho câu hỏi: Số đối của số thập phân 28,35 là?\par
\hspace*{1.5em}\textbf{A.} -28,53\quad\quad \textbf{B.} 28,53\quad\quad \textbf{C.} -28,35\quad\quad \textbf{D.} 0,2835\par\vspace{2pt}
\textbf{\color{primary} Câu 18.} Bài tập khởi động Chương VII: Làm tròn số thập phân 37,682 đến hàng phần mười ta được:\par
\hspace*{1.5em}\textbf{A.} 37,6\quad\quad \textbf{B.} 37,7\quad\quad \textbf{C.} 38\quad\quad \textbf{D.} 37,68\par\vspace{2pt}
\textbf{\color{primary} Câu 19.} Quan sát dữ kiện đề bài: Phân số thập phân -87/10 được viết dưới dạng số thập phân là. Hãy chọn phương án đúng:\par
\hspace*{1.5em}\textbf{A.} 8.7\quad\quad \textbf{B.} -8.7\quad\quad \textbf{C.} -0.87\quad\quad \textbf{D.} -87.0\par\vspace{2pt}
\textbf{\color{primary} Câu 20.} Sử dụng công thức hoặc tính chất đầu bài của Chương VII để giải quyết yêu cầu: Số đối của số thập phân 31,35 là:\par
\hspace*{1.5em}\textbf{A.} -31,53\quad\quad \textbf{B.} -31,35\quad\quad \textbf{C.} 31,53\quad\quad \textbf{D.} 0,3135\par\vspace{2pt}
\subsection*{\color{accent} Phần II. Mức độ Trung bình (Thông hiểu -- Câu 21 đến Câu 40)}
\textbf{\color{primary} Câu 21.} Tỉ số phần trăm của 3 và 20 bằng:\par
\hspace*{1.5em}\textbf{A.} 15\%\quad\quad \textbf{B.} 6\%\quad\quad \textbf{C.} 30\%\quad\quad \textbf{D.} 12\%\par\vspace{2pt}
\textbf{\color{primary} Câu 22.} Vận dụng quy tắc trong Chương VII (Số thập phân) để giải bài toán: Tính nhanh giá trị biểu thức: 5,5 $\cdot$ 3,8 + 5,5 $\cdot$ 6,2.:\par
\hspace*{1.5em}\textbf{A.} 49.5\quad\quad \textbf{B.} 550\quad\quad \textbf{C.} 55\quad\quad \textbf{D.} 44\par\vspace{2pt}
\textbf{\color{primary} Câu 23.} Biến đổi và tính toán theo kiến thức Chương VII: Một mặt hàng có giá niêm yết 240.000 đồng. Nhân dịp lễ, cửa hàng giảm giá 15\%. Số tiền được giảm là. Kết quả thu được là:\par
\hspace*{1.5em}\textbf{A.} 36.000 đồng\quad\quad \textbf{B.} 204.000 đồng\quad\quad \textbf{C.} 24.000 đồng\quad\quad \textbf{D.} 48.000 đồng\par\vspace{2pt}
\textbf{\color{primary} Câu 24.} Xét bài toán mức độ thông hiểu sau: Tỉ số phần trăm của 6 và 20 bằng. Đáp án đúng là:\par
\hspace*{1.5em}\textbf{A.} 24\%\quad\quad \textbf{B.} 12\%\quad\quad \textbf{C.} 60\%\quad\quad \textbf{D.} 30\%\par\vspace{2pt}
\textbf{\color{primary} Câu 25.} Để hoàn thành phiếu học tập Chương VII, hãy xác định: Tính nhanh giá trị biểu thức: 8,5 $\cdot$ 3,8 + 8,5 $\cdot$ 6,2.:\par
\hspace*{1.5em}\textbf{A.} 850\quad\quad \textbf{B.} 85\quad\quad \textbf{C.} 76.5\quad\quad \textbf{D.} 68\par\vspace{2pt}
\textbf{\color{primary} Câu 26.} Phân tích dữ kiện bài toán: Một mặt hàng có giá niêm yết 300.000 đồng. Nhân dịp lễ, cửa hàng giảm giá 15\%. Số tiền được giảm là. Giá trị cần tìm bằng:\par
\hspace*{1.5em}\textbf{A.} 30.000 đồng\quad\quad \textbf{B.} 45.000 đồng\quad\quad \textbf{C.} 60.000 đồng\quad\quad \textbf{D.} 255.000 đồng\par\vspace{2pt}
\textbf{\color{primary} Câu 27.} Áp dụng tính chất đã học ở Chương VII vào bài tập sau: Tỉ số phần trăm của 9 và 20 bằng:\par
\hspace*{1.5em}\textbf{A.} 18\%\quad\quad \textbf{B.} 36\%\quad\quad \textbf{C.} 90\%\quad\quad \textbf{D.} 45\%\par\vspace{2pt}
\textbf{\color{primary} Câu 28.} Thực hiện các bước suy luận hợp lý để giải quyết yêu cầu: Tính nhanh giá trị biểu thức: 11,5 $\cdot$ 3,8 + 11,5 $\cdot$ 6,2.:\par
\hspace*{1.5em}\textbf{A.} 92\quad\quad \textbf{B.} 1150\quad\quad \textbf{C.} 103.5\quad\quad \textbf{D.} 115\par\vspace{2pt}
\textbf{\color{primary} Câu 29.} Cho bài toán tự luyện (Số thập phân): Một mặt hàng có giá niêm yết 360.000 đồng. Nhân dịp lễ, cửa hàng giảm giá 15\%. Số tiền được giảm là. Phương án nào sau đây chính xác?\par
\hspace*{1.5em}\textbf{A.} 54.000 đồng\quad\quad \textbf{B.} 306.000 đồng\quad\quad \textbf{C.} 72.000 đồng\quad\quad \textbf{D.} 36.000 đồng\par\vspace{2pt}
\textbf{\color{primary} Câu 30.} Kết hợp các phép biến đổi trong Chương VII, hãy tìm kết quả của bài toán: Tỉ số phần trăm của 12 và 20 bằng:\par
\hspace*{1.5em}\textbf{A.} 48\%\quad\quad \textbf{B.} 24\%\quad\quad \textbf{C.} 120\%\quad\quad \textbf{D.} 60\%\par\vspace{2pt}
\textbf{\color{primary} Câu 31.} Một bài kiểm tra định kỳ nội dung Chương VII nêu câu hỏi: Tính nhanh giá trị biểu thức: 14,5 $\cdot$ 3,8 + 14,5 $\cdot$ 6,2.. Hãy chọn đáp số đúng:\par
\hspace*{1.5em}\textbf{A.} 130.5\quad\quad \textbf{B.} 116\quad\quad \textbf{C.} 1450\quad\quad \textbf{D.} 145\par\vspace{2pt}
\textbf{\color{primary} Câu 32.} Tìm lời giải đúng cho bài toán thông hiểu: Một mặt hàng có giá niêm yết 420.000 đồng. Nhân dịp lễ, cửa hàng giảm giá 15\%. Số tiền được giảm là:\par
\hspace*{1.5em}\textbf{A.} 63.000 đồng\quad\quad \textbf{B.} 84.000 đồng\quad\quad \textbf{C.} 357.000 đồng\quad\quad \textbf{D.} 42.000 đồng\par\vspace{2pt}
\textbf{\color{primary} Câu 33.} Sử dụng mối liên hệ giữa các đại lượng trong Chương VII: Tỉ số phần trăm của 15 và 20 bằng:\par
\hspace*{1.5em}\textbf{A.} 30\%\quad\quad \textbf{B.} 150\%\quad\quad \textbf{C.} 75\%\quad\quad \textbf{D.} 60\%\par\vspace{2pt}
\textbf{\color{primary} Câu 34.} Hãy tính toán và chọn phương án đúng đối với bài tập: Tính nhanh giá trị biểu thức: 17,5 $\cdot$ 3,8 + 17,5 $\cdot$ 6,2.:\par
\hspace*{1.5em}\textbf{A.} 1750\quad\quad \textbf{B.} 175\quad\quad \textbf{C.} 157.5\quad\quad \textbf{D.} 140\par\vspace{2pt}
\textbf{\color{primary} Câu 35.} Dựa trên các tính chất trọng tâm của Chương VII, xét bài toán: Một mặt hàng có giá niêm yết 480.000 đồng. Nhân dịp lễ, cửa hàng giảm giá 15\%. Số tiền được giảm là:\par
\hspace*{1.5em}\textbf{A.} 96.000 đồng\quad\quad \textbf{B.} 408.000 đồng\quad\quad \textbf{C.} 48.000 đồng\quad\quad \textbf{D.} 72.000 đồng\par\vspace{2pt}
\textbf{\color{primary} Câu 36.} Thực hành giải toán (Số thập phân): Tỉ số phần trăm của 18 và 20 bằng. Đáp án nào dưới đây thỏa mãn?\par
\hspace*{1.5em}\textbf{A.} 72\%\quad\quad \textbf{B.} 180\%\quad\quad \textbf{C.} 36\%\quad\quad \textbf{D.} 90\%\par\vspace{2pt}
\textbf{\color{primary} Câu 37.} Kiểm tra kỹ năng biến đổi và suy luận qua bài toán: Tính nhanh giá trị biểu thức: 20,5 $\cdot$ 3,8 + 20,5 $\cdot$ 6,2.:\par
\hspace*{1.5em}\textbf{A.} 164\quad\quad \textbf{B.} 2050\quad\quad \textbf{C.} 184.5\quad\quad \textbf{D.} 205\par\vspace{2pt}
\textbf{\color{primary} Câu 38.} Cho giả thiết thuộc phạm vi Chương VII: Một mặt hàng có giá niêm yết 540.000 đồng. Nhân dịp lễ, cửa hàng giảm giá 15\%. Số tiền được giảm là. Kết luận đúng là:\par
\hspace*{1.5em}\textbf{A.} 459.000 đồng\quad\quad \textbf{B.} 54.000 đồng\quad\quad \textbf{C.} 81.000 đồng\quad\quad \textbf{D.} 108.000 đồng\par\vspace{2pt}
\textbf{\color{primary} Câu 39.} Rút gọn hoặc tính giá trị theo yêu cầu sau của Chương VII: Tỉ số phần trăm của 21 và 20 bằng:\par
\hspace*{1.5em}\textbf{A.} 210\%\quad\quad \textbf{B.} 42\%\quad\quad \textbf{C.} 84\%\quad\quad \textbf{D.} 105\%\par\vspace{2pt}
\textbf{\color{primary} Câu 40.} Xác định phương án đúng sau khi thực hiện phép tính / suy luận: Tính nhanh giá trị biểu thức: 23,5 $\cdot$ 3,8 + 23,5 $\cdot$ 6,2.:\par
\hspace*{1.5em}\textbf{A.} 188\quad\quad \textbf{B.} 235\quad\quad \textbf{C.} 2350\quad\quad \textbf{D.} 211.5\par\vspace{2pt}
\subsection*{\color{accent} Phần III. Mức độ Khó (Vận dụng \& Vận dụng cao -- Câu 41 đến Câu 50)}
\textbf{\color{primary} Câu 41.} Sau khi giảm giá 20\%, một chiếc máy tính cầm tay có giá bán là 320.000 đồng. Giá niêm yết ban đầu của chiếc máy tính đó là:\par
\hspace*{1.5em}\textbf{A.} 384.000 đồng\quad\quad \textbf{B.} 400.000 đồng (1)\quad\quad \textbf{C.} 256.000 đồng\quad\quad \textbf{D.} 400.000 đồng\par\vspace{2pt}
\textbf{\color{primary} Câu 42.} Bài toán vận dụng (Chương VII – Số thập phân): Lãi suất tiết kiệm kỳ hạn 1 năm là 6\%/năm. Nếu gửi 15 triệu đồng thì sau 1 năm tổng số tiền cả vốn lẫn lãi nhận được là. Kết quả đúng là:\par
\hspace*{1.5em}\textbf{A.} 24.00 triệu đồng\quad\quad \textbf{B.} 16.80 triệu đồng\quad\quad \textbf{C.} 15.90 triệu đồng\quad\quad \textbf{D.} 0.90 triệu đồng\par\vspace{2pt}
\textbf{\color{primary} Câu 43.} Vận dụng kiến thức nâng cao trong phạm vi Chương VII để giải bài toán: Sau khi giảm giá 20\%, một chiếc máy tính cầm tay có giá bán là 400.000 đồng. Giá niêm yết ban đầu của chiếc máy tính đó là:\par
\hspace*{1.5em}\textbf{A.} 500.000 đồng\quad\quad \textbf{B.} 320.000 đồng\quad\quad \textbf{C.} 480.000 đồng\quad\quad \textbf{D.} 464.000 đồng\par\vspace{2pt}
\textbf{\color{primary} Câu 44.} Phân tích và giải bài toán tổng hợp sau của Chương VII: Lãi suất tiết kiệm kỳ hạn 1 năm là 6\%/năm. Nếu gửi 25 triệu đồng thì sau 1 năm tổng số tiền cả vốn lẫn lãi nhận được là:\par
\hspace*{1.5em}\textbf{A.} 40.00 triệu đồng\quad\quad \textbf{B.} 1.50 triệu đồng\quad\quad \textbf{C.} 28.00 triệu đồng\quad\quad \textbf{D.} 26.50 triệu đồng\par\vspace{2pt}
\textbf{\color{primary} Câu 45.} Thử thách tư duy trong phạm vi kiến thức Chương VII: Sau khi giảm giá 20\%, một chiếc máy tính cầm tay có giá bán là 480.000 đồng. Giá niêm yết ban đầu của chiếc máy tính đó là. Đáp số chính xác là:\par
\hspace*{1.5em}\textbf{A.} 384.000 đồng\quad\quad \textbf{B.} 544.000 đồng\quad\quad \textbf{C.} 576.000 đồng\quad\quad \textbf{D.} 600.000 đồng\par\vspace{2pt}
\textbf{\color{primary} Câu 46.} Sử dụng phương pháp giải toán của Chương VII để tìm đáp án cho bài toán: Lãi suất tiết kiệm kỳ hạn 1 năm là 6\%/năm. Nếu gửi 35 triệu đồng thì sau 1 năm tổng số tiền cả vốn lẫn lãi nhận được là:\par
\hspace*{1.5em}\textbf{A.} 56.00 triệu đồng\quad\quad \textbf{B.} 39.20 triệu đồng\quad\quad \textbf{C.} 2.10 triệu đồng\quad\quad \textbf{D.} 37.10 triệu đồng\par\vspace{2pt}
\textbf{\color{primary} Câu 47.} Xét bài toán vận dụng thực tiễn / nâng cao thuộc Chương VII: Sau khi giảm giá 20\%, một chiếc máy tính cầm tay có giá bán là 560.000 đồng. Giá niêm yết ban đầu của chiếc máy tính đó là:\par
\hspace*{1.5em}\textbf{A.} 624.000 đồng\quad\quad \textbf{B.} 448.000 đồng\quad\quad \textbf{C.} 672.000 đồng\quad\quad \textbf{D.} 700.000 đồng\par\vspace{2pt}
\textbf{\color{primary} Câu 48.} Tìm kết quả cuối cùng của bài toán phân loại học sinh khá giỏi (Chương VII): Lãi suất tiết kiệm kỳ hạn 1 năm là 6\%/năm. Nếu gửi 45 triệu đồng thì sau 1 năm tổng số tiền cả vốn lẫn lãi nhận được là:\par
\hspace*{1.5em}\textbf{A.} 47.70 triệu đồng\quad\quad \textbf{B.} 2.70 triệu đồng\quad\quad \textbf{C.} 50.40 triệu đồng\quad\quad \textbf{D.} 72.00 triệu đồng\par\vspace{2pt}
\textbf{\color{primary} Câu 49.} Kết hợp các tính chất đã học trong Chương VII, hãy giải quyết bài toán: Sau khi giảm giá 20\%, một chiếc máy tính cầm tay có giá bán là 640.000 đồng. Giá niêm yết ban đầu của chiếc máy tính đó là:\par
\hspace*{1.5em}\textbf{A.} 800.000 đồng\quad\quad \textbf{B.} 704.000 đồng\quad\quad \textbf{C.} 512.000 đồng\quad\quad \textbf{D.} 768.000 đồng\par\vspace{2pt}
\textbf{\color{primary} Câu 50.} Câu hỏi vận dụng cao cuối Chương VII: Lãi suất tiết kiệm kỳ hạn 1 năm là 6\%/năm. Nếu gửi 55 triệu đồng thì sau 1 năm tổng số tiền cả vốn lẫn lãi nhận được là. Phương án đúng là:\par
\hspace*{1.5em}\textbf{A.} 58.30 triệu đồng\quad\quad \textbf{B.} 61.60 triệu đồng\quad\quad \textbf{C.} 88.00 triệu đồng\quad\quad \textbf{D.} 3.30 triệu đồng\par\vspace{2pt}
\subsection*{\color{primary} Bảng đáp án Chương VII: Số thập phân}
\begin{center}
\renewcommand{\arraystretch}{1.25}
\begin{tabularx}{\textwidth}{*{10}{>{\centering\arraybackslash}X}}
\toprule
\textbf{\color{primary} Cột 1} & \textbf{\color{primary} Cột 2} & \textbf{\color{primary} Cột 3} & \textbf{\color{primary} Cột 4} & \textbf{\color{primary} Cột 5} & \textbf{\color{primary} Cột 6} & \textbf{\color{primary} Cột 7} & \textbf{\color{primary} Cột 8} & \textbf{\color{primary} Cột 9} & \textbf{\color{primary} Cột 10} \\
\midrule
\textbf{1.A} & \textbf{2.B} & \textbf{3.A} & \textbf{4.B} & \textbf{5.A} & \textbf{6.D} & \textbf{7.A} & \textbf{8.D} & \textbf{9.D} & \textbf{10.D} \\
\rowcolor{rowtint} \textbf{11.A} & \textbf{12.B} & \textbf{13.B} & \textbf{14.D} & \textbf{15.D} & \textbf{16.C} & \textbf{17.C} & \textbf{18.B} & \textbf{19.B} & \textbf{20.B} \\
\textbf{21.A} & \textbf{22.C} & \textbf{23.A} & \textbf{24.D} & \textbf{25.B} & \textbf{26.B} & \textbf{27.D} & \textbf{28.D} & \textbf{29.A} & \textbf{30.D} \\
\rowcolor{rowtint} \textbf{31.D} & \textbf{32.A} & \textbf{33.C} & \textbf{34.B} & \textbf{35.D} & \textbf{36.D} & \textbf{37.D} & \textbf{38.C} & \textbf{39.D} & \textbf{40.B} \\
\textbf{41.D} & \textbf{42.C} & \textbf{43.A} & \textbf{44.D} & \textbf{45.D} & \textbf{46.D} & \textbf{47.D} & \textbf{48.A} & \textbf{49.A} & \textbf{50.A} \\
\bottomrule
\end{tabularx}
\end{center}\vspace{10pt}
\section*{\color{primary} CHƯƠNG VIII: NHỮNG HÌNH HÌNH HỌC CƠ BẢN}
\textit{Trọng tâm kiến thức SGK: Điểm, đường thẳng, ba điểm thẳng hàng, hai đường thẳng cắt nhau và song song, tia, đoạn thẳng, độ dài đoạn thẳng, trung điểm của đoạn thẳng, góc và số đo góc.}\par\vspace{4pt}
\subsection*{\color{accent} Phần I. Mức độ Dễ (Nhận biết -- Câu 1 đến Câu 20)}
\textbf{\color{primary} Câu 1.} Cho đoạn thẳng AB có độ dài bằng 8 cm. Nếu M là trung điểm của đoạn thẳng AB thì độ dài đoạn AM bằng:\par
\hspace*{1.5em}\textbf{A.} 2 cm\quad\quad \textbf{B.} 16 cm\quad\quad \textbf{C.} 8 cm\quad\quad \textbf{D.} 4 cm\par\vspace{2pt}
\textbf{\color{primary} Câu 2.} Hãy trả lời câu hỏi cơ bản sau: Góc có số đo bằng 29° thuộc loại góc nào sau đây?\par
\hspace*{1.5em}\textbf{A.} Góc bẹt\quad\quad \textbf{B.} Góc tù\quad\quad \textbf{C.} Góc vuông\quad\quad \textbf{D.} Góc nhọn\par\vspace{2pt}
\textbf{\color{primary} Câu 3.} Trong bài học thuộc Chương VIII, xét nội dung: Qua hai điểm phân biệt A và B cho trước, ta vẽ được bao nhiêu đường thẳng đi qua cả hai điểm đó. Đáp án đúng là:\par
\begin{itemize}\setlength{\itemsep}{1pt}\setlength{\parskip}{0pt}
  \item[] \textbf{A.} Vô số đường thẳng
  \item[] \textbf{B.} Chỉ có đúng 1 đường thẳng
  \item[] \textbf{C.} Không vẽ được đường thẳng nào
  \item[] \textbf{D.} 2 đường thẳng
\end{itemize}\vspace{2pt}
\textbf{\color{primary} Câu 4.} Dựa vào kiến thức chuẩn SGK của Chương VIII (Những hình hình học cơ bản), cho đoạn thẳng AB có độ dài bằng 14 cm. Nếu M là trung điểm của đoạn thẳng AB thì độ dài đoạn AM bằng:\par
\hspace*{1.5em}\textbf{A.} 14 cm\quad\quad \textbf{B.} 5 cm\quad\quad \textbf{C.} 7 cm\quad\quad \textbf{D.} 28 cm\par\vspace{2pt}
\textbf{\color{primary} Câu 5.} Chọn phương án trả lời đúng nhất: Góc có số đo bằng 41° thuộc loại góc nào sau đây?\par
\hspace*{1.5em}\textbf{A.} Góc tù\quad\quad \textbf{B.} Góc nhọn\quad\quad \textbf{C.} Góc bẹt\quad\quad \textbf{D.} Góc vuông\par\vspace{2pt}
\textbf{\color{primary} Câu 6.} Khi làm bài tập cơ bản về những hình hình học cơ bản, một học sinh gặp câu hỏi: Qua hai điểm phân biệt A và B cho trước, ta vẽ được bao nhiêu đường thẳng đi qua cả hai điểm đó. Kết quả nào dưới đây đúng?\par
\begin{itemize}\setlength{\itemsep}{1pt}\setlength{\parskip}{0pt}
  \item[] \textbf{A.} Vô số đường thẳng
  \item[] \textbf{B.} 2 đường thẳng
  \item[] \textbf{C.} Chỉ có đúng 1 đường thẳng
  \item[] \textbf{D.} Không vẽ được đường thẳng nào
\end{itemize}\vspace{2pt}
\textbf{\color{primary} Câu 7.} Xác định giá trị hoặc khẳng định đúng cho bài toán sau: Cho đoạn thẳng AB có độ dài bằng 20 cm. Nếu M là trung điểm của đoạn thẳng AB thì độ dài đoạn AM bằng:\par
\hspace*{1.5em}\textbf{A.} 40 cm\quad\quad \textbf{B.} 8 cm\quad\quad \textbf{C.} 10 cm\quad\quad \textbf{D.} 20 cm\par\vspace{2pt}
\textbf{\color{primary} Câu 8.} Áp dụng trực tiếp định nghĩa và quy tắc trong Chương VIII: Góc có số đo bằng 53° thuộc loại góc nào sau đây:\par
\hspace*{1.5em}\textbf{A.} Góc nhọn\quad\quad \textbf{B.} Góc bẹt\quad\quad \textbf{C.} Góc vuông\quad\quad \textbf{D.} Góc tù\par\vspace{2pt}
\textbf{\color{primary} Câu 9.} Kiểm tra nhanh kiến thức nền tảng: Qua hai điểm phân biệt A và B cho trước, ta vẽ được bao nhiêu đường thẳng đi qua cả hai điểm đó?\par
\begin{itemize}\setlength{\itemsep}{1pt}\setlength{\parskip}{0pt}
  \item[] \textbf{A.} Vô số đường thẳng
  \item[] \textbf{B.} 2 đường thẳng
  \item[] \textbf{C.} Chỉ có đúng 1 đường thẳng
  \item[] \textbf{D.} Không vẽ được đường thẳng nào
\end{itemize}\vspace{2pt}
\textbf{\color{primary} Câu 10.} Cho dữ kiện nhận biết sau đây: Cho đoạn thẳng AB có độ dài bằng 26 cm. Nếu M là trung điểm của đoạn thẳng AB thì độ dài đoạn AM bằng. Phương án chính xác là:\par
\hspace*{1.5em}\textbf{A.} 13 cm\quad\quad \textbf{B.} 26 cm\quad\quad \textbf{C.} 11 cm\quad\quad \textbf{D.} 52 cm\par\vspace{2pt}
\textbf{\color{primary} Câu 11.} Hãy chỉ ra kết quả đúng trong bốn lựa chọn A, B, C, D đối với yêu cầu: Góc có số đo bằng 65° thuộc loại góc nào sau đây:\par
\hspace*{1.5em}\textbf{A.} Góc vuông\quad\quad \textbf{B.} Góc bẹt\quad\quad \textbf{C.} Góc nhọn\quad\quad \textbf{D.} Góc tù\par\vspace{2pt}
\textbf{\color{primary} Câu 12.} Câu hỏi ôn tập lý thuyết và kỹ năng cơ bản (Chương VIII): Qua hai điểm phân biệt A và B cho trước, ta vẽ được bao nhiêu đường thẳng đi qua cả hai điểm đó:\par
\begin{itemize}\setlength{\itemsep}{1pt}\setlength{\parskip}{0pt}
  \item[] \textbf{A.} 2 đường thẳng
  \item[] \textbf{B.} Chỉ có đúng 1 đường thẳng
  \item[] \textbf{C.} Vô số đường thẳng
  \item[] \textbf{D.} Không vẽ được đường thẳng nào
\end{itemize}\vspace{2pt}
\textbf{\color{primary} Câu 13.} Theo nội dung sách giáo khoa Kết nối tri thức với cuộc sống (Chương VIII), cho đoạn thẳng AB có độ dài bằng 32 cm. Nếu M là trung điểm của đoạn thẳng AB thì độ dài đoạn AM bằng:\par
\hspace*{1.5em}\textbf{A.} 16 cm\quad\quad \textbf{B.} 64 cm\quad\quad \textbf{C.} 32 cm\quad\quad \textbf{D.} 14 cm\par\vspace{2pt}
\textbf{\color{primary} Câu 14.} Tìm đáp số đúng của bài toán cơ bản sau: Góc có số đo bằng 77° thuộc loại góc nào sau đây:\par
\hspace*{1.5em}\textbf{A.} Góc bẹt\quad\quad \textbf{B.} Góc nhọn\quad\quad \textbf{C.} Góc vuông\quad\quad \textbf{D.} Góc tù\par\vspace{2pt}
\textbf{\color{primary} Câu 15.} Xét đối tượng toán học trong Chương VIII: Qua hai điểm phân biệt A và B cho trước, ta vẽ được bao nhiêu đường thẳng đi qua cả hai điểm đó. Giá trị hoặc tính chất tương ứng là:\par
\begin{itemize}\setlength{\itemsep}{1pt}\setlength{\parskip}{0pt}
  \item[] \textbf{A.} Vô số đường thẳng
  \item[] \textbf{B.} Không vẽ được đường thẳng nào
  \item[] \textbf{C.} 2 đường thẳng
  \item[] \textbf{D.} Chỉ có đúng 1 đường thẳng
\end{itemize}\vspace{2pt}
\textbf{\color{primary} Câu 16.} Hoàn thành phép xác định cơ bản sau đây: Cho đoạn thẳng AB có độ dài bằng 38 cm. Nếu M là trung điểm của đoạn thẳng AB thì độ dài đoạn AM bằng:\par
\hspace*{1.5em}\textbf{A.} 38 cm\quad\quad \textbf{B.} 17 cm\quad\quad \textbf{C.} 76 cm\quad\quad \textbf{D.} 19 cm\par\vspace{2pt}
\textbf{\color{primary} Câu 17.} Trong các khẳng định hoặc giá trị dưới đây, đâu là đáp án đúng cho câu hỏi: Góc có số đo bằng 89° thuộc loại góc nào sau đây?\par
\hspace*{1.5em}\textbf{A.} Góc tù\quad\quad \textbf{B.} Góc nhọn\quad\quad \textbf{C.} Góc vuông\quad\quad \textbf{D.} Góc bẹt\par\vspace{2pt}
\textbf{\color{primary} Câu 18.} Bài tập khởi động Chương VIII: Qua hai điểm phân biệt A và B cho trước, ta vẽ được bao nhiêu đường thẳng đi qua cả hai điểm đó:\par
\begin{itemize}\setlength{\itemsep}{1pt}\setlength{\parskip}{0pt}
  \item[] \textbf{A.} Không vẽ được đường thẳng nào
  \item[] \textbf{B.} Vô số đường thẳng
  \item[] \textbf{C.} 2 đường thẳng
  \item[] \textbf{D.} Chỉ có đúng 1 đường thẳng
\end{itemize}\vspace{2pt}
\textbf{\color{primary} Câu 19.} Quan sát dữ kiện đề bài: Cho đoạn thẳng AB có độ dài bằng 44 cm. Nếu M là trung điểm của đoạn thẳng AB thì độ dài đoạn AM bằng. Hãy chọn phương án đúng:\par
\hspace*{1.5em}\textbf{A.} 22 cm\quad\quad \textbf{B.} 88 cm\quad\quad \textbf{C.} 44 cm\quad\quad \textbf{D.} 20 cm\par\vspace{2pt}
\textbf{\color{primary} Câu 20.} Sử dụng công thức hoặc tính chất đầu bài của Chương VIII để giải quyết yêu cầu: Góc có số đo bằng 101° thuộc loại góc nào sau đây:\par
\hspace*{1.5em}\textbf{A.} Góc vuông\quad\quad \textbf{B.} Góc bẹt\quad\quad \textbf{C.} Góc tù\quad\quad \textbf{D.} Góc nhọn\par\vspace{2pt}
\subsection*{\color{accent} Phần II. Mức độ Trung bình (Thông hiểu -- Câu 21 đến Câu 40)}
\textbf{\color{primary} Câu 21.} Trên tia Ox lấy hai điểm A và B sao cho OA = 3 cm, OB = 9 cm. Độ dài đoạn thẳng AB bằng:\par
\hspace*{1.5em}\textbf{A.} 12 cm\quad\quad \textbf{B.} 3 cm\quad\quad \textbf{C.} 6 cm\quad\quad \textbf{D.} 9 cm\par\vspace{2pt}
\textbf{\color{primary} Câu 22.} Vận dụng quy tắc trong Chương VIII (Những hình hình học cơ bản) để giải bài toán: Cho điểm M nằm giữa hai điểm P và Q. Biết PM = 6 cm và PQ = 17 cm. Độ dài đoạn thẳng MQ là:\par
\hspace*{1.5em}\textbf{A.} 7 cm\quad\quad \textbf{B.} 10 cm\quad\quad \textbf{C.} 23 cm\quad\quad \textbf{D.} 11 cm\par\vspace{2pt}
\textbf{\color{primary} Câu 23.} Biến đổi và tính toán theo kiến thức Chương VIII: Sắp xếp các góc có số đo: $\alpha$ = 32°, $\beta$ = 90°, $\gamma$ = 112°, $\delta$ = 180° theo thứ tự số đo tăng dần. Kết quả thu được là:\par
\hspace*{1.5em}\textbf{A.} $\alpha$, $\beta$, $\gamma$, $\delta$\quad\quad \textbf{B.} $\beta$, $\alpha$, $\gamma$, $\delta$\quad\quad \textbf{C.} $\delta$, $\gamma$, $\beta$, $\alpha$\quad\quad \textbf{D.} $\alpha$, $\gamma$, $\beta$, $\delta$\par\vspace{2pt}
\textbf{\color{primary} Câu 24.} Xét bài toán mức độ thông hiểu sau: Trên tia Ox lấy hai điểm A và B sao cho OA = 6 cm, OB = 15 cm. Độ dài đoạn thẳng AB bằng. Đáp án đúng là:\par
\hspace*{1.5em}\textbf{A.} 6 cm\quad\quad \textbf{B.} 9 cm\quad\quad \textbf{C.} 15 cm\quad\quad \textbf{D.} 21 cm\par\vspace{2pt}
\textbf{\color{primary} Câu 25.} Để hoàn thành phiếu học tập Chương VIII, hãy xác định: Cho điểm M nằm giữa hai điểm P và Q. Biết PM = 9 cm và PQ = 26 cm. Độ dài đoạn thẳng MQ là:\par
\hspace*{1.5em}\textbf{A.} 17 cm (1)\quad\quad \textbf{B.} 35 cm\quad\quad \textbf{C.} 17 cm\quad\quad \textbf{D.} 13 cm\par\vspace{2pt}
\textbf{\color{primary} Câu 26.} Phân tích dữ kiện bài toán: Sắp xếp các góc có số đo: $\alpha$ = 35°, $\beta$ = 90°, $\gamma$ = 115°, $\delta$ = 180° theo thứ tự số đo tăng dần. Giá trị cần tìm bằng:\par
\hspace*{1.5em}\textbf{A.} $\beta$, $\alpha$, $\gamma$, $\delta$\quad\quad \textbf{B.} $\alpha$, $\beta$, $\gamma$, $\delta$\quad\quad \textbf{C.} $\delta$, $\gamma$, $\beta$, $\alpha$\quad\quad \textbf{D.} $\alpha$, $\gamma$, $\beta$, $\delta$\par\vspace{2pt}
\textbf{\color{primary} Câu 27.} Áp dụng tính chất đã học ở Chương VIII vào bài tập sau: Trên tia Ox lấy hai điểm A và B sao cho OA = 9 cm, OB = 21 cm. Độ dài đoạn thẳng AB bằng:\par
\hspace*{1.5em}\textbf{A.} 9 cm\quad\quad \textbf{B.} 30 cm\quad\quad \textbf{C.} 12 cm\quad\quad \textbf{D.} 21 cm\par\vspace{2pt}
\textbf{\color{primary} Câu 28.} Thực hiện các bước suy luận hợp lý để giải quyết yêu cầu: Cho điểm M nằm giữa hai điểm P và Q. Biết PM = 12 cm và PQ = 35 cm. Độ dài đoạn thẳng MQ là:\par
\hspace*{1.5em}\textbf{A.} 47 cm\quad\quad \textbf{B.} 16 cm\quad\quad \textbf{C.} 23 cm\quad\quad \textbf{D.} 19 cm\par\vspace{2pt}
\textbf{\color{primary} Câu 29.} Cho bài toán tự luyện (Những hình hình học cơ bản): Sắp xếp các góc có số đo: $\alpha$ = 38°, $\beta$ = 90°, $\gamma$ = 118°, $\delta$ = 180° theo thứ tự số đo tăng dần. Phương án nào sau đây chính xác?\par
\hspace*{1.5em}\textbf{A.} $\beta$, $\alpha$, $\gamma$, $\delta$\quad\quad \textbf{B.} $\alpha$, $\gamma$, $\beta$, $\delta$\quad\quad \textbf{C.} $\alpha$, $\beta$, $\gamma$, $\delta$\quad\quad \textbf{D.} $\delta$, $\gamma$, $\beta$, $\alpha$\par\vspace{2pt}
\textbf{\color{primary} Câu 30.} Kết hợp các phép biến đổi trong Chương VIII, hãy tìm kết quả của bài toán: Trên tia Ox lấy hai điểm A và B sao cho OA = 12 cm, OB = 27 cm. Độ dài đoạn thẳng AB bằng:\par
\hspace*{1.5em}\textbf{A.} 12 cm\quad\quad \textbf{B.} 39 cm\quad\quad \textbf{C.} 15 cm\quad\quad \textbf{D.} 27 cm\par\vspace{2pt}
\textbf{\color{primary} Câu 31.} Một bài kiểm tra định kỳ nội dung Chương VIII nêu câu hỏi: Cho điểm M nằm giữa hai điểm P và Q. Biết PM = 15 cm và PQ = 44 cm. Độ dài đoạn thẳng MQ là. Hãy chọn đáp số đúng:\par
\hspace*{1.5em}\textbf{A.} 59 cm\quad\quad \textbf{B.} 19 cm\quad\quad \textbf{C.} 25 cm\quad\quad \textbf{D.} 29 cm\par\vspace{2pt}
\textbf{\color{primary} Câu 32.} Tìm lời giải đúng cho bài toán thông hiểu: Sắp xếp các góc có số đo: $\alpha$ = 41°, $\beta$ = 90°, $\gamma$ = 121°, $\delta$ = 180° theo thứ tự số đo tăng dần:\par
\hspace*{1.5em}\textbf{A.} $\alpha$, $\gamma$, $\beta$, $\delta$\quad\quad \textbf{B.} $\delta$, $\gamma$, $\beta$, $\alpha$\quad\quad \textbf{C.} $\alpha$, $\beta$, $\gamma$, $\delta$\quad\quad \textbf{D.} $\beta$, $\alpha$, $\gamma$, $\delta$\par\vspace{2pt}
\textbf{\color{primary} Câu 33.} Sử dụng mối liên hệ giữa các đại lượng trong Chương VIII: Trên tia Ox lấy hai điểm A và B sao cho OA = 15 cm, OB = 33 cm. Độ dài đoạn thẳng AB bằng:\par
\hspace*{1.5em}\textbf{A.} 48 cm\quad\quad \textbf{B.} 15 cm\quad\quad \textbf{C.} 33 cm\quad\quad \textbf{D.} 18 cm\par\vspace{2pt}
\textbf{\color{primary} Câu 34.} Hãy tính toán và chọn phương án đúng đối với bài tập: Cho điểm M nằm giữa hai điểm P và Q. Biết PM = 18 cm và PQ = 53 cm. Độ dài đoạn thẳng MQ là:\par
\hspace*{1.5em}\textbf{A.} 71 cm\quad\quad \textbf{B.} 35 cm\quad\quad \textbf{C.} 31 cm\quad\quad \textbf{D.} 22 cm\par\vspace{2pt}
\textbf{\color{primary} Câu 35.} Dựa trên các tính chất trọng tâm của Chương VIII, xét bài toán: Sắp xếp các góc có số đo: $\alpha$ = 44°, $\beta$ = 90°, $\gamma$ = 124°, $\delta$ = 180° theo thứ tự số đo tăng dần:\par
\hspace*{1.5em}\textbf{A.} $\alpha$, $\gamma$, $\beta$, $\delta$\quad\quad \textbf{B.} $\delta$, $\gamma$, $\beta$, $\alpha$\quad\quad \textbf{C.} $\beta$, $\alpha$, $\gamma$, $\delta$\quad\quad \textbf{D.} $\alpha$, $\beta$, $\gamma$, $\delta$\par\vspace{2pt}
\textbf{\color{primary} Câu 36.} Thực hành giải toán (Những hình hình học cơ bản): Trên tia Ox lấy hai điểm A và B sao cho OA = 18 cm, OB = 39 cm. Độ dài đoạn thẳng AB bằng. Đáp án nào dưới đây thỏa mãn?\par
\hspace*{1.5em}\textbf{A.} 18 cm\quad\quad \textbf{B.} 57 cm\quad\quad \textbf{C.} 39 cm\quad\quad \textbf{D.} 21 cm\par\vspace{2pt}
\textbf{\color{primary} Câu 37.} Kiểm tra kỹ năng biến đổi và suy luận qua bài toán: Cho điểm M nằm giữa hai điểm P và Q. Biết PM = 21 cm và PQ = 62 cm. Độ dài đoạn thẳng MQ là:\par
\hspace*{1.5em}\textbf{A.} 41 cm\quad\quad \textbf{B.} 37 cm\quad\quad \textbf{C.} 25 cm\quad\quad \textbf{D.} 83 cm\par\vspace{2pt}
\textbf{\color{primary} Câu 38.} Cho giả thiết thuộc phạm vi Chương VIII: Sắp xếp các góc có số đo: $\alpha$ = 47°, $\beta$ = 90°, $\gamma$ = 127°, $\delta$ = 180° theo thứ tự số đo tăng dần. Kết luận đúng là:\par
\hspace*{1.5em}\textbf{A.} $\delta$, $\gamma$, $\beta$, $\alpha$\quad\quad \textbf{B.} $\alpha$, $\beta$, $\gamma$, $\delta$\quad\quad \textbf{C.} $\alpha$, $\gamma$, $\beta$, $\delta$\quad\quad \textbf{D.} $\beta$, $\alpha$, $\gamma$, $\delta$\par\vspace{2pt}
\textbf{\color{primary} Câu 39.} Rút gọn hoặc tính giá trị theo yêu cầu sau của Chương VIII: Trên tia Ox lấy hai điểm A và B sao cho OA = 21 cm, OB = 45 cm. Độ dài đoạn thẳng AB bằng:\par
\hspace*{1.5em}\textbf{A.} 21 cm\quad\quad \textbf{B.} 45 cm\quad\quad \textbf{C.} 24 cm\quad\quad \textbf{D.} 66 cm\par\vspace{2pt}
\textbf{\color{primary} Câu 40.} Xác định phương án đúng sau khi thực hiện phép tính / suy luận: Cho điểm M nằm giữa hai điểm P và Q. Biết PM = 24 cm và PQ = 71 cm. Độ dài đoạn thẳng MQ là:\par
\hspace*{1.5em}\textbf{A.} 28 cm\quad\quad \textbf{B.} 95 cm\quad\quad \textbf{C.} 47 cm\quad\quad \textbf{D.} 43 cm\par\vspace{2pt}
\subsection*{\color{accent} Phần III. Mức độ Khó (Vận dụng \& Vận dụng cao -- Câu 41 đến Câu 50)}
\textbf{\color{primary} Câu 41.} Cho 6 điểm phân biệt trong đó không có 3 điểm nào thẳng hàng. Cứ qua 2 điểm ta vẽ một đường thẳng. Số đường thẳng vẽ được là:\par
\hspace*{1.5em}\textbf{A.} 12 đường thẳng\quad\quad \textbf{B.} 15 đường thẳng\quad\quad \textbf{C.} 30 đường thẳng\quad\quad \textbf{D.} 13 đường thẳng\par\vspace{2pt}
\textbf{\color{primary} Câu 42.} Bài toán vận dụng (Chương VIII – Những hình hình học cơ bản): Cho đoạn thẳng AB = 16 cm. Gọi M là điểm nằm giữa A và B. Gọi I, K lần lượt là trung điểm của AM và MB. Độ dài đoạn thẳng IK bằng. Kết quả đúng là:\par
\hspace*{1.5em}\textbf{A.} 12 cm\quad\quad \textbf{B.} 16 cm\quad\quad \textbf{C.} 8 cm\quad\quad \textbf{D.} 4 cm\par\vspace{2pt}
\textbf{\color{primary} Câu 43.} Vận dụng kiến thức nâng cao trong phạm vi Chương VIII để giải bài toán: Cho 8 điểm phân biệt trong đó không có 3 điểm nào thẳng hàng. Cứ qua 2 điểm ta vẽ một đường thẳng. Số đường thẳng vẽ được là:\par
\hspace*{1.5em}\textbf{A.} 28 đường thẳng\quad\quad \textbf{B.} 56 đường thẳng\quad\quad \textbf{C.} 16 đường thẳng\quad\quad \textbf{D.} 26 đường thẳng\par\vspace{2pt}
\textbf{\color{primary} Câu 44.} Phân tích và giải bài toán tổng hợp sau của Chương VIII: Cho đoạn thẳng AB = 24 cm. Gọi M là điểm nằm giữa A và B. Gọi I, K lần lượt là trung điểm của AM và MB. Độ dài đoạn thẳng IK bằng:\par
\hspace*{1.5em}\textbf{A.} 12 cm\quad\quad \textbf{B.} 24 cm\quad\quad \textbf{C.} 6 cm\quad\quad \textbf{D.} 18 cm\par\vspace{2pt}
\textbf{\color{primary} Câu 45.} Thử thách tư duy trong phạm vi kiến thức Chương VIII: Cho 10 điểm phân biệt trong đó không có 3 điểm nào thẳng hàng. Cứ qua 2 điểm ta vẽ một đường thẳng. Số đường thẳng vẽ được là. Đáp số chính xác là:\par
\hspace*{1.5em}\textbf{A.} 90 đường thẳng\quad\quad \textbf{B.} 20 đường thẳng\quad\quad \textbf{C.} 43 đường thẳng\quad\quad \textbf{D.} 45 đường thẳng\par\vspace{2pt}
\textbf{\color{primary} Câu 46.} Sử dụng phương pháp giải toán của Chương VIII để tìm đáp án cho bài toán: Cho đoạn thẳng AB = 32 cm. Gọi M là điểm nằm giữa A và B. Gọi I, K lần lượt là trung điểm của AM và MB. Độ dài đoạn thẳng IK bằng:\par
\hspace*{1.5em}\textbf{A.} 24 cm\quad\quad \textbf{B.} 32 cm\quad\quad \textbf{C.} 8 cm\quad\quad \textbf{D.} 16 cm\par\vspace{2pt}
\textbf{\color{primary} Câu 47.} Xét bài toán vận dụng thực tiễn / nâng cao thuộc Chương VIII: Cho 12 điểm phân biệt trong đó không có 3 điểm nào thẳng hàng. Cứ qua 2 điểm ta vẽ một đường thẳng. Số đường thẳng vẽ được là:\par
\hspace*{1.5em}\textbf{A.} 132 đường thẳng\quad\quad \textbf{B.} 24 đường thẳng\quad\quad \textbf{C.} 64 đường thẳng\quad\quad \textbf{D.} 66 đường thẳng\par\vspace{2pt}
\textbf{\color{primary} Câu 48.} Tìm kết quả cuối cùng của bài toán phân loại học sinh khá giỏi (Chương VIII): Cho đoạn thẳng AB = 40 cm. Gọi M là điểm nằm giữa A và B. Gọi I, K lần lượt là trung điểm của AM và MB. Độ dài đoạn thẳng IK bằng:\par
\hspace*{1.5em}\textbf{A.} 20 cm\quad\quad \textbf{B.} 10 cm\quad\quad \textbf{C.} 30 cm\quad\quad \textbf{D.} 40 cm\par\vspace{2pt}
\textbf{\color{primary} Câu 49.} Kết hợp các tính chất đã học trong Chương VIII, hãy giải quyết bài toán: Cho 14 điểm phân biệt trong đó không có 3 điểm nào thẳng hàng. Cứ qua 2 điểm ta vẽ một đường thẳng. Số đường thẳng vẽ được là:\par
\hspace*{1.5em}\textbf{A.} 89 đường thẳng\quad\quad \textbf{B.} 182 đường thẳng\quad\quad \textbf{C.} 28 đường thẳng\quad\quad \textbf{D.} 91 đường thẳng\par\vspace{2pt}
\textbf{\color{primary} Câu 50.} Câu hỏi vận dụng cao cuối Chương VIII: Cho đoạn thẳng AB = 48 cm. Gọi M là điểm nằm giữa A và B. Gọi I, K lần lượt là trung điểm của AM và MB. Độ dài đoạn thẳng IK bằng. Phương án đúng là:\par
\hspace*{1.5em}\textbf{A.} 48 cm\quad\quad \textbf{B.} 24 cm\quad\quad \textbf{C.} 12 cm\quad\quad \textbf{D.} 36 cm\par\vspace{2pt}
\subsection*{\color{primary} Bảng đáp án Chương VIII: Những hình hình học cơ bản}
\begin{center}
\renewcommand{\arraystretch}{1.25}
\begin{tabularx}{\textwidth}{*{10}{>{\centering\arraybackslash}X}}
\toprule
\textbf{\color{primary} Cột 1} & \textbf{\color{primary} Cột 2} & \textbf{\color{primary} Cột 3} & \textbf{\color{primary} Cột 4} & \textbf{\color{primary} Cột 5} & \textbf{\color{primary} Cột 6} & \textbf{\color{primary} Cột 7} & \textbf{\color{primary} Cột 8} & \textbf{\color{primary} Cột 9} & \textbf{\color{primary} Cột 10} \\
\midrule
\textbf{1.D} & \textbf{2.D} & \textbf{3.B} & \textbf{4.C} & \textbf{5.B} & \textbf{6.C} & \textbf{7.C} & \textbf{8.A} & \textbf{9.C} & \textbf{10.A} \\
\rowcolor{rowtint} \textbf{11.C} & \textbf{12.B} & \textbf{13.A} & \textbf{14.B} & \textbf{15.D} & \textbf{16.D} & \textbf{17.B} & \textbf{18.D} & \textbf{19.A} & \textbf{20.C} \\
\textbf{21.C} & \textbf{22.D} & \textbf{23.A} & \textbf{24.B} & \textbf{25.C} & \textbf{26.B} & \textbf{27.C} & \textbf{28.C} & \textbf{29.C} & \textbf{30.C} \\
\rowcolor{rowtint} \textbf{31.D} & \textbf{32.C} & \textbf{33.D} & \textbf{34.B} & \textbf{35.D} & \textbf{36.D} & \textbf{37.A} & \textbf{38.B} & \textbf{39.C} & \textbf{40.C} \\
\textbf{41.B} & \textbf{42.C} & \textbf{43.A} & \textbf{44.A} & \textbf{45.D} & \textbf{46.D} & \textbf{47.D} & \textbf{48.A} & \textbf{49.D} & \textbf{50.B} \\
\bottomrule
\end{tabularx}
\end{center}\vspace{10pt}
\section*{\color{primary} CHƯƠNG IX: DỮ LIỆU VÀ XÁC SUẤT THỰC NGHIỆM}
\textit{Trọng tâm kiến thức SGK: Thu thập và phân loại dữ liệu, bảng thống kê, biểu đồ tranh, biểu đồ cột, biểu đồ cột kép, kết quả có thể và sự kiện trong trò chơi, xác suất thực nghiệm.}\par\vspace{4pt}
\subsection*{\color{accent} Phần I. Mức độ Dễ (Nhận biết -- Câu 1 đến Câu 20)}
\textbf{\color{primary} Câu 1.} Gieo một con xúc xắc 6 mặt cân đối. Số kết quả có thể xảy ra đối với số chấm xuất hiện trên mặt con xúc xắc là:\par
\hspace*{1.5em}\textbf{A.} 1 kết quả\quad\quad \textbf{B.} 12 kết quả\quad\quad \textbf{C.} 5 kết quả\quad\quad \textbf{D.} 6 kết quả\par\vspace{2pt}
\textbf{\color{primary} Câu 2.} Hãy trả lời câu hỏi cơ bản sau: Trong một biểu đồ tranh, nếu mỗi biểu tượng ngôi sao biểu diễn 6 học sinh thì 5 biểu tượng ngôi sao biểu diễn bao nhiêu học sinh?\par
\hspace*{1.5em}\textbf{A.} 11 học sinh\quad\quad \textbf{B.} 30 học sinh\quad\quad \textbf{C.} 35 học sinh\quad\quad \textbf{D.} 24 học sinh\par\vspace{2pt}
\textbf{\color{primary} Câu 3.} Trong bài học thuộc Chương IX, xét nội dung: Trong các dữ liệu sau, dữ liệu nào KHÔNG phải là số liệu (dữ liệu định lượng). Đáp án đúng là:\par
\begin{itemize}\setlength{\itemsep}{1pt}\setlength{\parskip}{0pt}
  \item[] \textbf{A.} Tên các môn thể thao yêu thích của học sinh lớp 6A
  \item[] \textbf{B.} Điểm kiểm tra môn Toán của học sinh tổ 2
  \item[] \textbf{C.} Chiều cao (tính bằng cm) của các bạn trong tổ 1
  \item[] \textbf{D.} Số quyển sách đọc được trong tháng của mỗi học sinh
\end{itemize}\vspace{2pt}
\textbf{\color{primary} Câu 4.} Dựa vào kiến thức chuẩn SGK của Chương IX (Dữ liệu và xác suất thực nghiệm), gieo một con xúc xắc 6 mặt cân đối. Số kết quả có thể xảy ra đối với số chấm xuất hiện trên mặt con xúc xắc là:\par
\hspace*{1.5em}\textbf{A.} 6 kết quả\quad\quad \textbf{B.} 5 kết quả\quad\quad \textbf{C.} 1 kết quả\quad\quad \textbf{D.} 12 kết quả\par\vspace{2pt}
\textbf{\color{primary} Câu 5.} Chọn phương án trả lời đúng nhất: Trong một biểu đồ tranh, nếu mỗi biểu tượng ngôi sao biểu diễn 9 học sinh thì 8 biểu tượng ngôi sao biểu diễn bao nhiêu học sinh?\par
\hspace*{1.5em}\textbf{A.} 77 học sinh\quad\quad \textbf{B.} 63 học sinh\quad\quad \textbf{C.} 72 học sinh\quad\quad \textbf{D.} 17 học sinh\par\vspace{2pt}
\textbf{\color{primary} Câu 6.} Khi làm bài tập cơ bản về dữ liệu và xác suất thực nghiệm, một học sinh gặp câu hỏi: Trong các dữ liệu sau, dữ liệu nào KHÔNG phải là số liệu (dữ liệu định lượng). Kết quả nào dưới đây đúng?\par
\begin{itemize}\setlength{\itemsep}{1pt}\setlength{\parskip}{0pt}
  \item[] \textbf{A.} Chiều cao (tính bằng cm) của các bạn trong tổ 1
  \item[] \textbf{B.} Tên các môn thể thao yêu thích của học sinh lớp 6A
  \item[] \textbf{C.} Điểm kiểm tra môn Toán của học sinh tổ 2
  \item[] \textbf{D.} Số quyển sách đọc được trong tháng của mỗi học sinh
\end{itemize}\vspace{2pt}
\textbf{\color{primary} Câu 7.} Xác định giá trị hoặc khẳng định đúng cho bài toán sau: Gieo một con xúc xắc 6 mặt cân đối. Số kết quả có thể xảy ra đối với số chấm xuất hiện trên mặt con xúc xắc là:\par
\hspace*{1.5em}\textbf{A.} 5 kết quả\quad\quad \textbf{B.} 1 kết quả\quad\quad \textbf{C.} 6 kết quả\quad\quad \textbf{D.} 12 kết quả\par\vspace{2pt}
\textbf{\color{primary} Câu 8.} Áp dụng trực tiếp định nghĩa và quy tắc trong Chương IX: Trong một biểu đồ tranh, nếu mỗi biểu tượng ngôi sao biểu diễn 12 học sinh thì 11 biểu tượng ngôi sao biểu diễn bao nhiêu học sinh:\par
\hspace*{1.5em}\textbf{A.} 120 học sinh\quad\quad \textbf{B.} 23 học sinh\quad\quad \textbf{C.} 137 học sinh\quad\quad \textbf{D.} 132 học sinh\par\vspace{2pt}
\textbf{\color{primary} Câu 9.} Kiểm tra nhanh kiến thức nền tảng: Trong các dữ liệu sau, dữ liệu nào KHÔNG phải là số liệu (dữ liệu định lượng)?\par
\begin{itemize}\setlength{\itemsep}{1pt}\setlength{\parskip}{0pt}
  \item[] \textbf{A.} Số quyển sách đọc được trong tháng của mỗi học sinh
  \item[] \textbf{B.} Điểm kiểm tra môn Toán của học sinh tổ 2
  \item[] \textbf{C.} Tên các môn thể thao yêu thích của học sinh lớp 6A
  \item[] \textbf{D.} Chiều cao (tính bằng cm) của các bạn trong tổ 1
\end{itemize}\vspace{2pt}
\textbf{\color{primary} Câu 10.} Cho dữ kiện nhận biết sau đây: Gieo một con xúc xắc 6 mặt cân đối. Số kết quả có thể xảy ra đối với số chấm xuất hiện trên mặt con xúc xắc là. Phương án chính xác là:\par
\hspace*{1.5em}\textbf{A.} 5 kết quả\quad\quad \textbf{B.} 1 kết quả\quad\quad \textbf{C.} 6 kết quả\quad\quad \textbf{D.} 12 kết quả\par\vspace{2pt}
\textbf{\color{primary} Câu 11.} Hãy chỉ ra kết quả đúng trong bốn lựa chọn A, B, C, D đối với yêu cầu: Trong một biểu đồ tranh, nếu mỗi biểu tượng ngôi sao biểu diễn 15 học sinh thì 14 biểu tượng ngôi sao biểu diễn bao nhiêu học sinh:\par
\hspace*{1.5em}\textbf{A.} 29 học sinh\quad\quad \textbf{B.} 210 học sinh\quad\quad \textbf{C.} 215 học sinh\quad\quad \textbf{D.} 195 học sinh\par\vspace{2pt}
\textbf{\color{primary} Câu 12.} Câu hỏi ôn tập lý thuyết và kỹ năng cơ bản (Chương IX): Trong các dữ liệu sau, dữ liệu nào KHÔNG phải là số liệu (dữ liệu định lượng):\par
\begin{itemize}\setlength{\itemsep}{1pt}\setlength{\parskip}{0pt}
  \item[] \textbf{A.} Số quyển sách đọc được trong tháng của mỗi học sinh
  \item[] \textbf{B.} Chiều cao (tính bằng cm) của các bạn trong tổ 1
  \item[] \textbf{C.} Tên các môn thể thao yêu thích của học sinh lớp 6A
  \item[] \textbf{D.} Điểm kiểm tra môn Toán của học sinh tổ 2
\end{itemize}\vspace{2pt}
\textbf{\color{primary} Câu 13.} Theo nội dung sách giáo khoa Kết nối tri thức với cuộc sống (Chương IX), gieo một con xúc xắc 6 mặt cân đối. Số kết quả có thể xảy ra đối với số chấm xuất hiện trên mặt con xúc xắc là:\par
\hspace*{1.5em}\textbf{A.} 1 kết quả\quad\quad \textbf{B.} 12 kết quả\quad\quad \textbf{C.} 5 kết quả\quad\quad \textbf{D.} 6 kết quả\par\vspace{2pt}
\textbf{\color{primary} Câu 14.} Tìm đáp số đúng của bài toán cơ bản sau: Trong một biểu đồ tranh, nếu mỗi biểu tượng ngôi sao biểu diễn 18 học sinh thì 17 biểu tượng ngôi sao biểu diễn bao nhiêu học sinh:\par
\hspace*{1.5em}\textbf{A.} 311 học sinh\quad\quad \textbf{B.} 35 học sinh\quad\quad \textbf{C.} 306 học sinh\quad\quad \textbf{D.} 288 học sinh\par\vspace{2pt}
\textbf{\color{primary} Câu 15.} Xét đối tượng toán học trong Chương IX: Trong các dữ liệu sau, dữ liệu nào KHÔNG phải là số liệu (dữ liệu định lượng). Giá trị hoặc tính chất tương ứng là:\par
\begin{itemize}\setlength{\itemsep}{1pt}\setlength{\parskip}{0pt}
  \item[] \textbf{A.} Chiều cao (tính bằng cm) của các bạn trong tổ 1
  \item[] \textbf{B.} Số quyển sách đọc được trong tháng của mỗi học sinh
  \item[] \textbf{C.} Điểm kiểm tra môn Toán của học sinh tổ 2
  \item[] \textbf{D.} Tên các môn thể thao yêu thích của học sinh lớp 6A
\end{itemize}\vspace{2pt}
\textbf{\color{primary} Câu 16.} Hoàn thành phép xác định cơ bản sau đây: Gieo một con xúc xắc 6 mặt cân đối. Số kết quả có thể xảy ra đối với số chấm xuất hiện trên mặt con xúc xắc là:\par
\hspace*{1.5em}\textbf{A.} 6 kết quả\quad\quad \textbf{B.} 12 kết quả\quad\quad \textbf{C.} 5 kết quả\quad\quad \textbf{D.} 1 kết quả\par\vspace{2pt}
\textbf{\color{primary} Câu 17.} Trong các khẳng định hoặc giá trị dưới đây, đâu là đáp án đúng cho câu hỏi: Trong một biểu đồ tranh, nếu mỗi biểu tượng ngôi sao biểu diễn 21 học sinh thì 20 biểu tượng ngôi sao biểu diễn bao nhiêu học sinh?\par
\hspace*{1.5em}\textbf{A.} 425 học sinh\quad\quad \textbf{B.} 399 học sinh\quad\quad \textbf{C.} 41 học sinh\quad\quad \textbf{D.} 420 học sinh\par\vspace{2pt}
\textbf{\color{primary} Câu 18.} Bài tập khởi động Chương IX: Trong các dữ liệu sau, dữ liệu nào KHÔNG phải là số liệu (dữ liệu định lượng):\par
\begin{itemize}\setlength{\itemsep}{1pt}\setlength{\parskip}{0pt}
  \item[] \textbf{A.} Số quyển sách đọc được trong tháng của mỗi học sinh
  \item[] \textbf{B.} Chiều cao (tính bằng cm) của các bạn trong tổ 1
  \item[] \textbf{C.} Tên các môn thể thao yêu thích của học sinh lớp 6A
  \item[] \textbf{D.} Điểm kiểm tra môn Toán của học sinh tổ 2
\end{itemize}\vspace{2pt}
\textbf{\color{primary} Câu 19.} Quan sát dữ kiện đề bài: Gieo một con xúc xắc 6 mặt cân đối. Số kết quả có thể xảy ra đối với số chấm xuất hiện trên mặt con xúc xắc là. Hãy chọn phương án đúng:\par
\hspace*{1.5em}\textbf{A.} 1 kết quả\quad\quad \textbf{B.} 5 kết quả\quad\quad \textbf{C.} 6 kết quả\quad\quad \textbf{D.} 12 kết quả\par\vspace{2pt}
\textbf{\color{primary} Câu 20.} Sử dụng công thức hoặc tính chất đầu bài của Chương IX để giải quyết yêu cầu: Trong một biểu đồ tranh, nếu mỗi biểu tượng ngôi sao biểu diễn 24 học sinh thì 23 biểu tượng ngôi sao biểu diễn bao nhiêu học sinh:\par
\hspace*{1.5em}\textbf{A.} 528 học sinh\quad\quad \textbf{B.} 557 học sinh\quad\quad \textbf{C.} 47 học sinh\quad\quad \textbf{D.} 552 học sinh\par\vspace{2pt}
\subsection*{\color{accent} Phần II. Mức độ Trung bình (Thông hiểu -- Câu 21 đến Câu 40)}
\textbf{\color{primary} Câu 21.} Tung một đồng xu cân đối 20 lần liên tiếp, thấy có 11 lần xuất hiện mặt Sấp (S). Xác suất thực nghiệm xuất hiện mặt Sấp bằng:\par
\hspace*{1.5em}\textbf{A.} 9/20\quad\quad \textbf{B.} 20/11\quad\quad \textbf{C.} 11/100\quad\quad \textbf{D.} 11/20\par\vspace{2pt}
\textbf{\color{primary} Câu 22.} Vận dụng quy tắc trong Chương IX (Dữ liệu và xác suất thực nghiệm) để giải bài toán: Gieo một con xúc xắc 40 lần, trong đó mặt 6 chấm xuất hiện 8 lần. Xác suất thực nghiệm xuất hiện mặt 6 chấm (viết dưới dạng phân số tối giản) là:\par
\hspace*{1.5em}\textbf{A.} 1/4\quad\quad \textbf{B.} 2/5\quad\quad \textbf{C.} 1/6\quad\quad \textbf{D.} 1/5\par\vspace{2pt}
\textbf{\color{primary} Câu 23.} Biến đổi và tính toán theo kiến thức Chương IX: Thống kê số lượng sách mượn của 4 tổ trong tuần lần lượt là 14, 17, 20 và 25 quyển. Tổng số sách cả 4 tổ đã mượn là. Kết quả thu được là:\par
\hspace*{1.5em}\textbf{A.} 80 quyển\quad\quad \textbf{B.} 76 quyển\quad\quad \textbf{C.} 72 quyển\quad\quad \textbf{D.} 68 quyển\par\vspace{2pt}
\textbf{\color{primary} Câu 24.} Xét bài toán mức độ thông hiểu sau: Tung một đồng xu cân đối 50 lần liên tiếp, thấy có 26 lần xuất hiện mặt Sấp (S). Xác suất thực nghiệm xuất hiện mặt Sấp bằng. Đáp án đúng là:\par
\hspace*{1.5em}\textbf{A.} 26/50\quad\quad \textbf{B.} 24/50\quad\quad \textbf{C.} 50/26\quad\quad \textbf{D.} 26/100\par\vspace{2pt}
\textbf{\color{primary} Câu 25.} Để hoàn thành phiếu học tập Chương IX, hãy xác định: Gieo một con xúc xắc 70 lần, trong đó mặt 6 chấm xuất hiện 14 lần. Xác suất thực nghiệm xuất hiện mặt 6 chấm (viết dưới dạng phân số tối giản) là:\par
\hspace*{1.5em}\textbf{A.} 2/5\quad\quad \textbf{B.} 1/6\quad\quad \textbf{C.} 1/4\quad\quad \textbf{D.} 1/5\par\vspace{2pt}
\textbf{\color{primary} Câu 26.} Phân tích dữ kiện bài toán: Thống kê số lượng sách mượn của 4 tổ trong tuần lần lượt là 17, 20, 23 và 28 quyển. Tổng số sách cả 4 tổ đã mượn là. Giá trị cần tìm bằng:\par
\hspace*{1.5em}\textbf{A.} 88 quyển\quad\quad \textbf{B.} 80 quyển\quad\quad \textbf{C.} 84 quyển\quad\quad \textbf{D.} 92 quyển\par\vspace{2pt}
\textbf{\color{primary} Câu 27.} Áp dụng tính chất đã học ở Chương IX vào bài tập sau: Tung một đồng xu cân đối 80 lần liên tiếp, thấy có 41 lần xuất hiện mặt Sấp (S). Xác suất thực nghiệm xuất hiện mặt Sấp bằng:\par
\hspace*{1.5em}\textbf{A.} 39/80\quad\quad \textbf{B.} 41/100\quad\quad \textbf{C.} 41/80\quad\quad \textbf{D.} 80/41\par\vspace{2pt}
\textbf{\color{primary} Câu 28.} Thực hiện các bước suy luận hợp lý để giải quyết yêu cầu: Gieo một con xúc xắc 100 lần, trong đó mặt 6 chấm xuất hiện 20 lần. Xác suất thực nghiệm xuất hiện mặt 6 chấm (viết dưới dạng phân số tối giản) là:\par
\hspace*{1.5em}\textbf{A.} 2/5\quad\quad \textbf{B.} 1/6\quad\quad \textbf{C.} 1/5\quad\quad \textbf{D.} 1/4\par\vspace{2pt}
\textbf{\color{primary} Câu 29.} Cho bài toán tự luyện (Dữ liệu và xác suất thực nghiệm): Thống kê số lượng sách mượn của 4 tổ trong tuần lần lượt là 20, 23, 26 và 31 quyển. Tổng số sách cả 4 tổ đã mượn là. Phương án nào sau đây chính xác?\par
\hspace*{1.5em}\textbf{A.} 96 quyển\quad\quad \textbf{B.} 104 quyển\quad\quad \textbf{C.} 92 quyển\quad\quad \textbf{D.} 100 quyển\par\vspace{2pt}
\textbf{\color{primary} Câu 30.} Kết hợp các phép biến đổi trong Chương IX, hãy tìm kết quả của bài toán: Tung một đồng xu cân đối 110 lần liên tiếp, thấy có 56 lần xuất hiện mặt Sấp (S). Xác suất thực nghiệm xuất hiện mặt Sấp bằng:\par
\hspace*{1.5em}\textbf{A.} 110/56\quad\quad \textbf{B.} 56/110\quad\quad \textbf{C.} 54/110\quad\quad \textbf{D.} 56/100\par\vspace{2pt}
\textbf{\color{primary} Câu 31.} Một bài kiểm tra định kỳ nội dung Chương IX nêu câu hỏi: Gieo một con xúc xắc 130 lần, trong đó mặt 6 chấm xuất hiện 26 lần. Xác suất thực nghiệm xuất hiện mặt 6 chấm (viết dưới dạng phân số tối giản) là. Hãy chọn đáp số đúng:\par
\hspace*{1.5em}\textbf{A.} 1/4\quad\quad \textbf{B.} 2/5\quad\quad \textbf{C.} 1/5\quad\quad \textbf{D.} 1/6\par\vspace{2pt}
\textbf{\color{primary} Câu 32.} Tìm lời giải đúng cho bài toán thông hiểu: Thống kê số lượng sách mượn của 4 tổ trong tuần lần lượt là 23, 26, 29 và 34 quyển. Tổng số sách cả 4 tổ đã mượn là:\par
\hspace*{1.5em}\textbf{A.} 104 quyển\quad\quad \textbf{B.} 116 quyển\quad\quad \textbf{C.} 112 quyển\quad\quad \textbf{D.} 108 quyển\par\vspace{2pt}
\textbf{\color{primary} Câu 33.} Sử dụng mối liên hệ giữa các đại lượng trong Chương IX: Tung một đồng xu cân đối 140 lần liên tiếp, thấy có 71 lần xuất hiện mặt Sấp (S). Xác suất thực nghiệm xuất hiện mặt Sấp bằng:\par
\hspace*{1.5em}\textbf{A.} 71/140\quad\quad \textbf{B.} 140/71\quad\quad \textbf{C.} 69/140\quad\quad \textbf{D.} 71/100\par\vspace{2pt}
\textbf{\color{primary} Câu 34.} Hãy tính toán và chọn phương án đúng đối với bài tập: Gieo một con xúc xắc 160 lần, trong đó mặt 6 chấm xuất hiện 32 lần. Xác suất thực nghiệm xuất hiện mặt 6 chấm (viết dưới dạng phân số tối giản) là:\par
\hspace*{1.5em}\textbf{A.} 1/6\quad\quad \textbf{B.} 2/5\quad\quad \textbf{C.} 1/4\quad\quad \textbf{D.} 1/5\par\vspace{2pt}
\textbf{\color{primary} Câu 35.} Dựa trên các tính chất trọng tâm của Chương IX, xét bài toán: Thống kê số lượng sách mượn của 4 tổ trong tuần lần lượt là 26, 29, 32 và 37 quyển. Tổng số sách cả 4 tổ đã mượn là:\par
\hspace*{1.5em}\textbf{A.} 120 quyển\quad\quad \textbf{B.} 124 quyển\quad\quad \textbf{C.} 128 quyển\quad\quad \textbf{D.} 116 quyển\par\vspace{2pt}
\textbf{\color{primary} Câu 36.} Thực hành giải toán (Dữ liệu và xác suất thực nghiệm): Tung một đồng xu cân đối 170 lần liên tiếp, thấy có 86 lần xuất hiện mặt Sấp (S). Xác suất thực nghiệm xuất hiện mặt Sấp bằng. Đáp án nào dưới đây thỏa mãn?\par
\hspace*{1.5em}\textbf{A.} 86/170\quad\quad \textbf{B.} 86/100\quad\quad \textbf{C.} 170/86\quad\quad \textbf{D.} 84/170\par\vspace{2pt}
\textbf{\color{primary} Câu 37.} Kiểm tra kỹ năng biến đổi và suy luận qua bài toán: Gieo một con xúc xắc 190 lần, trong đó mặt 6 chấm xuất hiện 38 lần. Xác suất thực nghiệm xuất hiện mặt 6 chấm (viết dưới dạng phân số tối giản) là:\par
\hspace*{1.5em}\textbf{A.} 2/5\quad\quad \textbf{B.} 1/5\quad\quad \textbf{C.} 1/4\quad\quad \textbf{D.} 1/6\par\vspace{2pt}
\textbf{\color{primary} Câu 38.} Cho giả thiết thuộc phạm vi Chương IX: Thống kê số lượng sách mượn của 4 tổ trong tuần lần lượt là 29, 32, 35 và 40 quyển. Tổng số sách cả 4 tổ đã mượn là. Kết luận đúng là:\par
\hspace*{1.5em}\textbf{A.} 128 quyển\quad\quad \textbf{B.} 140 quyển\quad\quad \textbf{C.} 132 quyển\quad\quad \textbf{D.} 136 quyển\par\vspace{2pt}
\textbf{\color{primary} Câu 39.} Rút gọn hoặc tính giá trị theo yêu cầu sau của Chương IX: Tung một đồng xu cân đối 200 lần liên tiếp, thấy có 101 lần xuất hiện mặt Sấp (S). Xác suất thực nghiệm xuất hiện mặt Sấp bằng:\par
\hspace*{1.5em}\textbf{A.} 101/100\quad\quad \textbf{B.} 99/200\quad\quad \textbf{C.} 101/200\quad\quad \textbf{D.} 200/101\par\vspace{2pt}
\textbf{\color{primary} Câu 40.} Xác định phương án đúng sau khi thực hiện phép tính / suy luận: Gieo một con xúc xắc 220 lần, trong đó mặt 6 chấm xuất hiện 44 lần. Xác suất thực nghiệm xuất hiện mặt 6 chấm (viết dưới dạng phân số tối giản) là:\par
\hspace*{1.5em}\textbf{A.} 2/5\quad\quad \textbf{B.} 1/4\quad\quad \textbf{C.} 1/6\quad\quad \textbf{D.} 1/5\par\vspace{2pt}
\subsection*{\color{accent} Phần III. Mức độ Khó (Vận dụng \& Vận dụng cao -- Câu 41 đến Câu 50)}
\textbf{\color{primary} Câu 41.} Một hộp có các quả bóng màu xanh, đỏ và vàng. Thực hiện lấy ngẫu nhiên 1 quả bóng rồi trả lại hộp 40 lần, thấy bóng xanh xuất hiện 15 lần, bóng đỏ xuất hiện 12 lần. Xác suất thực nghiệm xuất hiện bóng vàng là:\par
\hspace*{1.5em}\textbf{A.} 15/40\quad\quad \textbf{B.} 12/40\quad\quad \textbf{C.} 13/40\quad\quad \textbf{D.} 27/40\par\vspace{2pt}
\textbf{\color{primary} Câu 42.} Bài toán vận dụng (Chương IX – Dữ liệu và xác suất thực nghiệm): Gieo một con xúc xắc 60 lần, số lần xuất hiện mặt số chấm chẵn (2, 4, 6) lần lượt là 11, 13 và 12 lần. Xác suất thực nghiệm của sự kiện 'Số chấm xuất hiện là số lẻ' bằng. Kết quả đúng là:\par
\hspace*{1.5em}\textbf{A.} 1/3\quad\quad \textbf{B.} 3/5\quad\quad \textbf{C.} 1/2\quad\quad \textbf{D.} 2/5\par\vspace{2pt}
\textbf{\color{primary} Câu 43.} Vận dụng kiến thức nâng cao trong phạm vi Chương IX để giải bài toán: Một hộp có các quả bóng màu xanh, đỏ và vàng. Thực hiện lấy ngẫu nhiên 1 quả bóng rồi trả lại hộp 60 lần, thấy bóng xanh xuất hiện 21 lần, bóng đỏ xuất hiện 20 lần. Xác suất thực nghiệm xuất hiện bóng vàng là:\par
\hspace*{1.5em}\textbf{A.} 20/60\quad\quad \textbf{B.} 19/60\quad\quad \textbf{C.} 21/60\quad\quad \textbf{D.} 41/60\par\vspace{2pt}
\textbf{\color{primary} Câu 44.} Phân tích và giải bài toán tổng hợp sau của Chương IX: Gieo một con xúc xắc 60 lần, số lần xuất hiện mặt số chấm chẵn (2, 4, 6) lần lượt là 13, 15 và 8 lần. Xác suất thực nghiệm của sự kiện 'Số chấm xuất hiện là số lẻ' bằng:\par
\hspace*{1.5em}\textbf{A.} 1/2\quad\quad \textbf{B.} 2/5\quad\quad \textbf{C.} 1/3\quad\quad \textbf{D.} 3/5\par\vspace{2pt}
\textbf{\color{primary} Câu 45.} Thử thách tư duy trong phạm vi kiến thức Chương IX: Một hộp có các quả bóng màu xanh, đỏ và vàng. Thực hiện lấy ngẫu nhiên 1 quả bóng rồi trả lại hộp 80 lần, thấy bóng xanh xuất hiện 27 lần, bóng đỏ xuất hiện 28 lần. Xác suất thực nghiệm xuất hiện bóng vàng là. Đáp số chính xác là:\par
\hspace*{1.5em}\textbf{A.} 27/80\quad\quad \textbf{B.} 28/80\quad\quad \textbf{C.} 55/80\quad\quad \textbf{D.} 25/80\par\vspace{2pt}
\textbf{\color{primary} Câu 46.} Sử dụng phương pháp giải toán của Chương IX để tìm đáp án cho bài toán: Gieo một con xúc xắc 60 lần, số lần xuất hiện mặt số chấm chẵn (2, 4, 6) lần lượt là 15, 17 và 4 lần. Xác suất thực nghiệm của sự kiện 'Số chấm xuất hiện là số lẻ' bằng:\par
\hspace*{1.5em}\textbf{A.} 1/2\quad\quad \textbf{B.} 1/3\quad\quad \textbf{C.} 3/5\quad\quad \textbf{D.} 2/5\par\vspace{2pt}
\textbf{\color{primary} Câu 47.} Xét bài toán vận dụng thực tiễn / nâng cao thuộc Chương IX: Một hộp có các quả bóng màu xanh, đỏ và vàng. Thực hiện lấy ngẫu nhiên 1 quả bóng rồi trả lại hộp 100 lần, thấy bóng xanh xuất hiện 33 lần, bóng đỏ xuất hiện 36 lần. Xác suất thực nghiệm xuất hiện bóng vàng là:\par
\hspace*{1.5em}\textbf{A.} 36/100\quad\quad \textbf{B.} 69/100\quad\quad \textbf{C.} 33/100\quad\quad \textbf{D.} 31/100\par\vspace{2pt}
\textbf{\color{primary} Câu 48.} Tìm kết quả cuối cùng của bài toán phân loại học sinh khá giỏi (Chương IX): Gieo một con xúc xắc 60 lần, số lần xuất hiện mặt số chấm chẵn (2, 4, 6) lần lượt là 17, 19 và 0 lần. Xác suất thực nghiệm của sự kiện 'Số chấm xuất hiện là số lẻ' bằng:\par
\hspace*{1.5em}\textbf{A.} 2/5\quad\quad \textbf{B.} 1/3\quad\quad \textbf{C.} 1/2\quad\quad \textbf{D.} 3/5\par\vspace{2pt}
\textbf{\color{primary} Câu 49.} Kết hợp các tính chất đã học trong Chương IX, hãy giải quyết bài toán: Một hộp có các quả bóng màu xanh, đỏ và vàng. Thực hiện lấy ngẫu nhiên 1 quả bóng rồi trả lại hộp 120 lần, thấy bóng xanh xuất hiện 39 lần, bóng đỏ xuất hiện 44 lần. Xác suất thực nghiệm xuất hiện bóng vàng là:\par
\hspace*{1.5em}\textbf{A.} 37/120\quad\quad \textbf{B.} 44/120\quad\quad \textbf{C.} 39/120\quad\quad \textbf{D.} 83/120\par\vspace{2pt}
\textbf{\color{primary} Câu 50.} Câu hỏi vận dụng cao cuối Chương IX: Gieo một con xúc xắc 60 lần, số lần xuất hiện mặt số chấm chẵn (2, 4, 6) lần lượt là 19, 21 và -4 lần. Xác suất thực nghiệm của sự kiện 'Số chấm xuất hiện là số lẻ' bằng. Phương án đúng là:\par
\hspace*{1.5em}\textbf{A.} 1/3\quad\quad \textbf{B.} 3/5\quad\quad \textbf{C.} 1/2\quad\quad \textbf{D.} 2/5\par\vspace{2pt}
\subsection*{\color{primary} Bảng đáp án Chương IX: Dữ liệu và xác suất thực nghiệm}
\begin{center}
\renewcommand{\arraystretch}{1.25}
\begin{tabularx}{\textwidth}{*{10}{>{\centering\arraybackslash}X}}
\toprule
\textbf{\color{primary} Cột 1} & \textbf{\color{primary} Cột 2} & \textbf{\color{primary} Cột 3} & \textbf{\color{primary} Cột 4} & \textbf{\color{primary} Cột 5} & \textbf{\color{primary} Cột 6} & \textbf{\color{primary} Cột 7} & \textbf{\color{primary} Cột 8} & \textbf{\color{primary} Cột 9} & \textbf{\color{primary} Cột 10} \\
\midrule
\textbf{1.D} & \textbf{2.B} & \textbf{3.A} & \textbf{4.A} & \textbf{5.C} & \textbf{6.B} & \textbf{7.C} & \textbf{8.D} & \textbf{9.C} & \textbf{10.C} \\
\rowcolor{rowtint} \textbf{11.B} & \textbf{12.C} & \textbf{13.D} & \textbf{14.C} & \textbf{15.D} & \textbf{16.A} & \textbf{17.D} & \textbf{18.C} & \textbf{19.C} & \textbf{20.D} \\
\textbf{21.D} & \textbf{22.D} & \textbf{23.B} & \textbf{24.A} & \textbf{25.D} & \textbf{26.A} & \textbf{27.C} & \textbf{28.C} & \textbf{29.D} & \textbf{30.B} \\
\rowcolor{rowtint} \textbf{31.C} & \textbf{32.C} & \textbf{33.A} & \textbf{34.D} & \textbf{35.B} & \textbf{36.A} & \textbf{37.B} & \textbf{38.D} & \textbf{39.C} & \textbf{40.D} \\
\textbf{41.C} & \textbf{42.D} & \textbf{43.B} & \textbf{44.B} & \textbf{45.D} & \textbf{46.D} & \textbf{47.D} & \textbf{48.A} & \textbf{49.A} & \textbf{50.D} \\
\bottomrule
\end{tabularx}
\end{center}\vspace{10pt}
\end{document}
